\documentclass[12pt]{book}

%% ==================== 用户可配置选项 ====================

%% 1. 字体配置
\def\CJKMainFont{Noto Serif CJK SC}      % 中文字体
\def\CJKMainFontBold{Noto Sans CJK SC Bold}
\def\CJKSansFont{Noto Sans CJK SC}
\def\CJKMonoFont{Noto Sans Mono CJK SC}
\def\LatinMainFont{DejaVu Serif}        % 英文字体
\def\LatinSansFont{DejaVu Sans}
\def\LatinMonoFont{DejaVu Sans Mono}

%% 2. 页边距配置（单位：cm）
\def\PageMarginTop{2.5}
\def\PageMarginBottom{2.5}
\def\PageMarginLeft{2}
\def\PageMarginRight{2}

%% 3. 打印模式（双面打印留装订边）
\def\EnablePrintMode{false}  % true/false
\def\BindingOffset{0.5cm}

%% 4. 行间距
\def\LineStretch{1.2}

%% 5. 段间距
\def\ParagraphSkip{2pt}

%% ==================== 配置结束 ====================

%% 页边距设置
\usepackage{geometry}
\usepackage{ifthen}

\ifthenelse{\equal{\EnablePrintMode}{true}}{%
  %% 打印模式：双面，留装订边
  \geometry{top=\PageMarginTop cm, bottom=\PageMarginBottom cm,
             left=\PageMarginLeft cm, right=\PageMarginRight cm,
             bindingoffset=\BindingOffset,
             twoside}
}{%
  %% 普通模式：单面
  \geometry{top=\PageMarginTop cm, bottom=\PageMarginBottom cm,
             left=\PageMarginLeft cm, right=\PageMarginRight cm}
}

\usepackage{xeCJK}
\xeCJKsetup{AutoFallBack=true}
\setCJKmainfont{\CJKMainFont}[BoldFont=\CJKMainFontBold,ItalicFont=\CJKMainFont]
\setCJKsansfont{\CJKSansFont}[BoldFont=\CJKMainFontBold]
\setCJKmonofont{\CJKMonoFont}

%% 行间距设置
\usepackage{setspace}
\setstretch{\LineStretch}

%% 段间距设置
\setlength{\parskip}{\ParagraphSkip}

% 目录格式调整 - 修复编号宽度问题
\usepackage{tocloft}
\setlength{\cftsecnumwidth}{3em}  % section 编号宽度
\setlength{\cftsubsecnumwidth}{4em}  % subsection 编号宽度
\setlength{\cftchapnumwidth}{2em}  % chapter 编号宽度
\usepackage{listings}
\usepackage{graphicx}
\usepackage{xcolor}
\usepackage{url}
% \usepackage{showidx}
\usepackage{imakeidx}
\usepackage{booktabs}
\usepackage{url}
\usepackage{etoolbox}  % for showidx
\usepackage{soul}
% \usepackage{gnuplot-lua-tikz}  % 中文编译：gnuplot-lua-tikz.sty未安装
\usepackage{tikz}
\usetikzlibrary{arrows,positioning,calc}

\usepackage{fontspec}
\setmainfont{\LatinMainFont}
\setsansfont{\LatinSansFont}
\setmonofont{\LatinMonoFont}

\usepackage{hyperref}  % should be last
\hypersetup{unicode=true,
            pdfauthor={Russ Cox, Frans Kaashoek, Robert Morris},
            pdftitle={xv6: a simple, Unix-like teaching operating system},}

\lstset{basicstyle=\small\ttfamily}
\lstset{morecomment=[is]{[[[}{]]]}}
\lstset{escapeinside={(*@}{@*)}}
\lstset{xleftmargin=5.0ex}

\newcommand{\github}{https://github.com/mit-pdos/xv6-riscv/blob/riscv/}

\newcommand{\fileref}[1]{\small{\textcolor{red}{(#1)}}}
\newcommand{\lineref}[2]{\small{\textcolor{red}{(#1:#2)}}}
\newcommand{\linerefs}[3]{\small{\textcolor{red}{(#1:#2-#3)}}}

% MFK: this (or file:) gives a security alert
% \newcommand{\showlineref}[3]{\href{run:./xv6-riscv-src/#1:#2}{\small{(\textcolor{blue}{#3})}}}

\newcommand{\showfileref}[2]{\href{\github/#1}{\small{(\textcolor{blue}{#2})}}}
\newcommand{\showlineref}[3]{\href{\github/#1\#L#2}{\small{(\textcolor{blue}{#3})}}}
\newcommand{\showlinerefs}[5]{\href{\github/#1\#L#2-L#3}{\small\textcolor{blue}{(#4-#5)}}}

\newcommand{\indextext}[1]{\textit{#1}\index{#1}}
\newcommand{\indextextx}[1]{{#1}\index{#1}}
\newcommand{\indexcode}[1]{\lstinline{#1}\index{#1@\lstinline{#1}}}

%% editing markup
\newcommand{\insertnote}[3]{\noindent\textcolor{#1}{\textbf{#2:} #3}}
\newcommand{\note}[1]{\insertnote{blue}{NOTE}{#1}}
\newcommand{\rtm}[1]{\insertnote{red}{RTM}{#1}}
\newcommand{\mfk}[1]{\insertnote{red}{MFK}{#1}}
%% for publishing book without notes
%\renewcommand{\insertnote}[3]{}

\title{\textbf{xv6: a simple, Unix-like teaching operating system}}
\author{Russ Cox \and Frans Kaashoek \and Robert Morris}

\makeindex

\begin{document}

\maketitle

\tableofcontents

\chapter*{前言和致谢}


这是一本面向操作系统课程用的草稿教材。它通过研究一个名为 xv6 的内核来阐述操作系统的主要概念。Xv6 以 Dennis Ritchie 和 Ken Thompson 的 Unix Version 6 (v6) 为模型。Xv6 大致遵循 v6 的结构和风格，但使用 ANSI C 为 RISC-V 多核处理器实现。

本文应与 xv6 源代码一起阅读，这种方法受 John Lions 的《UNIX 第6版注释（Lions' Commentary on UNIX 6th Edition）》启发；文本有指向源代码的超链接，位于 \url{https://github.com/mit-pdos/xv6-riscv}。参见 \url{https://pdos.csail.mit.edu/6.1810} 以获取有关 v6 和 xv6 的额外在线资源，包括使用 xv6 的几个实验作业。

我们在麻省理工学院的操作系统课程 6.828 和 6.1810 中使用了本文。我们感谢这些课程的教师、助教和学生，他们都以直接或间接的方式为 xv6 做出了贡献。特别是，我们要感谢 Adam Belay、Austin Clements 和 Nickolai Zeldovich。最后，我们要感谢那些通过邮件向我们报告文本中的错误或提出改进建议的人：Abutalib Aghayev、Sebastian Boehm、brandb97、Anton Burtsev、Raphael Carvalho、Tej Chajed、Brendan Davidson、Rasit Eskicioglu、Color Fuzzy、Wojciech Gac、Giuseppe、Tao Guo、Haibo Hao、Naoki Hayama、Chris Henderson、Robert Hilderman、Eden Hochbaum、Wolfgang Keller、Paweł Kraszewski、Henry Laih、Jin Li、Austin Liew、lyazj@github.com、Pavan Maddamsetti、Jacek Masiulaniec、Michael McConville、m3hm00d、Mes0903、miguelgvieira、Mark Morrissey、Muhammed Mourad、Harry Pan、Harry Porter、Siyuan Qian、Zhefeng Qiao、Askar Safin、Salman Shah、Huang Sha、Vikram Shenoy、Adeodato Simó、Ruslan Savchenko、Pawel Szczurko、Warren Toomey、tyfkda、tzerbib、Vanush Vaswani、Chen Wang、Xi Wang、Zou Chang Wei、Sam Whitlock、Qiongsi Wu、LucyShawYang、ykf1114@gmail.com 和 Meng Zhou。

如果您发现错误或有改进建议，请发送邮件给 Frans Kaashoek 和 Robert Morris（kaashoek,rtm@csail.mit.edu）。

\chapter{操作系统接口}
\label{CH:UNIX}

操作系统的工作是在多个程序之间共享计算机，并提供比硬件本身支持的更有用的一组服务。操作系统管理和抽象底层硬件，这样，例如，文字处理器无需关心正在使用哪种类型的磁盘硬件。操作系统在多个程序之间共享硬件，使它们能够同时运行（或看起来同时运行）。最后，操作系统为程序提供受控的交互方式，使它们能够共享数据或协同工作。

操作系统通过接口向用户程序提供服务。
\index{interface design}
设计一个好的接口是困难的。一方面，我们希望接口简单且窄，因为这样更容易正确实现。另一方面，我们可能想向应用程序提供许多复杂的功能。解决这种张力的技巧是设计依赖于少数机制的接口，这些机制可以组合起来提供很大的通用性。

本书使用单一操作系统作为具体示例来说明操作系统概念。该操作系统 xv6 提供了 Ken Thompson 和 Dennis Ritchie 的 Unix 操作系统~\cite{unix} 引入的基本接口，并模仿 Unix 的内部设计。Unix 提供了一个窄接口，其机制组合良好，提供了令人惊讶的通用性。这个接口非常成功，以至于现代操作系统——BSD、Linux、macOS、Solaris，甚至较小程度上包括 Microsoft Windows——都具有类似 Unix 的接口。理解 xv6 是理解这些系统以及许多其他系统的良好开端。

如
Figure~\ref{fig:os}
所示，xv6 采用传统的
\indextext{kernel}（内核）
形式，这是一个向运行程序提供服务的特殊程序。每个运行程序称为一个
\indextext{process}（进程），
包含指令、数据和栈的内存。
指令实现程序的计算。数据是计算操作的变量。栈组织程序的过程调用。一台给定的计算机通常有许多进程，但只有一个内核。

当进程需要调用内核服务时，它调用
\indextext{system call}（系统调用），
即操作系统接口中的一个调用。系统调用进入内核；内核执行服务并返回。因此，进程在
\indextext{user space}（用户空间）
和
\indextext{kernel space}（内核空间）
之间交替执行。

如后续章节详细描述，内核使用 CPU 提供的硬件保护机制\footnote{
本文通常用术语 \indextext{CPU}（中央处理单元的缩写）来指执行计算的硬件元素。其他文档（例如 RISC-V 规范）也使用处理器、核心和 hart 等词代替 CPU。
}
来确保在用户空间中执行的每个进程只能访问自己的内存。内核以实施这些保护所需的硬件特权执行；用户程序在没有这些特权的情况下执行。当用户程序调用系统调用时，硬件提升特权级别并开始执行内核中预先安排的函数。

\begin{figure}[t]
\centering
\includegraphics[scale=0.5]{fig/os.pdf}
\caption{内核和两个用户进程。}
\label{fig:os}
\end{figure}

内核提供的系统调用集合是用户程序看到的接口。xv6 内核提供了 Unix 内核传统提供的服务和系统调用的一个子集。
Figure~\ref{fig:api}
列出了 xv6 的所有系统调用。

本章其余部分概述了 xv6 的服务——进程、内存、文件描述符（file descriptor）、管道（pipe）和文件系统——并通过代码片段和讨论
\indextext{shell}（shell，
Unix 的命令行用户界面）如何使用它们来说明。shell 对系统调用的使用说明了它们被设计得多么仔细。

shell 是一个普通程序，从用户读取命令并执行它们。shell 是一个用户程序，而不是内核的一部分，这一事实说明了系统调用接口的强大：shell 没有什么特别之处。这也意味着 shell 很容易被替换；因此，现代 Unix 系统有多种 shell 可供选择，每种都有自己的用户界面和脚本功能。xv6 shell 是 Unix Bourne shell 本质的简单实现。

xv6 shell 的实现可以在
\fileref{user/sh.c}
找到。（\textcolor{blue}{链接}
是指向
\url{https://github.com/mit-pdos/xv6-riscv/}
上相关 xv6 源代码的超链接，具体数字指的是
\href{xv6-src-booklet.pdf}{xv6-src-booklet.pdf}
中的页码和行号，如 Lions' Commentary on UNIX 6th Edition~\cite{lions} 中所示。一个好的做法是先在你自己的开发环境中（在 github 上或 PDF 查看器中）阅读源代码，然后再回到本书。到本书结束时，你应该能够理解 xv6 源代码的每一行，而无需查阅本书。


%%
%%	Processes and memory
%%
\section{进程和内存}

xv6 进程由用户空间内存（指令、数据和栈）和内核私有的每进程状态组成。Xv6
\indextext{time-share}s
分时（time-sharing）进程：它在等待执行的进程集合之间透明地切换可用的 CPU。当进程不执行时，xv6 保存进程的 CPU 寄存器，在下次运行该进程时恢复它们。内核为每个进程关联一个进程标识符，或
\indexcode{PID}（PID）。

\begin{figure}[t]
\centering
\begin{tabular}{ll}
{\bf 系统调用} & {\bf 描述} \\
\midrule
int fork() & 创建进程，返回子进程的 PID。 \\
int exit(int status) & 终止当前进程；状态报告给 wait()。不返回。 \\
int wait(int *status) & 等待子进程退出；退出状态在 *status 中；返回子进程 PID。 \\
int kill(int pid) & 终止进程 PID。返回 0，出错返回 -1。 \\
int getpid() & 返回当前进程的 PID。 \\
int pause(int n) & 暂停 n 个时钟滴答。 \\
int exec(char *file, char *argv[]) & 加载文件并用参数执行；仅出错时返回。 \\
char *sbrk(int n) & 将进程内存增加 n 个零字节。返回新内存的起始地址。 \\
int open(char *file, int flags) & 打开文件；标志指示读/写；返回 fd（文件描述符）。 \\
int write(int fd, char *buf, int n) & 将 n 字节从 buf 写入文件描述符 fd；返回 n。 \\
int read(int fd, char *buf, int n) & 将 n 字节读入 buf；返回读取的数量；或 0 如果文件结束。 \\
int close(int fd) & 释放打开的文件 fd。 \\
int dup(int fd) & 返回引用与 fd 相同文件的新文件描述符。\\
int pipe(int p[]) & 创建管道，将读/写文件描述符放入 p[0] 和 p[1]。 \\
int chdir(char *dir) & 更改当前目录。 \\
int mkdir(char *dir) & 创建新目录。 \\
int mknod(char *file, int, int) & 创建设备文件。 \\
int fstat(int fd, struct stat *st) & 将打开文件的信息放入 *st。 \\
int link(char *file1, char *file2) & 为文件 file1 创建另一个名称（file2）。 \\
int unlink(char *file) & 删除文件。 \\
\end{tabular}
\caption{Xv6 系统调用。除非另有说明，这些调用在无错误时返回 0，出错时返回 -1。}
\label{fig:api}
\end{figure}

进程可以使用
\indexcode{fork}
系统调用创建新进程。
\lstinline{fork}
给新进程一个调用进程内存的精确副本：\lstinline{fork} 将调用进程的指令、数据和栈复制到新进程的内存中。
\lstinline{fork}
在原始进程和新进程中都返回。在原始进程中，\lstinline{fork} 返回新进程的 PID。在新进程中，\lstinline{fork} 返回零。原始进程和新进程通常称为
\indextext{parent}（父进程）
和
\indextext{child}（子进程）。

例如，考虑以下用 C 编程语言~\cite{kernighan}编写的程序片段：
\begin{lstlisting}[]
int pid = fork();
if(pid > 0){
  printf("parent: child=%d\n", pid);
  pid = wait((int *) 0);
  printf("child %d is done\n", pid);
} else if(pid == 0){
  printf("child: exiting\n");
  exit(0);
} else {
  printf("fork error\n");
}
\end{lstlisting}
\indexcode{exit}
系统调用导致调用进程停止执行并释放资源，如内存和打开的文件。Exit 接受一个整数状态参数，通常 0 表示成功，1 表示失败。
\indexcode{wait}
系统调用返回当前进程的已退出（或被杀死）子进程的 PID，并将子进程的退出状态复制到传递给 wait 的地址；如果调用者的子进程都没有退出，
\indexcode{wait}
等待其中一个退出。如果调用者没有子进程，\lstinline{wait} 立即返回 -1。如果父进程不关心子进程的退出状态，它可以向
\lstinline{wait}
传递 0 地址。

在示例中，输出行
\begin{lstlisting}[]
parent: child=1234
child: exiting
\end{lstlisting}
可能以任意顺序出现（甚至交错），取决于父进程或子进程谁先到达其
\indexcode{printf}
调用。子进程退出后，父进程的
\indexcode{wait}
返回，导致父进程打印
\begin{lstlisting}[]
parent: child 1234 is done
\end{lstlisting}
虽然子进程以父进程内存的副本开始，但父进程和子进程使用独立的内存和独立的寄存器执行：在一个进程中更改变量不会影响另一个进程。例如，当
\lstinline{wait}
的返回值存储到父进程中的
\lstinline{pid}
时，它不会更改子进程中的变量
\lstinline{pid}。子进程中
\lstinline{pid}
的值仍将是零。

\indexcode{exec}
系统调用用从文件系统存储的文件加载的新内存映像替换调用进程的内存。文件必须具有特定格式，指定文件的哪部分保存指令、哪部分是数据、从哪条指令开始执行等。Xv6 使用 ELF 格式，Chapter~\ref{CH:MEM} 将更详细地讨论。通常，文件是编译程序源代码的结果。当
\indexcode{exec}
成功时，它不会返回到调用程序；相反，从文件加载的指令开始在 ELF 头中声明的入口点执行。
\lstinline{exec}
接受两个参数：包含可执行文件的文件名和参数字符串数组。例如：
\begin{lstlisting}[]
char *argv[3];

argv[0] = "echo";
argv[1] = "hello";
argv[2] = 0;
exec("/bin/echo", argv);
printf("exec error\n");
\end{lstlisting}
此片段用
\lstinline{/bin/echo}
程序的一个实例替换调用程序，该实例以参数列表
\lstinline{echo}
\lstinline{hello}
运行。大多数程序忽略参数数组的第一个元素，它通常是程序的名称。

xv6 shell 使用上述调用来代表用户运行程序。shell 的主要结构很简单；参见
\lstinline{main}
\lineref{user/sh.c:/main/}。
主循环使用
\indexcode{getcmd}
从用户读取一行输入。然后它调用
\lstinline{fork}，
它创建 shell 进程的副本。父进程调用
\lstinline{wait}，
而子进程运行命令。例如，如果用户向 shell 输入了
``\lstinline{echo hello}''，
\lstinline{runcmd}
将被调用，参数为
``\lstinline{echo hello}''。
\lstinline{runcmd}
\lineref{user/sh.c:/runcmd/}
运行实际命令。对于
``\lstinline{echo hello}''，
它将调用
\lstinline{exec}
\lineref{user/sh.c:/exec.ecmd/}。
如果
\lstinline{exec}
成功，则子进程将执行来自
\lstinline{echo}
而不是
\lstinline{runcmd}
的指令。
在某个时刻
\lstinline{echo}
将调用
\lstinline{exit}，
这将导致父进程从
\lstinline{main}
\lineref{user/sh.c:/main/}
中的
\lstinline{wait}
返回。

你可能想知道为什么
\indexcode{fork}
和
\indexcode{exec}
没有合并为一个调用；稍后我们将看到 shell 在其实现 I/O 重定向时利用了这种分离。为了避免创建重复进程然后立即替换它（使用 \lstinline{exec}）的浪费，操作系统内核通过使用虚拟内存技术（如写时复制，见 Section~\ref{CH:PGFAULTS}）来优化此用例的
\lstinline{fork}
实现。

Xv6 隐式分配大多数用户空间内存：
\indexcode{fork}
分配子进程复制父进程内存所需的内存，而
\indexcode{exec}
分配足够的内存来容纳可执行文件。在运行时需要更多内存的进程（可能用于
\indexcode{malloc}）
可以调用
\lstinline{sbrk(n)}
将其数据内存增加
\lstinline{n}
个零字节；
\indexcode{sbrk}
返回新内存的位置。

%%
%%	I/O and File descriptors
%%
\section{I/O 和文件描述符}

\indextext{file descriptor}（文件描述符）
是一个小整数，表示进程可以从中读取或写入的内核管理对象。进程可以通过打开文件、目录或设备，或创建管道，或复制现有描述符来获得文件描述符。为简单起见，我们通常将文件描述符引用的对象称为``文件''；文件描述符接口抽象了文件、管道和设备之间的差异，使它们看起来都像字节流。我们将输入和输出称为 \indextext{I/O}（I/O）。

在内部，xv6 内核使用文件描述符作为每进程表的索引，这样每个进程都有从零开始的私有文件描述符空间。按照惯例，进程从文件描述符 0（标准输入）读取，向文件描述符 1（标准输出）写入输出，并向文件描述符 2（标准错误）写入错误消息。正如我们将看到的，shell 利用这一惯例来实现 I/O 重定向和管道。shell 确保它始终有三个文件描述符打开
\lineref{user/sh.c:/open..console/}，
默认情况下它们是控制台（console）的文件描述符。

\lstinline{read}
和
\lstinline{write}
系统调用从文件描述符命名的打开文件读取字节和写入字节。调用
\lstinline{read(fd},
\lstinline{buf},
\lstinline{n)}
从文件描述符
\lstinline{fd}
读取最多
\lstinline{n}
字节，将它们复制到
\lstinline{buf}，
并返回读取的字节数。引用文件的每个文件描述符都有一个与之关联的偏移量。
\lstinline{read}
从当前文件偏移量读取数据，然后将该偏移量增加读取的字节数：后续的
\lstinline{read}
将返回第一个
\lstinline{read}
返回的字节之后的字节。当没有更多字节可读时，
\lstinline{read}
返回零表示文件结束。

调用
\lstinline{write(fd}、
\lstinline{buf}、
\lstinline{n)}
将
\lstinline{n}
字节从
\lstinline{buf}
写入文件描述符
\lstinline{fd}
并返回写入的字节数。只有当发生错误时，写入的字节数才会小于
\lstinline{n}。
与
\lstinline{read}
一样，
\lstinline{write}
在当前文件偏移量处写入数据，然后将该偏移量增加写入的字节数：每个
\lstinline{write}
从上一个写入结束的地方开始。

以下程序片段（形成程序 \lstinline{cat} 的本质）将数据从其标准输入复制到其标准输出。如果发生错误，它会向标准错误写入一条消息。
\begin{lstlisting}[]
char buf[512];
int n;

for(;;){
  n = read(0, buf, sizeof buf);
  if(n == 0)
    break;
  if(n < 0){
    fprintf(2, "read error\n");
    exit(1);
  }
  if(write(1, buf, n) != n){
    fprintf(2, "write error\n");
    exit(1);
  }
}
\end{lstlisting}
代码片段中需要注意的是
\lstinline{cat}
不知道是正在从文件、控制台还是管道读取。同样，
\lstinline{cat}
不知道是正在打印到控制台、文件还是其他什么。使用文件描述符以及文件描述符 0 是输入、文件描述符 1 是输出的惯例允许
\lstinline{cat}
的简单实现。

\lstinline{close}
系统调用释放文件描述符，使其可供将来的
\lstinline{open}、
\lstinline{pipe}
或
\lstinline{dup}
系统调用重用（见下文）。新分配的文件描述符始终是当前进程的最低未使用描述符。

文件描述符和
\indexcode{fork}
相互作用，使 I/O 重定向易于实现。
\lstinline{fork}
复制父进程的文件描述符表及其内存，因此子进程以与父进程完全相同的打开文件开始。
\indexcode{exec}
系统调用替换调用进程的内存但保留其文件表。这种行为允许 shell 通过 fork、在子进程中关闭和重新打开选定的文件描述符，然后调用 \lstinline{exec} 运行新程序来实现 \indextext{I/O redirection}（I/O 重定向）。以下是 shell 为命令
\lstinline{cat}
\lstinline{<}
\lstinline{input.txt}
运行的代码的简化版本：
\begin{lstlisting}[]
char *argv[2];

argv[0] = "cat";
argv[1] = 0;
if(fork() == 0) {
  close(0);
  open("input.txt", O_RDONLY);
  exec("cat", argv);
}
\end{lstlisting}
子进程关闭文件描述符 0 后，
\lstinline{open}
保证将该文件描述符用于新打开的
\lstinline{input.txt}：
0 将是最小的可用文件描述符。
\lstinline{cat}
然后以文件描述符 0（标准输入）引用
\lstinline{input.txt}
执行。父进程的文件描述符不受此序列的影响，因为它只修改子进程的描述符。

xv6 shell 中 I/O 重定向的代码正是以这种方式工作
\lineref{user/sh.c:/case.REDIR/}。
回想一下，此时代码中 shell 已经 fork 了子 shell，并且
\lstinline{runcmd}
将调用
\lstinline{exec}
加载新程序。

\lstinline{open}
的第二个参数由一组标志组成，以位表示，控制 \lstinline{open} 的行为。可能的值在文件控制（fcntl）头文件中定义
\linerefs{kernel/fcntl.h:/RDONLY/,/TRUNC/}：
\lstinline{O_RDONLY}、
\lstinline{O_WRONLY}、
\lstinline{O_RDWR}、
\lstinline{O_CREATE}
和
\lstinline{O_TRUNC}，
它们指示 \lstinline{open}
打开文件进行读取、
写入、
同时读取和写入、
在文件不存在时创建文件，
以及将文件截断为零长度。

现在应该清楚为什么
\lstinline{fork}
和
\lstinline{exec}
是分开的调用是有帮助的：在这两者之间，shell 有机会重定向子进程的 I/O，而不会干扰主 shell 的 I/O 设置。可以想象一个假设的组合
\lstinline{forkexec} 系统调用，
但使用这种调用来进行 I/O 重定向的选项似乎很尴尬。shell 可以在调用 \lstinline{forkexec} 之前修改自己的 I/O 设置（然后撤消这些修改）；或者
\lstinline{forkexec}
可以接受 I/O 重定向的指令作为参数；或者（最不具吸引力的）每个像 \lstinline{cat} 这样的程序都可以被教导自己做 I/O 重定向。

尽管
\lstinline{fork}
复制文件描述符表，但每个底层文件偏移量在父进程和子进程之间共享。考虑这个例子：
\begin{lstlisting}[]
if(fork() == 0) {
  write(1, "hello ", 6);
  exit(0);
} else {
  wait(0);
  write(1, "world\n", 6);
}
\end{lstlisting}
在此片段结束时，附加到文件描述符 1 的文件将包含数据
\lstinline{hello}
\lstinline{world}。
父进程中的
\lstinline{write}
（由于
\lstinline{wait}，
仅在子进程完成后运行）从子进程的
\lstinline{write}
停止的地方继续。这种行为有助于从 shell 命令序列生成顺序输出，如
\lstinline{(echo}
\lstinline{hello};
\lstinline{echo}
\lstinline{world)}
\lstinline{>output.txt}。

\lstinline{dup}
系统调用复制现有文件描述符，返回一个引用相同底层 I/O 对象的新描述符。两个文件描述符共享偏移量，就像
\lstinline{fork}
复制的文件描述符一样。这是将
\lstinline{hello}
\lstinline{world}
写入文件的另一种方式：
\begin{lstlisting}[]
fd = dup(1);
write(1, "hello ", 6);
write(fd, "world\n", 6);
\end{lstlisting}

如果两个文件描述符是通过一系列
\lstinline{fork}
和
\lstinline{dup}
调用从同一原始文件描述符派生的，则它们共享偏移量。否则文件描述符不共享偏移量，即使它们是由同一文件的
\lstinline{open}
调用产生的。
\lstinline{dup}
允许 shell 实现如下命令：
\lstinline{ls}
\lstinline{existing-file}
\lstinline{non-existing-file}
\lstinline{>}
\lstinline{tmp1}
\lstinline{2>&1}。
\lstinline{2>&1}
告诉 shell 给命令一个文件描述符 2，它是描述符 1 的副本。现有文件的名称和不存在的文件的错误消息都将显示在文件
\lstinline{tmp1}
中。Xv6 shell 不支持错误文件描述符的 I/O 重定向，但现在你知道如何实现它了。

文件描述符是一个强大的抽象，因为它们隐藏了它们连接到的内容的细节：向文件描述符 1 写入的进程可能正在写入文件、像控制台这样的设备，或管道。
%%
%%	Pipes
%%
\section{管道}

\indextext{pipe}（管道）
是一个小的内核缓冲区，作为一对文件描述符暴露给进程，一个用于读取，一个用于写入。向管道一端写入数据使该数据可从管道的另一端读取。管道为进程提供了一种通信方式。

下面的示例代码运行程序
\lstinline{wc}
，标准输入连接到管道的读取端。
\begin{lstlisting}[]
int p[2];
char *argv[2];

argv[0] = "wc";
argv[1] = 0;

pipe(p);
if(fork() == 0) {
  close(0);
  dup(p[0]);
  close(p[0]);
  close(p[1]);
  exec("/bin/wc", argv);
} else {
  close(p[0]);
  write(p[1], "hello world\n", 12);
  close(p[1]);
}
\end{lstlisting}
程序调用
\lstinline{pipe}，
它创建一个新管道，并将读取和写入文件描述符记录在数组
\lstinline{p}
中。
\lstinline{fork}
之后，父进程和子进程都有引用管道的文件描述符。子进程调用 \lstinline{close} 和 \lstinline{dup} 使文件描述符零引用管道的读取端，关闭
\lstinline{p}
中的文件描述符，并调用 \lstinline{exec} 运行
\lstinline{wc}。
当
\lstinline{wc}
从其标准输入读取时，它从管道读取。父进程关闭管道的读取端，写入管道，然后关闭写入端。

如果没有数据可用，管道上的
\lstinline{read}
将等待数据写入或所有引用写入端的文件描述符关闭；在后一种情况下，
\lstinline{read}
将返回 0，就像已到达数据文件末尾一样。
\lstinline{read}
阻塞直到不可能有新数据到达这一事实，是子进程在执行
\lstinline{wc}
之前关闭管道写入端的重要原因：如果
\lstinline{wc}
的文件描述符之一引用管道的写入端，
\lstinline{wc}
将永远不会看到文件结束。

xv6 shell 实现管道的方式与上述代码类似，如
\lstinline{grep fork sh.c | wc -l}
\lineref{user/sh.c:/case.PIPE/}。
子进程创建一个管道来连接管道的左端和右端。然后它调用
\lstinline{fork}
和
\lstinline{runcmd}
处理管道的左端，以及
\lstinline{fork}
和
\lstinline{runcmd}
处理管道的右端，并等待两者完成。管道的右端可能是一个本身包含管道的命令（例如，
\lstinline{a}
\lstinline{|}
\lstinline{b}
\lstinline{|}
\lstinline{c)}，
它本身 fork 两个新进程（一个用于
\lstinline{b}，
一个用于
\lstinline{c}）。
因此，shell 可能创建一棵进程树。这棵树的叶子是命令，内部节点是等待左右子进程完成的进程。

% In principle, one could have the interior nodes run the left end of a
% pipeline, but doing so correctly would complicate the
% implementation. Consider making just the following modification:
% change \lstinline{sh.c} to not fork for \lstinline{p->left} and run
% \lstinline{runcmd(p->left)} in the interior process. Then, for
% example, \lstinline{echo hi | wc} won't produce output, because when
% \lstinline{echo hi} exits in \lstinline{runcmd}, the interior process
% exits and never calls fork to run the right end of the pipe.  This
% incorrect behavior could be fixed by not calling \lstinline{exit} in
% \lstinline{runcmd} for interior processes, but this fix complicates
% the code: now \lstinline{runcmd} needs to know if it's in an interior
% process or not.  Complications also arise when not forking for
% \lstinline{runcmd(p->right)}.  For example, with just that
% modification, \lstinline{sleep 10 | echo hi} will immediately print
% ``hi' and a new prompt, instead of after 10 seconds; this happens because \lstinline{echo} runs
% immediately and exits, not waiting for \lstinline{sleep} to finish.
% Since the goal of the \lstinline{sh.c} is to be as simple as possible,
% it doesn't try to avoid creating interior processes.

管道在这种情况下似乎不比临时文件更强大：管道
\begin{lstlisting}[]
echo hello world | wc
\end{lstlisting}
可以不用管道实现为
\begin{lstlisting}[]
echo hello world >/tmp/xyz; wc </tmp/xyz
\end{lstlisting}
在这种情况下，管道相对于临时文件至少有三个优势。首先，管道自动清理自己；使用文件重定向时，shell 必须小心在完成时删除
\lstinline{/tmp/xyz}。
其次，管道可以传递任意长的数据流，而文件重定向需要足够的磁盘空闲空间来存储所有数据。第三，管道允许管道阶段的并行执行，而文件方法要求第一个程序在第二个程序开始之前完成。
% Fourth, if you are implementing inter-process communication,
% pipes' blocking reads and writes are more efficient
% than the non-blocking semantics of files.
%%
%%	File system
%%
\section{文件系统}

xv6 文件系统提供数据文件，包含未解释的字节数组，以及目录，包含对数据文件和其他目录的命名引用。目录形成一棵树，从一个称为
\indextext{root}（根）
的特殊目录开始。像
\lstinline{/a/b/c}
这样的
\indextext{path}（路径）
指的是根目录
\lstinline{/}
中名为
\lstinline{a}
的目录中名为
\lstinline{b}
的目录中名为
\lstinline{c}
的文件或目录。不以
\lstinline{/}
开头的路径相对于调用进程的
\indextext{current directory}（当前目录）
进行评估，可以使用
\lstinline{chdir}
系统调用更改。这两个代码片段打开相同的文件（假设所有涉及的目录都存在）：
\begin{lstlisting}[]
chdir("/a");
chdir("b");
open("c", O_RDONLY);

open("/a/b/c", O_RDONLY);
\end{lstlisting}
第一个片段将进程的当前目录更改为
\lstinline{/a/b}；
第二个既不引用也不更改进程的当前目录。

有系统调用来创建新文件和目录：
\lstinline{mkdir}
创建新目录，
带有
\lstinline{O_CREATE}
标志的
\lstinline{open}
创建新数据文件，
\lstinline{mknod}
创建设备文件。这个例子说明了所有三种：
\begin{lstlisting}[]
mkdir("/dir");
fd = open("/dir/file", O_CREATE|O_WRONLY);
close(fd);
mknod("/console", 1, 1);
\end{lstlisting}
\lstinline{mknod}
创建一个引用设备的特殊文件。与设备文件相关联的是主设备号和次设备号（
\lstinline{mknod}
的两个参数），它们唯一标识一个内核设备。当进程稍后打开设备文件时，内核将
\lstinline{read}
和
\lstinline{write}
系统调用转移到内核设备实现，而不是将它们传递给文件系统。

文件的名称与文件本身不同；相同的底层文件，称为
\indextext{inode}（inode），
可以有多个名称，称为
\indextext{links}（链接）。
每个链接由目录中的一个条目组成；该条目包含文件名和对 inode 的引用。Inode 保存有关文件的
\indextext{metadata}（元数据），
包括其类型（文件或目录或设备）、其长度、文件内容在磁盘上的位置，以及指向该文件的链接数。

\lstinline{fstat}
系统调用从文件描述符引用的 inode 检索信息。它填充一个
\lstinline{struct}
\lstinline{stat}，
在
\lstinline{stat.h} \fileref{kernel/stat.h}
中定义为：
\begin{lstlisting}[]
#define T_DIR     1   // 目录
#define T_FILE    2   // 文件
#define T_DEVICE  3   // 设备

struct stat {
  int dev;     // 文件系统的磁盘设备
  uint ino;    // Inode 号
  short type;  // 文件类型
  short nlink; // 指向文件的链接数
  uint64 size; // 文件大小（字节）
};
\end{lstlisting}

\lstinline{link}
系统调用创建另一个文件系统名称，引用与现有文件相同的 inode。这个片段创建一个新文件，名为
\lstinline{a}
和
\lstinline{b}。
\begin{lstlisting}[]
open("a", O_CREATE|O_WRONLY);
link("a", "b");
\end{lstlisting}
从
\lstinline{a}
读取或写入与从
\lstinline{b}
读取或写入相同。
每个 inode 由唯一的
\textit{inode}
\textit{号}
标识。在上述代码序列之后，可以通过检查
\lstinline{fstat}
的结果来确定
\lstinline{a}
和
\lstinline{b}
引用相同的底层内容：两者都将返回相同的 inode 号
（\lstinline{ino}），
并且
\lstinline{nlink}
计数将设置为 2。

\lstinline{unlink}
系统调用从文件系统中删除一个名称。只有当文件的链接计数为零且没有文件描述符引用它时，文件的 inode 和保存其内容的磁盘空间才会被释放。
因此，将
\begin{lstlisting}[]
unlink("a");
\end{lstlisting}
添加到最后一个代码序列会使 inode 和文件内容可以作为
\lstinline{b}
访问。此外，
\begin{lstlisting}[]
fd = open("/tmp/xyz", O_CREATE|O_RDWR);
unlink("/tmp/xyz");
\end{lstlisting}
是创建一个没有名称的临时 inode 的惯用方式，当进程关闭
\lstinline{fd}
或退出时将被清理。

Unix 提供可从 shell 调用的文件实用程序作为用户级程序，例如
\lstinline{mkdir}、
\lstinline{ln}
和
\lstinline{rm}。
这种设计允许任何人通过添加新的用户级程序来扩展命令行界面。事后看来这个计划似乎显而易见，但在 Unix 时代设计的其他系统通常将这种命令内置到 shell 中（并将 shell 内置到内核中）。

一个例外是
\lstinline{cd}，
它内置在 shell 中
\lineref{user/sh.c:/if.cmd.0..==..c./}。
\lstinline{cd}
必须更改 shell 本身的当前工作目录。如果
\lstinline{cd}
作为常规命令运行，那么 shell 会 fork 一个子进程，子进程会运行
\lstinline{cd}，
\lstinline{cd}
会更改
\textit{子进程}
的工作目录。父进程（即 shell）的工作目录不会改变。
%%
%%	Real world
%%
\section{现实世界}

Unix 的``标准''文件描述符、管道以及用于操作它们的方便的 shell 语法的结合，是编写通用可重用程序的重大进步。这一想法激发了一种``软件工具''文化，这种文化对 Unix 的强大和流行负有责任，shell 是第一个所谓的``脚本语言''。Unix 系统调用接口今天存在于 BSD、Linux 和 macOS 等系统中。

Unix 系统调用接口已通过可移植操作系统接口（POSIX）标准标准化。Xv6
\textit{不是}
POSIX 兼容的：它缺少许多系统调用（包括基本的如
\lstinline{lseek}），
并且它提供的许多系统调用与标准不同。我们对 xv6 的主要目标是简单和清晰，同时提供一个简单的类 Unix 系统调用接口。一些人通过添加更多系统调用和一个简单的 C 库来扩展 xv6，以运行基本的 Unix 程序。然而，现代内核提供了比 xv6 更多的系统调用，以及更多种类的内核服务。例如，它们支持网络、窗口系统、用户级线程、许多设备的驱动程序等。现代内核不断快速演进，并提供 POSIX 之外的功能。

Unix 用单一的文件名和文件描述符接口统一了对多种资源类型（文件、目录和设备）的访问。这个想法可以扩展到更多种类的资源；一个很好的例子是 Plan 9~\cite{Presotto91plan9}，它将``资源即文件''的概念应用于网络、图形等。然而，大多数 Unix 衍生的操作系统没有走这条路。

文件系统和文件描述符一直是强大的抽象。即便如此，还有其他操作系统接口模型。Multics，Unix 的前身，以一种使文件存储看起来像内存的方式抽象文件存储，产生了一种非常不同风格的接口。Multics 设计的复杂性对 Unix 的设计者产生了直接影响，他们的目标是构建更简单的东西。

Xv6 不提供用户或保护用户免受其他用户侵害的概念；用 Unix 术语来说，所有 xv6 进程都以 root 运行。

本书研究 xv6 如何实现其类 Unix 接口，但这些想法和概念适用于不仅仅是 Unix。任何操作系统都必须将进程多路复用到底层硬件上，隔离进程彼此，并提供受控的进程间通信机制。研究 xv6 后，你应该能够查看其他更复杂的操作系统，并在这些系统中看到 xv6 背后的概念。

%%
\section{练习}
%%

\begin{enumerate}

\item 编写一个程序，使用 UNIX 系统调用通过一对管道在两个进程之间``乒乓''传输一个字节，每个方向一个。测量程序的性能，以每秒交换次数为单位。

\end{enumerate}

%   terminology:
%     process: refers to execution in user space, or maybe struct proc &c
%     process memory: the lower part of the address space
%     process has one thread with two stacks (one for in kernel mode and one for
%     in user mode)
% talk a little about initial page table conditions:
%     paging not on, but virtual mostly mapped direct to physical,
%     which is what things look like when we turn paging on as well
%     since paging is turned on after we create first process.
%   mention why still have SEG_UCODE/SEG_UDATA?
%   do we ever really say what the low two bits of %cs do?
%     in particular their interaction with PTE_U
%   sidebar about why it is extern char[]

\chapter{操作系统组织结构}
\label{CH:FIRST}

操作系统的关键需求是同时支持多个活动。例如，可以使用
Chapter~\ref{CH:UNIX}
中的 {\tt fork} 和 {\tt exec} 系统调用（system call）来启动编译器和文本编辑器作为进程（process）。操作系统必须在多个进程之间
\indextext{time-share}
分时（time-sharing）共享 CPU 和内存等资源。操作系统还必须安排
\indextext{isolation}
隔离（isolation）机制，使各进程相互独立。如果一个进程出现错误并发生故障，不应影响其他无关进程。然而，完全隔离又过于严格，因为进程应该能够有意地进行交互；管道（pipeline）就是一个例子。因此，操作系统必须满足三个要求：多路复用（multiplexing）、隔离（isolation）和交互（interaction）。

本章概述了操作系统如何组织以实现这三个要求。实现方式有很多，但本文聚焦于围绕
\indextext{monolithic kernel}
宏内核（monolithic kernel）的主流设计，许多 Unix 操作系统都采用这种架构。本章还概述了 xv6 进程，即 xv6 中的隔离单元。

Xv6 运行在
\indextext{multi-core}
多核（multi-core）RISC-V 微处理器上，
其许多底层功能（例如进程实现）是特定于 RISC-V 的。RISC-V 是 64 位 CPU，xv6 使用 ``LP64'' C 语言编写，这意味着 C 语言中的 long（L）和指针（P）是 64 位，但 {\tt int} 是 32 位。本书假设读者在某些架构上有一定的机器级编程经验，并将在需要时介绍 RISC-V 特定的概念。
用户级 ISA~\cite{riscv:user} 和特权架构~\cite{riscv:priv} 文档是完整的规范。
读者还可以参考
``The RISC-V Reader: An Open Architecture
Atlas''~\cite{riscv}。

完整计算机中的 CPU 被支持硬件包围，其中大部分以 I/O 接口的形式存在。Xv6 为 qemu 的 ``-machine virt'' 选项模拟的支持硬件而编写。这包括 RAM、包含启动代码的 ROM、与用户键盘/屏幕的串行连接，以及用于存储的磁盘。

%%
\section{抽象物理资源}
%%

人们遇到操作系统时可能会问的第一个问题是：为什么需要操作系统？也就是说，可以将
Figure~\ref{fig:api}
中的系统调用实现为一个库，应用程序与之链接。在这种方案中，每个应用程序甚至可以有自己的库，根据其需求定制。应用程序可以直接与硬件资源交互，并以最适合应用的方式使用这些资源（例如，实现高性能或可预测的性能）。一些用于嵌入式设备或实时系统的操作系统就是以这种方式组织的。

这种库方法的缺点是，如果有多个应用程序运行，这些应用程序必须表现良好。例如，每个应用程序必须定期放弃 CPU，以便其他应用程序可以运行。如果所有应用程序相互信任且没有错误，这种
\textit{协作式}
分时方案可能是可行的。更典型的情况是应用程序互不信任，且存在错误，因此通常需要比协作方案更强的隔离。

为了实现强隔离，最好禁止应用程序直接访问敏感的硬件资源，而是将资源抽象为服务。例如，Unix 应用程序仅通过文件系统的
\lstinline{open}、
\lstinline{read}、
\lstinline{write}
和
\lstinline{close}
系统调用与存储交互，而不是直接读写磁盘。这为应用程序提供了路径名的便利，并允许操作系统（提供该接口）管理磁盘。即使隔离不是关注点，有意交互（或只是希望互不干扰）的程序也可能发现文件系统是比直接使用磁盘更方便的抽象。

同样，Unix 透明地在进程之间切换硬件 CPU，根据需要保存和恢复寄存器状态，使应用程序无需感知分时。这种透明性允许操作系统共享 CPU，即使某些应用程序处于无限循环中。

再举一个例子，Unix 进程使用
\lstinline{exec}
构建其内存映像，而不是直接与物理内存交互。这允许操作系统决定将进程放置在内存的何处；如果内存紧张，操作系统甚至可能将进程的某些数据存储在磁盘上。
\lstinline{exec}
还为用户提供了使用文件系统存储可执行程序映像的便利。

Unix 进程之间的许多交互形式通过文件描述符（file descriptor）进行。文件描述符不仅抽象了许多细节（例如，管道或文件中的数据存储位置），而且其定义方式也简化了交互。例如，如果管道中的一个应用程序退出或失败，内核会自动为管道中的下一个进程生成文件结束信号。

Figure~\ref{fig:api}
中的系统调用接口经过精心设计，既提供了程序员便利，又提供了强隔离的可能性。Unix 接口不是抽象资源的唯一方式，但已被证明是一种良好的方式。

%%
\section{用户模式、监管者模式和系统调用}
%%

强隔离需要在应用程序和操作系统之间建立严格的边界。即使应用程序有错误或恶意行为，也不应允许其干扰操作系统或其他程序的运行。为了实现强隔离，操作系统必须确保应用程序不能修改（甚至不能读取）操作系统的数据结构和指令，并且应用程序不能访问其他进程的内存。

CPU 为强隔离提供硬件支持。例如，RISC-V 有三个特权级别（privilege level），限制代码可以执行的操作：
\indextext{machine mode}（机器模式）、
\indextext{supervisor mode}（监管者模式）和
\indextext{user mode}（用户模式）。在机器模式下执行的指令拥有完全特权；CPU 启动时处于机器模式。机器模式主要用于在启动期间配置计算机。Xv6 在机器模式下执行一小段时间，然后切换到监管者模式。

在监管者模式下，CPU 被允许执行
\indextext{privileged instructions}（特权指令）：例如，启用和禁用中断、读写保存页表地址的寄存器等。如果用户模式中的应用程序尝试执行特权指令，CPU 不会执行该指令，而是会``陷入''（trap）到监管者模式中的特殊代码，该代码可以终止应用程序。
Chapter~\ref{CH:UNIX}
中的 Figure~\ref{fig:os} 说明了这种组织结构。应用程序只能执行用户模式指令（例如，数字相加等），被称为在
\indextext{user space}（用户空间）
中运行，而监管者模式中的软件还可以执行特权指令，被称为在
\indextext{kernel space}（内核空间）
中运行。在内核空间（或监管者模式）中运行的软件称为
\indextext{kernel}（内核）。

应用程序通过系统调用（system call）与内核交互，例如 {\tt read}。应用程序不允许直接调用内核函数或访问内核内存。RISC-V 提供了 \indexcode{ecall} 指令用于系统调用；它将 CPU 从用户模式切换到监管者模式，并跳转到内核指定的入口点。一旦 CPU 切换到监管者模式，内核就可以验证系统调用的参数（例如，检查传递给系统调用的地址是否是应用程序内存的一部分），决定应用程序是否被允许执行请求的操作（例如，检查应用程序是否被允许写入指定文件），然后拒绝或执行该操作。内核控制进入监管者模式的入口点非常重要；如果应用程序可以决定内核入口点，恶意应用程序可能会在内核跳过参数验证的位置进入，例如。

%%
\section{内核组织结构}
%% 

一个关键的设计问题是操作系统的哪部分应该在监管者模式下运行。一种可能是整个操作系统驻留在内核中，这样所有系统调用的实现都在监管者模式下运行。这种组织称为
\indextext{monolithic kernel}（宏内核）。

在宏内核组织中，整个操作系统由一个在监管者模式下运行的单一程序组成。这种组织方便的一个原因是操作系统设计者不需要将代码划分为需要监管者特权和不需要的部分。此外，操作系统不同部分之间的协作很容易，因为它们是同一程序的一部分。例如，宏内核可能与文件系统和虚拟内存系统共享磁盘块缓存。

缺点是宏内核往往变得庞大而复杂，以至于没有一个开发者理解代码不同部分之间的所有交互；这是产生错误的根源。内核中的错误特别麻烦，因为它可能导致整个计算机崩溃，或导致许多应用程序故障，或使整个计算机容易受到安全攻击。

\indextext{microkernel}（微内核）
旨在减少内核中错误的发生。其想法是将绝对最少的功能放入内核本身，这样少量代码在监管者模式下执行，内核易于理解和分析其正确性。操作系统的大部分作为用户级服务器进程运行。例如，文件系统代码将作为服务器进程执行，在用户模式而非监管者模式下。

\begin{figure}[t]
\center
\includegraphics[scale=0.5]{fig/mkernel.pdf}
\caption{A microkernel with a file-system server}
\label{fig:mkernel}
\end{figure}

Figure~\ref{fig:mkernel}
说明了这种微内核设计。在图中，文件系统作为用户级服务器进程运行。为了允许应用程序与文件服务器交互，内核提供了一种进程间通信（inter-process communication）机制，用于将消息从一个用户模式进程发送到另一个用户模式进程。例如，如果像 shell 这样的应用程序想要读取或写入文件，它会向文件服务器发送消息并等待响应。

在微内核中，内核接口由一些低级函数组成，用于启动应用程序、发送消息、访问设备硬件等。这种组织允许内核相对简单，因为大部分操作系统驻留在用户级服务器中。

在现实世界中，宏内核和微内核都很流行。许多 Unix 内核是宏内核。例如，Linux 具有宏内核，尽管某些操作系统功能作为用户级服务器运行（例如，窗口系统）。Linux 为操作系统密集型应用程序提供高性能，部分原因是内核的子系统可以紧密集成。

Minix、L4 和 QNX 等操作系统被组织为带有服务器的微内核，并已在嵌入式环境中广泛部署。L4 的一个变体 seL4 足够小，已被验证具有内存安全性和其他安全属性~\cite{sel4}。

关于哪种组织更好的争论在操作系统开发者之间存在很大争议，但没有确凿证据表明哪一种更好。此外，这很大程度上取决于``更好''的含义：更快的性能、更小的代码大小、内核的可靠性、完整操作系统的可靠性（包括用户级服务）等。

还有一些实际考虑可能比组织问题更重要。某些操作系统具有微内核，但出于性能原因，某些用户级服务在内核空间中运行。某些操作系统具有宏内核，因为它们就是这样开始的，并且几乎没有动力转向纯微内核组织，因为新功能可能比重写现有操作系统以适应微内核设计更重要。

从本书的角度来看，微内核和宏内核操作系统共享许多关键思想。它们实现系统调用，使用页表（page table），处理中断（interrupt），支持进程，使用锁（lock）进行并发控制，实现文件系统（file system）等。本书专注于这些核心思想。

Xv6 被实现为宏内核，与大多数 Unix 操作系统一样。因此，xv6 内核接口对应于操作系统接口，内核实现完整的操作系统。由于 xv6 不提供许多服务，其内核比某些微内核更小，但从概念上讲 xv6 是宏内核。  

\section{代码：xv6 组织结构}

\begin{figure}[t]
\center
\begin{tabular}{l|l|l}
& {\bf 文件} & {\bf 描述} \\
\midrule
启动 & entry.S & 最初的启动指令。 \\
 & main.c & 控制其他模块的初始化。 \\
 & start.c & 早期的机器模式启动代码。 \\
进程 & exec.c & exec() 系统调用。 \\
 & proc.c & 进程和调度。 \\
 & swtch.S & 线程切换。 \\
 & sysproc.c & 进程相关的系统调用。 \\
陷入 & kernelvec.S & 处理来自内核代码的陷入。 \\
 & trampoline.S & 处理来自用户代码的陷入。 \\
 & trap.c & 处理和从陷入与中断返回的 C 代码。 \\
 & syscall.c & 将系统调用分派给处理函数。 \\
内存 & vm.c & 管理页表和地址空间。 \\
 & kalloc.c & 物理页分配器。 \\
设备 & console.c & 连接用户键盘和屏幕。 \\
 & plic.c & RISC-V 中断控制器。 \\
 & printf.c & 格式化输出到控制台。 \\
 & uart.c & 串口控制台设备驱动。 \\
 & virtio\_disk.c & 磁盘设备驱动。 \\
文件系统 & bio.c & 文件系统的磁盘块缓存。 \\
 & file.c & 文件描述符支持。 \\
 & fs.c & 文件系统。 \\
 & log.c & 文件系统日志和崩溃恢复。 \\
 & sysfile.c & 文件相关的系统调用。 \\
 & pipe.c & 管道。 \\
其他 & sleeplock.c & 让出 CPU 的锁。 \\
 & spinlock.c & 不让出 CPU 的锁。 \\
 & string.c & C 字符串和字节数组库。 \\
\end{tabular}
\caption{Xv6 内核源文件。}
\label{fig:source}
\end{figure}

xv6 内核源代码位于 {\tt kernel/} 子目录中。
Figure~\ref{fig:source} 列出了这些文件，按内核的主要职责领域划分：启动系统（booting）、创建和控制进程、处理陷入（中断和系统调用）、分配内存和配置虚拟地址、控制设备，以及管理文件系统。

%% 
\section{Process overview}
%% 

The unit of isolation in xv6 (as in other Unix operating systems) is a 
\indextext{process}.
The process abstraction prevents one process from wrecking or spying on
another process's memory, CPU, file descriptors, etc.  It also prevents a process
from wrecking the kernel itself, so that a process can't subvert the kernel's
isolation mechanisms.
The kernel must implement the process abstraction with care because
a buggy or malicious application may trick the kernel or hardware into doing
something bad (e.g., circumventing isolation).  The mechanisms used by
the kernel to implement processes include the user/supervisor mode flag, address spaces,
and time-slicing of threads.

To help enforce isolation, the process abstraction provides the
illusion to a program that it has its own private machine.  A process provides
a program with what appears to be a private memory system, or
\indextext{address space}, 
which other processes cannot read or write.
A process also provides the program with what appears to be its own
CPU to execute the program's instructions.

Xv6 uses page tables (which are implemented by hardware) to give each process
its own address space. The RISC-V page table
translates (or ``maps'') a
\indextext{virtual address}
(the address that an RISC-V instruction manipulates) to a
\indextext{physical address}
(an address that the CPU sends to main memory).

\begin{figure}[t]
\centering
\includegraphics[scale=0.5]{fig/as.pdf}
\caption{Layout of a process's virtual address space}
\label{fig:as}
\end{figure}

Xv6 maintains a separate page table for each process that defines that process's
address space. As illustrated in 
Figure~\ref{fig:as},
an address space includes the process's
\indextext{user memory}
starting at virtual address zero. Instructions come first,
followed by global variables, then the stack,
and finally a ``heap'' area (for malloc)
that the process can expand as needed.
There are a number of factors that limit the
maximum size of a process's address space:
pointers on the RISC-V are 64 bits wide;
the hardware uses only the low 39 bits when
looking up virtual addresses in page tables;
and xv6 uses only 38 of those 39 bits.
Thus, the maximum address is $2^{38}-1$ =
0x3fffffffff, which is \lstinline{MAXVA}~\lineref{kernel/riscv.h:/define.MAXVA/}.
At the top of the address space xv6 places
a \indextext{trampoline} page (4096 bytes)
and a \indextext{trapframe} page.  Xv6 uses these two pages
to transition into the kernel and back;
the trampoline page contains the code to transition in and out
of the kernel, and the trapframe is where the kernel saves
the process's user registers,
as Chapter~\ref{CH:TRAP} explains.

The xv6 kernel maintains many pieces of state for each process,
which it gathers into a
\indexcode{struct proc}
\lineref{kernel/proc.h:/^struct.proc/}.
A process's most important pieces of kernel state are its 
page table, its kernel stack, and its run state.
We'll use the notation
\indexcode{p->xxx}
to refer to elements of the
\lstinline{proc}
structure; for example,
\indexcode{p->pagetable} is a pointer to the process's page table.

At this point, please read {\tt kernel/proc.h}, which defines {\tt
  struct proc}. The xv6 code is more important for you to understand
than this book; you should prioritize the code, and consult this book
as needed to clarify the code. The purpose of some of the code may not
be apparent at first, but further reading and searching the code will
help. Feel free to explore and modify the code.

Each process has a thread of control (or 
\indextext{thread}
for short) that holds the state needed to execute the process.
At any given time, a thread might be executing on a CPU,
or suspended (not executing, but capable of resuming executing
in the future).
To switch a CPU between processes,
the kernel suspends the thread currently running on that CPU
and saves its state,
and restores the state of another process's previously-suspended
thread.  Much of the state of a thread (local variables, function call return
addresses) is stored on the thread's stacks.
Each process has two stacks: a user stack and a kernel stack
(\indexcode{p->kstack}).
When the process is executing user instructions, only its user stack
is in use, and its kernel stack is empty.
When the process enters the kernel (for a system call or interrupt),
the kernel code executes on the process's kernel stack; while
a process is in the kernel, its user stack still contains saved
data, but isn't actively used.
A process's thread alternates between actively using its user stack
and its kernel stack. The kernel stack is separate (and protected from
user code) so that the kernel
can execute even if a process has wrecked its user stack.

A process's user code
can make a system call by executing the RISC-V \indexcode{ecall}
instruction. This instruction switches to supervisor mode and
changes the program counter to a kernel-defined entry point.  The code
at the entry point switches to the process's kernel stack and executes the kernel
instructions that implement the system call.  When the system call
completes, the kernel returns to
user space by executin the \indexcode{sret} instruction, which switches
to user mode and resumes executing user instructions
just after the system call instruction.  A process's thread can
``block'' in the kernel to wait for I/O, and resume where it left off
when the I/O has finished.

\indexcode{p->state} 
indicates whether the process is allocated, ready
to run, currently running on a CPU, waiting for I/O, or exiting.

\indexcode{p->pagetable}
holds the process's page table, in the format
that the RISC-V hardware expects.
Xv6 causes the paging hardware to use a process's
\lstinline{p->pagetable}
when executing that process in user space.
A process's page table also serves as the record of the
addresses of the physical pages allocated to store the process's
memory.

In summary, a process bundles two design ideas: an address space to
give a process the illusion of its own memory, and a thread to give
the process the illusion of its own CPU.  In xv6, a process consists
of one address space and one thread.  In real operating systems a
process may have more than one thread to take advantage of multiple CPUs.

%% 
\section{Code: starting xv6, the first process and system call}
%% 
To make xv6 more concrete, we'll outline how the kernel starts and
runs the first process. The subsequent chapters will describe the
mechanisms that show up in this overview in more detail.
Please read {\tt kernel/entry.S}, {\tt kernel/start.c},
{\tt kernel/main.c}, and {\tt user/init.c}.

When the RISC-V computer powers on, it initializes
itself and runs a boot loader which is stored in read-only
memory.  The boot loader copies the xv6 kernel into memory
at physical address
\texttt{0x80000000}.
The reason it places the kernel at
\texttt{0x80000000}
rather than
\texttt{0x0}
is because the address range
\texttt{0x0:0x80000000}
contains I/O devices.

Then the boot loader jumps to xv6 starting at
\indexcode{_entry}
\lineref{kernel/entry.S:/^.entry:/}.
The RISC-V starts with paging hardware disabled:
virtual addresses map directly to physical addresses.
The instructions at
\lstinline{_entry}
set up a stack so that xv6 can run C code.
Xv6 declares space for this stack,
\lstinline{stack0},
in the file
\lstinline{start.c}
\lineref{kernel/start.c:/stack0/}.
The code at
\lstinline{_entry}
loads the stack pointer register
\texttt{sp}
with the address
\lstinline{stack0+4096},
the top of the stack, because the stack
on RISC-V grows down.
Now that the kernel has a stack,
\lstinline{_entry}
calls into C code at
\lstinline{start}
\lineref{kernel/start.c:/^start/}.

The function
\lstinline{start}
performs some setup that the CPU only allows in machine mode,
most crucially programming the clock chip to generate
timer interrupts.
Then {\tt start} uses the RISC-V {\tt mret}
instruction to switch to supervisor mode and
jump to 
\lstinline{main}
\lineref{kernel/main.c:/^main/}.
{\tt mret} requires
a bit of setup:
{\tt start} sets the previous privilege mode to
supervisor in the register
\lstinline{mstatus},
sets the destination address to
\lstinline{main}
by writing
\lstinline{main}'s
address into
the register
\lstinline{mepc},
disables virtual address translation in supervisor mode
by writing
\lstinline{0}
into the page-table register
\lstinline{satp},
and delegates all interrupts and exceptions
to supervisor mode.

After
\lstinline{main}
\lineref{kernel/main.c:/^main/}  
initializes several devices and subsystems, 
it creates the first process by calling 
\lstinline{userinit}
\lineref{kernel/proc.c:/^userinit/}.
All newly created processes start executing in the kernel in
\lstinline{forkret}
\lineref{kernel/proc.c:/^forkret/}.
As a special case for the first process, {\tt forkret}
calls {\tt kexec} to load the user program {\tt /init}.

After calling
\lstinline{kexec},
{\tt forkret} returns to user space in
the \lstinline{/init} process.
\lstinline{init}
\lineref{user/init.c:/^main/}
creates a new console device file
if needed
and then opens it as file descriptors 0, 1, and 2.
Then it starts a shell on the console.
The system is up.

\section{Security Model}

You may wonder how the operating system deals with buggy or malicious
code. Because coping with malice is strictly harder than dealing with
accidental bugs, it's reasonable to 
focus mostly on providing security against malice.
Here's a high-level view of typical security assumptions and
goals in operating system design.

The operating system must assume that a process's user-level code will
do its best to wreck the kernel or other processes. User code may try
to dereference pointers outside its allowed address space; it may
attempt to execute instructions not intended
for user code; it may try to read and write RISC-V control
registers; it may try to access device hardware;
and it may pass clever values to system calls in an attempt
to trick the kernel into crashing or doing something stupid.

The
kernel's goal is to restrict each user process so that it can only
access its own user memory, use the 32 general-purpose
RISC-V registers, and affect the kernel and other processes in the
ways that system calls are intended to allow. The kernel must prevent
any other actions. These are typically absolute requirements in kernel
design.

Expectations for the kernel's own code are different. Kernel
code is assumed to be written by well-meaning and careful programmers,
to be bug-free, and to contain
nothing malicious. This assumption affects how we analyze kernel code.
For example, there are many internal kernel functions (e.g., the spin
locks) that would cause serious problems if kernel code used them
incorrectly. We assume,
however, that the kernel uses its own functions correctly.
At the hardware level, the RISC-V CPU, RAM, disk, etc. are
assumed to operate as advertised in the documentation, with no
hardware bugs.

Real life is not so straightforward. It's
difficult to prevent abusive user programs from 
calling system calls in a way that makes the system unusable
by consuming kernel-protected resources:
disk space, CPU time, process table slots, etc. It's usually
impossible to write 100\% bug-free kernel code or design bug-free hardware; if the
writers of malicious user code are aware of kernel or hardware bugs,
they will exploit them. Even in mature, widely-used kernels, such as
Linux, people often discover previously-unknown
vulnerabilities~\cite{mitre:cves}.
Finally, the distinction between
user and kernel code is sometimes blurred: some privileged user-level
processes may provide essential services and effectively be part of
the operating system, and in some operating systems privileged user
code can insert new code into the kernel (as with Linux's loadable
kernel modules and eBPF).

As a partial defense against kernel bugs, xv6 code includes checks for
inconsistencies and unrecoverable errors, and will ``panic'' in
response, by calling {\tt panic()}. This function prints an error
message and halts the system. Panicking is not desirable, but is
preferable to continuing execution. Typically a panic results from a
kernel bug that causes kernel data to be incorrect or causes the
kernel to perform an illegal action such as referencing non-existent
memory; in such a situation it is safer to halt execution with {\tt
  panic()} than to try to continue in an inconsistent state. A kernel
developer would react to a panic by working to identify and fix the
underlying code bug.

%% 
\section{Real world}
%% 

Most operating systems have adopted the process concept, and most
processes look similar to xv6's.  Modern operating systems, however,
support several threads within a process, to allow a single process to
exploit multiple CPUs.  Supporting multiple threads in a
process involves quite a bit of machinery that xv6 doesn't have,
often including interface changes (e.g., Linux's
\lstinline{clone},
a variant of
\lstinline{fork}),
to control which aspects of
a process threads share.
%% 
\section{Exercises}
%%

\begin{enumerate}

\item Add a system call to xv6
  that returns the amount of free memory available.

% break *0x3ffffff000
% disas 0x3ffffff000,+8
% set disassemble-next-line auto
% x/i $pc
% break *0x3ffffff10e
% print/x $sepc
% print/x $pc

\end{enumerate}

% talk a little about initial page table conditions:
%     paging not on, but virtual mostly mapped direct to physical,
%     which is what things look like when we turn paging on as well
%     since paging is turned on after we create first process.
%   mention why still have SEG_UCODE/SEG_UDATA?
%   do we ever really say what the low two bits of %cs do?
%     in particular their interaction with PTE_U
%   sidebar about why it is extern char[]
\chapter{页表}
\label{CH:MEM}

页表（page table）是操作系统为每个进程提供私有地址空间和内存的最流行机制。页表决定内存地址的含义，以及哪些物理内存区域可被访问。它们允许 xv6 隔离不同进程的地址空间，并将其多路复用到单个物理内存上。页表提供了一层间接性，使操作系统能够执行许多有用的技巧。Xv6 实现了一些：在多个地址空间中映射相同的内存（trampoline 页），用未映射的页保护内核和用户栈，以及延迟分配用户堆内存。本章的其余部分将解释 RISC-V 硬件提供的页表以及 xv6 如何使用它们。

%%
\section{分页硬件}
%%
提醒一下，RISC-V 指令（用户态和内核态）操作虚拟地址（virtual address）。机器的 RAM（即物理内存，physical memory）使用物理地址索引。RISC-V 页表硬件通过将每个虚拟地址映射到物理地址来连接这两种地址。

\begin{figure}[t]
\centering
\includegraphics[scale=0.5]{fig/riscv_address.pdf}
\caption{平坦页表将虚拟地址映射到物理地址的抽象视图。}
\label{fig:riscv_address}
\end{figure}

Xv6 使用 RISC-V 的 Sv39 模式，这意味着 64 位虚拟地址只有低 39 位被使用，高 25 位不使用。在此 Sv39 配置中，RISC-V 页表逻辑上是 $2^{27}$ (134,217,728) 个 \indextext{page table entries (PTEs)}（页表项）的数组。每个 PTE 包含一个 44 位的物理页号（PPN）和一些标志。分页硬件使用 39 位中的高 27 位索引到页表中找到 PTE，并生成一个 56 位的物理地址，其高 44 位来自 PTE 中的 PPN，低 12 位从原始虚拟地址复制。
Figure~\ref{fig:riscv_address}
显示了此过程，页表的逻辑视图是 PTE 的简单数组（RISC-V 页表实际上是一棵树；参见 Figure~\ref{fig:riscv_pagetable} 了解完整故事）。页表使操作系统能够以 4096 ($2^{12}$) 字节对齐块的粒度控制虚拟到物理地址转换。这样的块称为 \indextext{page}（页）。

\begin{figure}[t]
\centering
\includegraphics[scale=0.5]{fig/riscv_pagetable.pdf}
\caption{RISC-V 地址转换细节。}
\label{fig:riscv_pagetable}
\end{figure}

RISC-V 的设计为虚拟地址和物理地址的扩展留下了空间。如果需要更多虚拟地址空间，RISC-V 支持 Sv48 模式，具有 48 位虚拟地址~\cite{riscv:priv}。物理地址也有扩展空间：PTE 格式中有空间让物理页号再增长 10 位。RISC-V 的设计者根据技术预测选择了地址大小。$2^{48}$ 字节即 262,144 GB，远大于当今任何应用程序可能使用的用户虚拟地址空间。$2^{56}$ 字节即 65,536 TB 的物理地址空间，远超目前任何计算机可配备的 RAM。

如
Figure~\ref{fig:riscv_pagetable}
所示，RISC-V CPU 页表作为三级树存储在物理内存中。树的根是一个 4096 字节的页表页，包含 512 个 PTE，其中包含树下一级页表页的物理地址。这些页中的每一个都包含树最后一级的 512 个 PTE。分页硬件使用 27 位中的高 9 位在根页表页中选择 PTE，中间 9 位在树下一级的页表页中选择 PTE，低 9 位选择最终的 PTE。（在 Sv48 RISC-V 中，页表有四级，虚拟地址的 39 到 47 位索引到顶级。）

如果转换地址所需的三个 PTE 中任何一个不存在，分页硬件会引发 \indextext{page-fault exception}（页错误异常），由内核处理页错误（参见 Chapters~\ref{CH:TRAP} 和~\ref{CH:PGFAULTS}）。

与 Figure~\ref{fig:riscv_address} 的单级设计相比，Figure~\ref{fig:riscv_pagetable} 的三级结构以更节省内存的方式记录 PTE。在虚拟地址大范围无映射的常见情况下，三级结构可以省略整个页目录。例如，如果应用程序只使用从地址零开始的少数几页，那么顶级页目录的条目 1 到 511 无效，内核无需为这 511 个中间页目录分配页面。此外，内核也无需为这 511 个中间页目录的底层页目录分配页面。因此，在此例中，三级设计节省了 511 个中间页目录页面和 $511\times512$ 个底层页目录页面。

尽管 CPU 在执行加载或存储指令时会以硬件方式遍历三级结构，但三级的一个潜在缺点是 CPU 必须从内存加载三个 PTE 才能将加载/存储指令中的虚拟地址转换为物理地址。为了避免从物理内存加载 PTE 的开销，RISC-V CPU 将页表项缓存在
\indextext{Translation Look-aside Buffer (TLB)}（转换后备缓冲器）中。

每个 PTE 包含标志位，告知分页硬件如何使用关联的虚拟地址。
\indexcode{PTE_V}
指示 PTE 是否存在：若未设置，对页面的引用会导致页错误（即不允许）。
\indexcode{PTE_R}
控制是否允许指令读取页面。
\indexcode{PTE_W}
控制是否允许指令写入页面。
\indexcode{PTE_X}
控制 CPU 是否可将页面内容解释为指令并执行。
\indexcode{PTE_U}
控制用户模式中的指令是否允许访问页面；若
\indexcode{PTE_U}
未设置，PTE 只能在监管者模式中使用。
Figure~\ref{fig:riscv_pagetable}
显示了标志位在 PTE 中的位置。标志和所有其他与页面硬件相关的结构定义在
\fileref{kernel/riscv.h}
中。

为了让 CPU 使用页表，内核必须将根页表页的物理地址写入
\texttt{satp}\index{satp@\lstinline{satp}}
寄存器。CPU 将使用其 \texttt{satp} 指向的页表翻译后续指令生成的所有地址。每个 CPU 都有自己的 \texttt{satp}，因此不同的 CPU 可以运行不同的进程，每个进程都有自己的页表描述的私有地址空间。

从内核的角度看，页表是存储在内存中的数据，内核使用类似于树形数据结构的代码创建和修改页表。

关于本书中使用的一些术语的说明。\textit{物理内存}（Physical memory）指 RAM 中的存储单元。物理内存的字节有一个地址，称为 \textit{物理地址}（physical address）。解引用地址的指令（如加载、存储、跳转和函数调用）只使用虚拟地址（virtual address），分页硬件将其转换为物理地址，然后发送到 RAM 硬件读取或写入存储。\textit{地址空间}（address space）是给定页表中有效的虚拟地址集合；每个 xv6 进程有独立的用户地址空间，xv6 内核也有自己的地址空间。\textit{用户内存}（User memory）指进程的用户地址空间加上页表允许进程访问的物理内存。\textit{虚拟内存}（Virtual memory）指与页表管理及使用它们实现隔离等相关的概念和技术。

\begin{figure}[h]
\centering
 \includegraphics[scale=0.65]{fig/xv6_layout.pdf}
\caption{左侧，xv6 的内核虚拟地址空间。
{\sf \small{RWX}}
指 PTE 的读、写和执行权限。
右侧，xv6 期望看到的 RISC-V 物理地址空间。}
\label{fig:xv6_layout}
\end{figure}

%%
\section{内核地址空间}
%%
启动时，xv6 创建一个描述内核地址空间的页表。内核配置其地址空间布局，使自己能够访问物理内存和各种硬件资源的可预测虚拟地址。
Figure~\ref{fig:xv6_layout}
显示了此布局如何将内核虚拟地址映射到物理地址。文件
\fileref{kernel/memlayout.h}
声明了 xv6 内核内存布局的常量。

QEMU 模拟一台计算机，包括从物理地址 \texttt{0x80000000} 开始并至少持续到 \texttt{0x88000000} 的 RAM（物理内存），xv6 称之为 \texttt{PHYSTOP}。QEMU 模拟还包括 I/O 设备，如磁盘接口。QEMU 将设备接口作为位于物理地址空间 \texttt{0x80000000} 以下的 \indextext{memory-mapped}（内存映射）控制寄存器暴露给软件。内核可以通过读写这些特殊物理地址与设备交互；此类读写与设备硬件通信而非 RAM。Chapter~\ref{CH:TRAP} 解释 xv6 如何与设备交互。

内核将所有物理 RAM 和设备寄存器映射到等于物理地址的虚拟地址。这称为``直接映射''，允许内核通过简单地加载或存储虚拟地址 $x$ 来读取或写入物理地址 $x$。内核代码本身位于虚拟地址空间和物理内存中的 \lstinline{KERNBASE=0x80000000}。当 \lstinline{kfork}
\lineref{kernel/proc.c:/^kfork/}
为子进程分配用户内存时，分配器返回该内存的物理地址；\lstinline{kfork} 在将父进程的用户内存复制到子进程时直接使用该地址作为虚拟地址。

有几个内核虚拟地址不是直接映射的：

\begin{itemize}
  
\item trampoline 页。它映射在虚拟地址空间的顶部；用户页表具有相同的映射。Chapter~\ref{CH:TRAP} 讨论了 trampoline 页的作用，但我们在这里看到了页表的一个有趣用例：一个物理页（保存 trampoline 代码）在内核的虚拟地址空间中映射了两次：一次在虚拟地址空间的顶部，一次使用直接映射。

\item 内核栈页。每个进程都有自己的内核栈，映射在较高的内核虚拟地址处，因此在其下方 xv6 可以保留一个未映射的\indextext{保护页}。保护页的 PTE 是无效的（即\lstinline{PTE_V}未设置），因此如果内核溢出内核栈，很可能会引发页故障，导致内核崩溃。如果没有保护页，溢出的栈将覆盖其他内核内存，导致错误操作。崩溃重启是可取的。

\end{itemize}

虽然内核通过高内存映射使用其栈，但每个栈也可通过直接映射地址供内核访问。另一种设计可能只有直接映射，并在直接映射地址使用栈。然而，在这种安排中，提供保护页将涉及取消映射否则会引用物理内存的虚拟地址，这将很难使用。

内核将 trampoline 和内核文本页映射为具有 \lstinline{PTE_R} 和 \lstinline{PTE_X} 权限，但没有 \lstinline{PTE_W}。内核将其他页映射为具有 \lstinline{PTE_R} 和 \lstinline{PTE_W} 权限，但没有 \lstinline{PTE_X}。保护页映射是无效的。这些受限权限旨在帮助捕获以意外方式访问页的内核错误，例如若内核代码意外尝试写入内核指令。

内核创建一个单一的内核页表，所有 CPU 在内核中执行时使用它。xv6 在最初创建后不会修改内核页表。

%%
\section{代码：创建地址空间}
%%

请在继续之前阅读 {\tt kernel/vm.c} 到 {\tt mappages()} 的结尾。

大部分用于操作地址空间和页表的 xv6 代码位于 {\tt vm.c}
\fileref{kernel/vm.c}。
中心数据结构是 {\tt pagetable\_t}，它实际上是指向 RISC-V 根页表页的指针；{\tt pagetable\_t} 可以是内核页表或每个进程的页表之一。中心函数是 {\tt walk}，它查找虚拟地址的 PTE，以及 {\tt mappages}，它安装新映射的 PTE。以 {\tt kvm} 开头的函数操作内核页表；以 {\tt uvm} 开头的函数操作用户页表；其他函数用于两者。{\tt copyout} 和 {\tt copyin} 将数据复制到作为系统调用参数提供的用户虚拟地址或从该地址复制；它们在 {\tt vm.c} 中，因为它们需要显式转换这些地址以找到相应的物理内存。

在启动序列的早期，
\indexcode{main}
调用
\indexcode{kvminit}
\lineref{kernel/vm.c:/^kvminit/}
使用
\indexcode{kvmmake}
\lineref{kernel/vm.c:/^kvmmake/}
创建内核的页表。此调用发生在 xv6 在 RISC-V 上启用分页之前，因此地址直接引用物理内存。
\lstinline{kvmmake}
首先分配一页物理内存来保存根页表页。然后它调用
\indexcode{kvmmap}
来安装内核需要的转换。转换包括内核的指令和数据、高达
\indexcode{PHYSTOP}
的物理内存，以及实际上是设备的内存范围。
\indexcode{proc\_mapstacks}
\lineref{kernel/proc.c:/^proc\_mapstacks/}
为每个进程分配一个内核栈。它调用 \lstinline{kvmmap} 将每个栈映射到由
\lstinline{KSTACK}
生成的虚拟地址，这为无效的栈保护页留出了空间。

\indexcode{kvmmap} 
\lineref{kernel/vm.c:/^kvmmap/}
calls
\indexcode{mappages}
\lineref{kernel/vm.c:/^mappages/}，
它在页表中为一组虚拟地址安装映射到相应的物理地址范围。它以页间隔对范围内的每个虚拟地址执行此操作。对于要映射的每个虚拟地址，
\lstinline{mappages}
调用
\indexcode{walk}
来查找该地址的 PTE 地址。然后它初始化 PTE 以保存相关的物理页号、所需的权限
（\lstinline{PTE_W}、
\lstinline{PTE_X}
和/或
\lstinline{PTE_R}），
以及
\lstinline{PTE_V}
以将 PTE 标记为有效
\lineref{kernel/vm.c:/perm...PTE_V/}。

\indexcode{walk}
\lineref{kernel/vm.c:/^walk/}
模仿 RISC-V 分页硬件查找虚拟地址的 PTE（参见
Figure~\ref{fig:riscv_pagetable}）。
\lstinline{walk}
逐级下降页表树，使用每级的 9 位虚拟地址索引到相关的页目录页。在每一级，它找到下一级页目录页的 PTE 或最终页的 PTE
\lineref{kernel/vm.c:/pte.=..pagetable/}。
如果第一级或第二级页目录页中的 PTE 无效，则所需的目录页尚未分配；如果
\lstinline{alloc}
参数已设置，
\lstinline{walk}
分配一个新的页表页并将其物理地址放入 PTE。它返回树中最底层 PTE 的地址
\lineref{kernel/vm.c:/return..pagetable/}。

上述代码依赖于物理内存直接映射到内核虚拟地址空间。例如，当 \lstinline{walk} 逐级下降页表时，它从 PTE 中提取下一级页表的（物理）地址 \lineref{kernel/vm.c:/pagetable.=..pa.*E2P/}，然后使用该地址作为虚拟地址来获取下一级 PTE
\lineref{kernel/vm.c:/t..pte.=..paget/}。

在每个 CPU 上，
\indexcode{main}
调用
\indexcode{kvminithart}
\lineref{kernel/vm.c:/^kvminithart/}
来安装内核页表，将根页表页的物理地址放入 CPU 的
\texttt{satp}
寄存器。此后，CPU 使用内核页表翻译地址。内核继续正确执行，因为内核页表是直接映射的，因此地址在此更改前后引用 RAM 中的相同位置。

每个 RISC-V CPU 在
\indextext{Translation Look-aside Buffer (TLB)}（转换后备缓冲器）中缓存页表项，当 xv6 更改页表时，它必须告诉 CPU 使相应的缓存 TLB 条目无效。如果没有这样做，那么在之后的某个时刻，TLB 可能会使用旧的缓存映射，指向在此期间已分配给另一个进程的物理页，结果进程可能会在其他进程的内存上涂鸦。RISC-V 有一条指令 \indexcode{sfence.vma} 刷新当前 CPU 的 TLB。Xv6 在重新加载
\texttt{satp}
寄存器后在 {\tt kvminithart} 中执行 {\tt sfence.vma}，以及在
\lstinline{uservec}
和 \lstinline{userret} 的 trampoline 代码中执行。

在更改
\texttt{satp}
之前，还需要发出 \texttt{sfence.vma}，以等待所有未完成的加载和存储完成。此等待确保对页表的先前更新已完成，并确保先前的加载和存储使用旧页表，而不是新页表。

\section{物理内存分配}

内核必须在运行时分配和释放物理内存，用于页表、用户内存、内核栈和管道缓冲区。

Xv6 使用内核末尾和 \indexcode{PHYSTOP} 之间的物理内存进行运行时分配。它每次分配和释放整个 4096 字节的页。它通过在页面本身中内嵌链表来跟踪哪些页面是空闲的。分配即把一页从链表中删除；释放即将页面添加到链表中。

请阅读 {\tt kernel/kalloc.c}。

%%
\section{代码：物理内存分配器}
%%

分配器位于 {\tt kalloc.c} \fileref{kernel/kalloc.c}。其数据结构是一个 \textit{空闲列表}（free list），包含可用于分配的物理内存页。每个空闲页存储其``下一个''指针的位置在 \indexcode{struct run} \lineref{kernel/kalloc.c:/^struct.run/}。分配器将每个空闲页的 \lstinline{run} 结构存储在空闲页本身中，因为页面空闲时无其他内容存储。空闲列表受自旋锁保护 \linerefs{kernel/kalloc.c:/^struct.{/,/}/}。列表和锁包装在一个结构体中，以明确锁保护结构体中的字段。现在，忽略锁和对 \lstinline{acquire} 与 \lstinline{release} 的调用；Chapter~\ref{CH:LOCK} 将详细讨论锁定。

函数
\indexcode{main}
调用
\indexcode{kinit}
来初始化分配器
\lineref{kernel/kalloc.c:/^kinit/}。
\lstinline{kinit} 将空闲列表初始化为包含内核末尾和 {\tt PHYSTOP} 之间的每个物理 RAM 页。Xv6 本应该通过解析硬件提供的配置信息来确定有多少物理内存可用。但实际上，xv6 假设机器有 128 MB 的 RAM。
\lstinline{kinit}
调用
\indexcode{freerange}
通过每页调用
\indexcode{kfree}
来将内存添加到空闲列表。PTE 只能引用对齐到 4096 字节边界（是 4096 的倍数）的物理地址，因此
\lstinline{freerange}
使用
\indexcode{PGROUNDUP}
来确保它仅释放对齐的物理地址。分配器从没有内存开始；这些对
\lstinline{kfree}
的调用给它一些内存来管理。

分配器有时将地址视为整数以便对它们执行算术运算（例如，在 \lstinline{freerange} 中遍历所有页），有时使用地址作为指针来读写内存（例如，操作存储在每页中的 \lstinline{run} 结构）；地址的这种双重用途是分配器代码充满 C 类型转换的主要原因。
\index{type cast}

函数
\lstinline{kfree}
\lineref{kernel/kalloc.c:/^kfree/}
首先将正在释放的内存中的每个字节设置为值 1。这将导致在释放内存后使用内存的代码（使用``悬空引用''）读取垃圾而不是旧的 valid 内容；希望这将导致此类代码更快地崩溃。然后
\lstinline{kfree}
将页面前置到空闲列表：它将
\lstinline{pa}
转换为指向
\lstinline{struct}
\lstinline{run}
的指针，将空闲列表的旧开始记录在
\lstinline{r->next}，
并将空闲列表设置为
\lstinline{r}。
\indexcode{kalloc}
从空闲列表中移除并返回第一个元素。

\section{进程地址空间}

每个进程都有自己的页表，当 xv6 在进程之间切换时，它也会更改页表。
Figure~\ref{fig:processlayout} 比 Figure~\ref{fig:as} 更详细地显示了进程的地址空间。进程的用户地址空间从零开始，理论上结束于
\texttt{MAXVA}
({\tt 0x4000000000})
\lineref{kernel/riscv.h:/MAXVA/}，尽管实际上只有其中一小部分映射到物理内存。

进程的地址空间由以下页面组成：包含程序文本的页面（xv6 使用权限 \lstinline{PTE_R}、\lstinline{PTE_X} 和 \lstinline{PTE_U} 映射）、包含程序预初始化数据的页面、栈页面以及堆页面。Xv6 使用权限 \lstinline{PTE_R}、\lstinline{PTE_W} 和 \lstinline{PTE_U} 映射数据、栈和堆。

在用户地址空间内使用权限是强化用户进程的常用技术。如果文本使用 \lstinline{PTE_W} 映射，那么进程可能会意外修改自己的程序；例如，编程错误可能导致程序写入空指针，修改地址 0 处的指令，然后继续运行，可能造成更大的破坏。为了立即检测此类错误，xv6 在不使用 \lstinline{PTE_W} 的情况下映射文本；如果程序意外尝试存储到地址 0，硬件将拒绝执行存储并引发页错误（参见 Chapter~\ref{CH:TRAP}）。然后内核杀死进程并打印出信息性消息，以便开发人员可以追踪问题。

同样，通过在不使用 \lstinline{PTE_X} 的情况下映射数据，用户程序不能意外地跳转到程序数据中的地址并从该地址开始执行。

在现实世界中，通过仔细设置权限来强化进程也有助于防御安全攻击。攻击者可能向程序（例如 Web 服务器）提供精心构造的输入，触发程序中的错误，希望将该错误转化为漏洞利用~\cite{aleph:smashing}。仔细设置权限和其他技术（如用户地址空间布局随机化）使此类攻击更难。

栈是单页，显示为 {\tt exec} 系统调用创建的初始内容。包含命令行参数的字符串以及指向它们的指针数组位于栈的最顶部。就在那下面是允许程序在
\lstinline{main}
开始执行，就好像函数 \lstinline{main(argc}、\lstinline{argv)} 刚刚被调用一样。

为了检测用户栈溢出分配的栈内存，xv6 通过清除 \lstinline{PTE_U} 标志在栈正下方放置一个不可访问的保护页。如果用户栈溢出并且进程尝试使用栈下方的地址，硬件将生成页错误异常，因为保护页对用户模式下运行的程序不可访问。现实世界的操作系统可能会在栈溢出时自动为用户栈分配更多内存。

我们在这里看到了一些页表使用的很好的例子。首先，不同进程的页表将用户地址转换为不同的物理内存页，因此每个进程都有私有的用户内存。其次，每个进程将其内存视为具有从零开始的连续虚拟地址，而进程的物理内存可以是非连续的。第三，内核在用户地址空间顶部映射一个具有 trampoline 代码的页（没有 \lstinline{PTE_U}），因此单个物理内存页出现在所有地址空间中，但只能由内核使用。

\begin{figure}[t]
\centering
\includegraphics[scale=0.5]{fig/processlayout.png}
\caption{进程的用户地址空间及其初始栈。}
\label{fig:processlayout}
\end{figure}


%%
\section{代码：exec}
%%

请阅读 {\tt kernel/exec.c} 和 {\tt kernel/vm.c}，从 {\tt uvmcreate()} 开始。

\lstinline{exec} 是一个系统调用，用从文件读取的数据（称为二进制或可执行文件）替换进程的用户地址空间。二进制通常是编译器和链接器的输出，保存机器指令和程序数据。
\lstinline{kexec} \lineref{kernel/exec.c:/^kexec/}，
{\tt exec} 的内核内部实现，使用 \indexcode{namei}
\lineref{kernel/exec.c:/namei/}
打开命名的二进制文件 \lstinline{path}，这在 Chapter~\ref{CH:FS} 中解释。然后，它读取 ELF 头。Xv6 二进制文件以广泛使用的 \indextext{ELF format}（ELF 格式）格式化，定义在
\fileref{kernel/elf.h}。ELF 二进制文件由 ELF 头
\indexcode{struct elfhdr} \lineref{kernel/elf.h:/^struct.elfhdr/}、
后跟一系列程序段头
\lstinline{struct proghdr} \lineref{kernel/elf.h:/^struct.proghdr/} 组成。每个
\lstinline{proghdr}
描述必须加载到内存中的应用段；xv6 程序有两个程序段头：一个用于指令，一个用于数据。

第一步是快速检查文件是否可能包含 ELF 二进制文件。ELF 二进制文件以四字节``魔数''
\lstinline{0x7F}、
\lstinline{`E'}、
\lstinline{`L'}、
\lstinline{`F'}
或
\indexcode{ELF_MAGIC}
\lineref{kernel/elf.h:/ELF_MAGIC/} 开头。
如果 ELF 头具有正确的魔数，
\lstinline{kexec}
假设二进制文件格式正确。

\lstinline{kexec}
使用
\indexcode{proc\_pagetable}
\lineref{kernel/exec.c:/proc\_pagetable/}
分配一个没有用户映射的新页表，使用
\indexcode{uvmalloc}
\lineref{kernel/exec.c:/uvmalloc/}
为每个 ELF 段分配内存，并使用
\indexcode{loadseg}
\lineref{kernel/exec.c:/loadseg/}
将每个段加载到内存中。
\lstinline{loadseg}
使用
\indexcode{walkaddr}
来查找分配的内存的物理地址，以便写入
ELF 段的每一页，并使用
\indexcode{readi}
从文件中读取。

\indexcode{/init}
的程序段头，
用
\lstinline{exec}
创建的第一个用户程序，
看起来像这样：
\begin{footnotesize}
\begin{verbatim}
# objdump -p user/_init

user/_init:     file format elf64-little

Program Header:
0x70000003 off    0x0000000000006bb0 vaddr 0x0000000000000000
                                       paddr 0x0000000000000000 align 2**0
         filesz 0x000000000000004a memsz 0x0000000000000000 flags r--
    LOAD off    0x0000000000001000 vaddr 0x0000000000000000
                                       paddr 0x0000000000000000 align 2**12
         filesz 0x0000000000001000 memsz 0x0000000000001000 flags r-x
    LOAD off    0x0000000000002000 vaddr 0x0000000000001000
                                       paddr 0x0000000000001000 align 2**12
         filesz 0x0000000000000010 memsz 0x0000000000000030 flags rw-
   STACK off    0x0000000000000000 vaddr 0x0000000000000000
                                       paddr 0x0000000000000000 align 2**4
         filesz 0x0000000000000000 memsz 0x0000000000000000 flags rw-
\end{verbatim}
\end{footnotesize}

我们看到文本段应该加载到内存中的虚拟地址 0x0（无写权限），从文件中偏移量 0x1000 处的内容加载。我们还看到数据应该加载到地址 0x1000，该地址位于页边界，并且没有可执行权限。

程序段头的
\lstinline{filesz}
可能小于
\lstinline{memsz}，
表示它们之间的间隙应该用零填充（对于 C 全局变量）而不是从文件中读取。
对于
\lstinline{/init}，数据
\lstinline{filesz}
为 0x10 字节，
\lstinline{memsz}
为 0x30 字节，
因此
\indexcode{uvmalloc}
分配足够的物理内存以容纳 0x30 字节，但仅从文件读取 0x10 字节
\lstinline{/init}。

现在
\indexcode{kexec}
分配并初始化用户栈。它只分配一个栈页。
\lstinline{kexec}
将参数字符串一次一个地复制到栈顶，在
\indexcode{ustack}
中记录指向它们的指针。它在
\indexcode{argv}
列表的末尾放置一个空指针，该列表将传递给
\lstinline{main}。
\indexcode{argc}
和 \lstinline{argv} 的值通过系统调用返回路径传递给 \lstinline{main}：\lstinline{argc} 通过系统调用返回值传递，放在 {\tt a0} 中，\lstinline{argv} 通过进程 trapframe 的 {\tt a1} 条目传递。

\lstinline{kexec}
在栈页正下方放置一个不可访问的页，因此尝试使用多个页的程序将出错。此不可访问的页还允许
\lstinline{kexec}
处理太大的参数；在这种情况下，
\lstinline{kexec}
使用的
\indexcode{copyout}
\lineref{kernel/vm.c:/^copyout/}
函数将注意到目标页不可访问，并将返回 -1。

在准备新内存映像期间，如果
\lstinline{kexec}
检测到错误（如无效程序段），它会跳转到标签
\lstinline{bad}，
释放新映像，并返回 -1。
\lstinline{kexec}
必须等到确定系统调用成功后才能释放旧映像：如果旧映像已消失，系统调用无法向它返回 -1。
\lstinline{kexec}
中的唯一错误情况发生在映像创建期间。一旦映像完成，
\lstinline{kexec}
就可以提交到新页表
\lineref{kernel/exec.c:/pagetable.=.pagetable/}
并释放旧页表
\lineref{kernel/exec.c:/proc_freepagetable/}。

\lstinline{exec} 系统调用将字节从 ELF 文件加载到 ELF 文件指定的内存地址。用户或进程可以在 ELF 文件中放置任意地址。因此 \lstinline{exec} 是有风险的，因为 ELF 文件中的地址可能意外或故意引用内核。对粗心的内核而言，后果可能从崩溃到恶意颠覆内核的隔离机制（即安全漏洞）。Xv6 执行许多检查以避免这些风险。例如，\lstinline{if(ph.vaddr + ph.memsz < ph.vaddr)} 检查总和是否溢出 64 位整数。危险在于用户可以构造一个 ELF 二进制文件，其 \lstinline{ph.vaddr} 指向用户选择的地址，且 \lstinline{ph.memsz} 足够大，使总和溢出到 0x1000，这看起来像一个有效值。在 xv6 的旧版本中，用户地址空间也包含内核（但在用户模式下不可读/写），用户可以选择对应于内核内存的地址，从而将数据从 ELF 二进制文件复制到内核中。在 RISC-V 版本的 xv6 中，这不可能发生，因为内核有自己的独立页表；
\lstinline{loadseg}
加载到进程的页表中，而不是内核的页表中。

内核开发人员很容易省略关键检查，现实世界的内核长期缺少可以被用户程序利用来获取内核权限的检查。xv6 可能未完全验证提供给内核的用户级数据，恶意用户程序可能利用这一点来绕过 xv6 的隔离。
%%
\section{现实世界}
%%

与大多数操作系统一样，xv6 使用分页硬件进行内存保护和映射。大多数操作系统通过结合分页和页错误异常更复杂地使用分页，我们将在 Chapter~\ref{CH:TRAP} 中讨论。

Xv6 因内核使用虚拟地址和物理地址之间的直接映射，以及假设物理 RAM 位于地址 0x80000000（内核期望加载的位置）而简化。这适用于 QEMU，但在真实硬件上效果不佳；真实硬件将 RAM 和设备放置在不可预测的物理地址，因此（例如）0x80000000 处可能没有 RAM，而 xv6 期望能够在该处存储内核。更严肃的内核设计利用页表将任意的硬件物理内存布局转换为可预测的内核虚拟地址布局。

RISC-V 支持物理地址级别的保护，但 xv6 不使用该功能。

在内存较多的机器上，使用 RISC-V 对``超级页''的支持可能是有意义的。小页面在物理内存较小时有意义，以允许细粒度分配和换出到磁盘。例如，如果程序仅使用 8 KB 内存，为其分配整个 4 MB 的超级物理内存页是浪费的。较大的页面在 RAM 较多的机器上有意义，并可能减少页表操作的开销。

为了避免在更改页表时必须刷新整个 TLB，RISC-V CPU 可能支持地址空间标识符（ASIDs）~\cite{riscv:priv}。然后内核可以只刷新特定地址空间的 TLB 条目。Xv6 不使用此功能。

xv6 内核缺乏类似 {\tt malloc} 的分配器来为小对象提供内存，这阻止了内核使用需要动态分配的复杂数据结构。更复杂的内核可能会分配许多不同大小的小块，而不是（如 xv6）仅 4096 字节块；真正的内核分配器需要处理小分配和大分配。

内存分配是一个永恒的热门话题，基本问题是有效使用有限内存和为未知的未来请求做准备~\cite{knuth}。今天人们更关心速度而不是空间效率。
%%
\section{练习}
%%

\begin{enumerate}

\item 解析 RISC-V 的设备树以查找计算机拥有的物理内存量。

\item 函数 {\tt copyin} 和 {\tt copyinstr} 在软件中遍历用户页表。设置内核页表，使内核映射了用户程序，并且 {\tt copyin} 和 {\tt copyinstr} 可以使用 {\tt memcpy} 将系统调用参数复制到内核空间，依靠硬件进行页表遍历。

\item 修改 xv6 以在内核中使用超级页。

\item Unix 的
\lstinline{exec}
实现对 shell 脚本有特殊处理。如果要执行的文件以文本
\lstinline{#!}
开头，则第一行被视为运行来解释文件的程序。例如，如果
\lstinline{exec}
被调用以运行
\lstinline{myprog}
\lstinline{arg1}
并且
\lstinline{myprog}
的第一行是
\lstinline{#!/interp}，
则
\lstinline{exec}
运行
\lstinline{/interp}
，命令行为
\lstinline{/interp}
\lstinline{myprog}
\lstinline{arg1}。
在 xv6 中实现对此约定的支持。

\item 为内核实现地址空间布局随机化。

\end{enumerate}

\chapter{陷入和系统调用}
\label{CH:TRAP}

有三种事件会导致 CPU 暂停指令的正常执行，并强制将控制权转移到处理该事件的特殊内核代码。一种情况是系统调用，当用户程序执行 {\tt ecall} 指令请求内核为它做某事时。另一种情况是 \indextext{exception}（异常）：指令（用户或内核）做了非法的事情，例如从无效的虚拟地址加载。第三种情况是设备 \indextext{interrupt}（中断），当设备发出需要关注的信号时，例如磁盘硬件完成读取或写入请求时。

本书使用 \indextext{trap}（陷入）作为这些情况的统称。通常，在陷入时执行的任何代码稍后都需要恢复，且不应意识到发生了什么特殊的事情。也就是说，我们希望陷入是透明的；这对于设备中断尤其重要，因为被中断的代码通常不期望中断的发生。陷入的处理流程如下：强制将控制权转移到内核；内核保存寄存器和其他状态以便恢复执行；内核执行适当的处理程序代码（例如，系统调用实现或设备驱动）；内核恢复保存的状态并从陷入返回；原始代码从停止的地方继续执行。

Xv6 在内核中处理所有陷入；陷入不会传递给用户代码。在内核中处理系统调用是顺理成章的。对于中断而言，这也是合理的，因为隔离机制要求只有内核才能使用设备，而且内核能够在多个进程之间共享设备。对于异常而言这同样合理，因为内核可能能够处理来自用户空间的异常（例如参见 Chapter~\ref{CH:PGFAULTS}）或通过终止违规程序来响应。

Xv6 陷入处理分为四个阶段：RISC-V CPU 采取的硬件操作、为内核 C 代码铺平道路的一些汇编指令、决定如何处理陷入的 C 函数，以及系统调用或设备驱动服务例程。虽然三种陷入类型之间的共性表明内核可以用单一代码路径处理所有陷入，但事实证明为两种不同情况设置单独的代码是方便的：来自用户空间的陷入和来自内核空间的陷入。处理陷入的内核代码（汇编或 C）通常称为 \indextext{handler}（处理程序）；处理程序的第一条指令通常用汇编（而非 C）编写，有时称为 \indextext{vector}（向量入口）。

在继续之前，请阅读 {\tt kernel/trampoline.S} 以及 {\tt kernel/trap.c} 中的 {\tt usertrap()} 和 {\tt prepare\_return()}。

\section{RISC-V 陷入机制}

每个 RISC-V CPU 都有一组硬件控制寄存器，内核写入这些寄存器以告诉 CPU 如何处理陷入，并且内核可以读取这些寄存器以了解已发生的陷入。RISC-V 文档中有完整说明~\cite{riscv:priv}。{\tt riscv.h}
\fileref{kernel/riscv.h}
包含 xv6 使用的定义。以下是最重要的寄存器的概述：

\begin{itemize}

\item \indexcode{stvec}：内核在此处写入其陷入处理程序代码的地址；RISC-V 跳转到 {\tt stvec} 中的地址以处理陷入。

\item \indexcode{sepc}：当发生陷入时，RISC-V 在此处保存程序计数器（因为 {\tt pc} 随后被 {\tt stvec} 中的值覆盖）。{\tt sret}（从陷入返回）指令将 {\tt sepc} 复制到 {\tt pc}。内核可以写入 {\tt sepc} 以控制 {\tt sret} 的去向。

\item \indexcode{scause}：RISC-V 在此处放置一个数字，描述陷入的原因。

\item \indexcode{sscratch}：内核陷入处理程序代码使用 {\tt sscratch} 来帮助它在保存用户寄存器之前避免覆盖它们。

\item \indexcode{sstatus}：{\tt sstatus} 中的 SIE（Supervisor Interrupt Enable，监督模式中断使能）位控制是否启用设备中断。如果内核清除 SIE，RISC-V 将推迟设备中断，直到内核重新设置 SIE。SPP（Supervisor Previous Privilege，监督模式先前特权）位指示陷入来自用户模式还是监督模式，并控制 {\tt sret} 返回到哪种模式。

\end{itemize}

上述寄存器只能在监管者模式下（即内核）访问；CPU 阻止用户代码读取或写入它们。

多核芯片上的每个 CPU 都有自己的一组这些寄存器，并且可能有多个 CPU 在任何给定时间处理陷入。

当 RISC-V 硬件强制陷入时，会执行以下操作：

\begin{enumerate}

\item 如果陷入是设备中断，且 {\tt sstatus} 的 SIE 位被清除，则不执行以下任何操作。

\item 通过清除 {\tt sstatus} 中的 SIE 位来禁用中断。

\item 将 {\tt pc} 复制到 {\tt sepc}。

\item 将当前模式（用户或监管者）保存在 {\tt sstatus} 的 SPP 位中。

\item 将 {\tt scause} 设置为一个表示陷入原因的数字。

\item 将模式设置为监管者模式。

\item 将 {\tt stvec} 复制到 {\tt pc}。

\item 开始在新的 {\tt pc} 处执行。

\end{enumerate}

CPU 不会切换到内核页表，不会切换到内核栈，也不会保存除 {\tt pc} 之外的任何寄存器。这些任务必须由内核软件完成。CPU 在陷入期间只做最少工作的原因之一是为软件提供灵活性；例如，某些操作系统在某些情况下会省略页表切换以提高陷入处理的性能。

值得思考的是，上面列出的任何步骤是否可以省略，也许是为了更快的陷入。虽然在某些情况下更简单的序列可以工作，但许多步骤在一般情况下省略是危险的。例如，假设 CPU 没有切换程序计数器。那么来自用户空间的陷入可能会在仍处于用户指令时切换到监管者模式。这些用户指令可能会破坏用户/内核隔离，例如通过修改 {\tt satp} 寄存器以指向允许访问所有物理内存的页表。因此，CPU 切换到内核指定的指令地址，即 {\tt stvec}，这一点很重要。

\section{来自用户空间的陷入}
\begin{figure}[t]
  \centering
    \chapter{陷入和系统调用}
\label{CH:TRAP}

有三种事件会导致 CPU 暂停指令的正常执行，并强制将控制权转移到处理该事件的特殊内核代码。一种情况是系统调用，当用户程序执行 {\tt ecall} 指令请求内核为它做某事时。另一种情况是 \indextext{exception}（异常）：指令（用户或内核）做了非法的事情，例如从无效的虚拟地址加载。第三种情况是设备 \indextext{interrupt}（中断），当设备发出需要关注的信号时，例如磁盘硬件完成读取或写入请求时。

本书使用 \indextext{trap}（陷入）作为这些情况的统称。通常，在陷入时执行的任何代码稍后都需要恢复，且不应意识到发生了什么特殊的事情。也就是说，我们希望陷入是透明的；这对于设备中断尤其重要，因为被中断的代码通常不期望中断的发生。陷入的处理流程如下：强制将控制权转移到内核；内核保存寄存器和其他状态以便恢复执行；内核执行适当的处理程序代码（例如，系统调用实现或设备驱动）；内核恢复保存的状态并从陷入返回；原始代码从停止的地方继续执行。

Xv6 在内核中处理所有陷入；陷入不会传递给用户代码。在内核中处理系统调用是顺理成章的。对于中断而言，这也是合理的，因为隔离机制要求只有内核才能使用设备，而且内核能够在多个进程之间共享设备。对于异常而言这同样合理，因为内核可能能够处理来自用户空间的异常（例如参见 Chapter~\ref{CH:PGFAULTS}）或通过终止违规程序来响应。

Xv6 陷入处理分为四个阶段：RISC-V CPU 采取的硬件操作、为内核 C 代码铺平道路的一些汇编指令、决定如何处理陷入的 C 函数，以及系统调用或设备驱动服务例程。虽然三种陷入类型之间的共性表明内核可以用单一代码路径处理所有陷入，但事实证明为两种不同情况设置单独的代码是方便的：来自用户空间的陷入和来自内核空间的陷入。处理陷入的内核代码（汇编或 C）通常称为 \indextext{handler}（处理程序）；处理程序的第一条指令通常用汇编（而非 C）编写，有时称为 \indextext{vector}（向量入口）。

在继续之前，请阅读 {\tt kernel/trampoline.S} 以及 {\tt kernel/trap.c} 中的 {\tt usertrap()} 和 {\tt prepare\_return()}。

\section{RISC-V 陷入机制}

每个 RISC-V CPU 都有一组硬件控制寄存器，内核写入这些寄存器以告诉 CPU 如何处理陷入，并且内核可以读取这些寄存器以了解已发生的陷入。RISC-V 文档中有完整说明~\cite{riscv:priv}。{\tt riscv.h}
\fileref{kernel/riscv.h}
包含 xv6 使用的定义。以下是最重要的寄存器的概述：

\begin{itemize}

\item \indexcode{stvec}：内核在此处写入其陷入处理程序代码的地址；RISC-V 跳转到 {\tt stvec} 中的地址以处理陷入。

\item \indexcode{sepc}：当发生陷入时，RISC-V 在此处保存程序计数器（因为 {\tt pc} 随后被 {\tt stvec} 中的值覆盖）。{\tt sret}（从陷入返回）指令将 {\tt sepc} 复制到 {\tt pc}。内核可以写入 {\tt sepc} 以控制 {\tt sret} 的去向。

\item \indexcode{scause}：RISC-V 在此处放置一个数字，描述陷入的原因。

\item \indexcode{sscratch}：内核陷入处理程序代码使用 {\tt sscratch} 来帮助它在保存用户寄存器之前避免覆盖它们。

\item \indexcode{sstatus}：{\tt sstatus} 中的 SIE（Supervisor Interrupt Enable，监督模式中断使能）位控制是否启用设备中断。如果内核清除 SIE，RISC-V 将推迟设备中断，直到内核重新设置 SIE。SPP（Supervisor Previous Privilege，监督模式先前特权）位指示陷入来自用户模式还是监督模式，并控制 {\tt sret} 返回到哪种模式。

\end{itemize}

上述寄存器只能在监管者模式下（即内核）访问；CPU 阻止用户代码读取或写入它们。

多核芯片上的每个 CPU 都有自己的一组这些寄存器，并且可能有多个 CPU 在任何给定时间处理陷入。

当 RISC-V 硬件强制陷入时，会执行以下操作：

\begin{enumerate}

\item 如果陷入是设备中断，且 {\tt sstatus} 的 SIE 位被清除，则不执行以下任何操作。

\item 通过清除 {\tt sstatus} 中的 SIE 位来禁用中断。

\item 将 {\tt pc} 复制到 {\tt sepc}。

\item 将当前模式（用户或监管者）保存在 {\tt sstatus} 的 SPP 位中。

\item 将 {\tt scause} 设置为一个表示陷入原因的数字。

\item 将模式设置为监管者模式。

\item 将 {\tt stvec} 复制到 {\tt pc}。

\item 开始在新的 {\tt pc} 处执行。

\end{enumerate}

CPU 不会切换到内核页表，不会切换到内核栈，也不会保存除 {\tt pc} 之外的任何寄存器。这些任务必须由内核软件完成。CPU 在陷入期间只做最少工作的原因之一是为软件提供灵活性；例如，某些操作系统在某些情况下会省略页表切换以提高陷入处理的性能。

值得思考的是，上面列出的任何步骤是否可以省略，也许是为了更快的陷入。虽然在某些情况下更简单的序列可以工作，但许多步骤在一般情况下省略是危险的。例如，假设 CPU 没有切换程序计数器。那么来自用户空间的陷入可能会在仍处于用户指令时切换到监管者模式。这些用户指令可能会破坏用户/内核隔离，例如通过修改 {\tt satp} 寄存器以指向允许访问所有物理内存的页表。因此，CPU 切换到内核指定的指令地址，即 {\tt stvec}，这一点很重要。

\section{来自用户空间的陷入}
\begin{figure}[t]
  \centering
    \chapter{陷入和系统调用}
\label{CH:TRAP}

有三种事件会导致 CPU 暂停指令的正常执行，并强制将控制权转移到处理该事件的特殊内核代码。一种情况是系统调用，当用户程序执行 {\tt ecall} 指令请求内核为它做某事时。另一种情况是 \indextext{exception}（异常）：指令（用户或内核）做了非法的事情，例如从无效的虚拟地址加载。第三种情况是设备 \indextext{interrupt}（中断），当设备发出需要关注的信号时，例如磁盘硬件完成读取或写入请求时。

本书使用 \indextext{trap}（陷入）作为这些情况的统称。通常，在陷入时执行的任何代码稍后都需要恢复，且不应意识到发生了什么特殊的事情。也就是说，我们希望陷入是透明的；这对于设备中断尤其重要，因为被中断的代码通常不期望中断的发生。陷入的处理流程如下：强制将控制权转移到内核；内核保存寄存器和其他状态以便恢复执行；内核执行适当的处理程序代码（例如，系统调用实现或设备驱动）；内核恢复保存的状态并从陷入返回；原始代码从停止的地方继续执行。

Xv6 在内核中处理所有陷入；陷入不会传递给用户代码。在内核中处理系统调用是顺理成章的。对于中断而言，这也是合理的，因为隔离机制要求只有内核才能使用设备，而且内核能够在多个进程之间共享设备。对于异常而言这同样合理，因为内核可能能够处理来自用户空间的异常（例如参见 Chapter~\ref{CH:PGFAULTS}）或通过终止违规程序来响应。

Xv6 陷入处理分为四个阶段：RISC-V CPU 采取的硬件操作、为内核 C 代码铺平道路的一些汇编指令、决定如何处理陷入的 C 函数，以及系统调用或设备驱动服务例程。虽然三种陷入类型之间的共性表明内核可以用单一代码路径处理所有陷入，但事实证明为两种不同情况设置单独的代码是方便的：来自用户空间的陷入和来自内核空间的陷入。处理陷入的内核代码（汇编或 C）通常称为 \indextext{handler}（处理程序）；处理程序的第一条指令通常用汇编（而非 C）编写，有时称为 \indextext{vector}（向量入口）。

在继续之前，请阅读 {\tt kernel/trampoline.S} 以及 {\tt kernel/trap.c} 中的 {\tt usertrap()} 和 {\tt prepare\_return()}。

\section{RISC-V 陷入机制}

每个 RISC-V CPU 都有一组硬件控制寄存器，内核写入这些寄存器以告诉 CPU 如何处理陷入，并且内核可以读取这些寄存器以了解已发生的陷入。RISC-V 文档中有完整说明~\cite{riscv:priv}。{\tt riscv.h}
\fileref{kernel/riscv.h}
包含 xv6 使用的定义。以下是最重要的寄存器的概述：

\begin{itemize}

\item \indexcode{stvec}：内核在此处写入其陷入处理程序代码的地址；RISC-V 跳转到 {\tt stvec} 中的地址以处理陷入。

\item \indexcode{sepc}：当发生陷入时，RISC-V 在此处保存程序计数器（因为 {\tt pc} 随后被 {\tt stvec} 中的值覆盖）。{\tt sret}（从陷入返回）指令将 {\tt sepc} 复制到 {\tt pc}。内核可以写入 {\tt sepc} 以控制 {\tt sret} 的去向。

\item \indexcode{scause}：RISC-V 在此处放置一个数字，描述陷入的原因。

\item \indexcode{sscratch}：内核陷入处理程序代码使用 {\tt sscratch} 来帮助它在保存用户寄存器之前避免覆盖它们。

\item \indexcode{sstatus}：{\tt sstatus} 中的 SIE（Supervisor Interrupt Enable，监督模式中断使能）位控制是否启用设备中断。如果内核清除 SIE，RISC-V 将推迟设备中断，直到内核重新设置 SIE。SPP（Supervisor Previous Privilege，监督模式先前特权）位指示陷入来自用户模式还是监督模式，并控制 {\tt sret} 返回到哪种模式。

\end{itemize}

上述寄存器只能在监管者模式下（即内核）访问；CPU 阻止用户代码读取或写入它们。

多核芯片上的每个 CPU 都有自己的一组这些寄存器，并且可能有多个 CPU 在任何给定时间处理陷入。

当 RISC-V 硬件强制陷入时，会执行以下操作：

\begin{enumerate}

\item 如果陷入是设备中断，且 {\tt sstatus} 的 SIE 位被清除，则不执行以下任何操作。

\item 通过清除 {\tt sstatus} 中的 SIE 位来禁用中断。

\item 将 {\tt pc} 复制到 {\tt sepc}。

\item 将当前模式（用户或监管者）保存在 {\tt sstatus} 的 SPP 位中。

\item 将 {\tt scause} 设置为一个表示陷入原因的数字。

\item 将模式设置为监管者模式。

\item 将 {\tt stvec} 复制到 {\tt pc}。

\item 开始在新的 {\tt pc} 处执行。

\end{enumerate}

CPU 不会切换到内核页表，不会切换到内核栈，也不会保存除 {\tt pc} 之外的任何寄存器。这些任务必须由内核软件完成。CPU 在陷入期间只做最少工作的原因之一是为软件提供灵活性；例如，某些操作系统在某些情况下会省略页表切换以提高陷入处理的性能。

值得思考的是，上面列出的任何步骤是否可以省略，也许是为了更快的陷入。虽然在某些情况下更简单的序列可以工作，但许多步骤在一般情况下省略是危险的。例如，假设 CPU 没有切换程序计数器。那么来自用户空间的陷入可能会在仍处于用户指令时切换到监管者模式。这些用户指令可能会破坏用户/内核隔离，例如通过修改 {\tt satp} 寄存器以指向允许访问所有物理内存的页表。因此，CPU 切换到内核指定的指令地址，即 {\tt stvec}，这一点很重要。

\section{来自用户空间的陷入}
\begin{figure}[t]
  \centering
    \input{fig/trap.tex}
\caption{处理来自用户代码的陷入的概要。}
\label{fig:usertrap}
\end{figure}

Xv6 根据陷入发生在内核执行还是用户代码执行中而有所不同地处理陷入。以下是用户代码陷入的处理流程；Section~\ref{s:ktraps} 将介绍内核代码的陷入处理。

如果用户程序进行系统调用（{\tt ecall} 指令），或执行了非法操作，或发生设备中断，则可能在用户空间执行时发生陷入。如 Figure~\ref{fig:usertrap} 所示，来自用户空间的高级别陷入路径是
{\tt uservec}
\lineref{kernel/trampoline.S:/^uservec/}，
然后是 {\tt usertrap}
\lineref{kernel/trap.c:/^usertrap/}；
当内核准备返回时，
{\tt usertrap}
返回给
{\tt userret}
\lineref{kernel/trampoline.S:/^userret/}，
它执行 {\tt sret}
返回到用户空间。

% talk about why RISC-V doesn't switch page tables

xv6 陷入处理设计的一个主要限制是 RISC-V 硬件在强制陷入时不会切换页表。这意味着 {\tt stvec} 中的陷入处理程序地址必须在用户页表中有有效映射，因为这是在陷入处理代码开始执行时生效的页表。此外，xv6 的陷入处理代码需要切换到内核页表；为了能够在切换后继续执行，内核页表也必须具有 {\tt stvec} 指向的处理程序的映射。

Xv6 使用 \indextext{trampoline}（跳板）页来满足这些要求。此页包含 {\tt uservec}，即 {\tt stvec} 指向的 xv6 陷入处理代码。trampoline 页在每个进程的页表中映射到虚拟地址 {\tt 0x3ffffff000}（称为 \indexcode{TRAMPOLINE}），这是虚拟地址空间中的最后一页，以便它位于程序自己使用的内存之上。trampoline 页在内核页表中映射到相同的虚拟地址。参见 Figure~\ref{fig:as} 和 Figure~\ref{fig:xv6_layout}。因为 trampoline 页映射在用户页表中，所以陷入可以在监管者模式下开始在那里执行。因为 trampoline 页在内核地址空间中映射到相同的地址，所以陷入处理程序可以在切换到内核页表后继续执行。

{\tt uservec} 陷阱处理程序的代码位于 {\tt trampoline.S}
\lineref{kernel/trampoline.S:/^uservec/}。
当 {\tt uservec} 开始时，所有 32 个寄存器都包含被中断用户代码的值。这 32 个值需要保存在内存中的某个位置，以便内核在稍后返回用户空间之前恢复它们。存储到内存需要使用寄存器来保存目标地址，但此时没有通用寄存器可用！幸运的是，RISC-V 以 {\tt sscratch} 寄存器的形式提供了帮助。{\tt uservec} 开头的 {\tt csrw} 指令将 {\tt a0} 保存在 {\tt sscratch} 中。现在 {\tt uservec} 有一个寄存器（{\tt a0}）可以使用了。

{\tt uservec} 的下一个任务是保存 32 个用户寄存器。内核为每个进程分配一页内存作为 {\tt trapframe} 结构，该结构（除其他外）有空间保存 32 个用户寄存器
\lineref{kernel/proc.h:/^struct.trapframe/}。因为 {\tt satp} 仍然引用用户页表，所以 {\tt uservec} 需要 trapframe 映射在用户地址空间中。Xv6 将每个进程的 trapframe 映射到该进程用户页表中的虚拟地址 {\tt TRAPFRAME}（{\tt 0x3fffffe000}）；在 {\tt TRAMPOLINE} 下方一页。每个进程的 {\tt p->trapframe}
包含进程 trapframe 的内核虚拟地址。

{\tt uservec} 将寄存器 {\tt a0} 设置为地址 {\tt TRAPFRAME} 并将所有用户寄存器保存在那里。然后它从 {\tt sscratch} 检索用户 {\tt a0} 并将其保存在 trapframe 中。

内核先前初始化 trapframe 以包含对 {\tt uservec} 有用的一些值：当前进程的内核栈地址、当前 CPU 的 hartid、{\tt usertrap} 函数的地址以及内核页表的地址。{\tt uservec} 检索这些值，将 {\tt satp} 切换到内核页表，并跳转到 {\tt usertrap}，一个 C 函数。

{\tt usertrap} 的工作是确定陷入的原因，进行处理，然后返回
\lineref{kernel/trap.c:/^usertrap/}。
它首先将 {\tt stvec} 指向 {\tt kernelvec}，以便内核中的陷入将由 {\tt kernelvec} 而非 {\tt uservec} 处理。它保存 {\tt sepc} 寄存器（保存的用户程序计数器），以便将来返回用户空间时使用。如果陷入是系统调用，{\tt usertrap} 调用 {\tt syscall} 处理；如果是设备中断，调用 {\tt devintr}；如果是页错误，调用 {\tt vmfault}；否则就是异常（例如，使用了无效地址），内核会终止故障进程。系统调用路径将保存的用户程序计数器加 4，因为在系统调用的情况下，RISC-V 使程序指针指向 {\tt ecall} 指令，但用户代码需要从后续指令恢复执行。{\tt usertrap} 检查进程是否已被终止或应该让出 CPU（如果此陷入是定时器中断）。

返回用户空间的第一步是调用 {\tt prepare\_return}
{\lineref{kernel/trap.c:/^prepare/}}。
此函数设置 RISC-V 控制寄存器，为将来从用户空间陷入做准备：将 {\tt stvec}
设置为 {\tt uservec}，并准备 {\tt uservec} 依赖的 trapframe 字段。{\tt prepare\_return} 将 {\tt sepc} 设置为先前保存的用户程序计数器。最后，{\tt usertrap}
返回到 trampoline 页中的 {\tt userret}
\lineref{kernel/trampoline.S:/^userret/}，
在 {\tt a0} 中传递指向用户页表的指针。

{\tt userret} 将 {\tt satp} 切换到进程的用户页表。回想一下，用户页表映射了 trampoline 页和 {\tt TRAPFRAME}，但不映射内核的其他内容。trampoline 页在用户和内核页表中映射到相同的虚拟地址，这使得 {\tt userret} 在更改 {\tt satp} 后可以继续执行。从此刻起，{\tt userret} 唯一能使用的数据是寄存器内容和 trapframe 内容。{\tt userret} 将 {\tt TRAPFRAME} 地址加载到 {\tt a0} 中，通过 {\tt a0} 从 trapframe 恢复保存的用户寄存器，恢复保存的用户 {\tt a0}，然后执行 {\tt sret} 返回用户空间。

{\tt uservec} 和 {\tt userret} 用汇编语言编写，因为很难编写能够保存或恢复所有寄存器，或在切换页表后继续执行的 C 代码。

\section{代码：调用系统调用}

用户程序调用库函数以进行系统调用。例如，shell 使用此函数调用显示提示符（在 {\tt user/sh.c} 中）：

\begin{verbatim}
  write(2, "$ ", 2);
\end{verbatim}

这是库函数，在 {\tt user/usys.S} 中：

\begin{verbatim}
write:
 li a7, SYS_write
 ecall
 ret
\end{verbatim}

C 编译器为函数调用生成的代码将三个参数加载到寄存器 {\tt a0}、{\tt a1} 和 {\tt a2} 中。然后 {\tt write()} 函数将系统调用号 {\tt SYS\_write} (16) 加载到 {\tt a7} 中。内核将查看这些寄存器以找出打算进行什么系统调用以及参数是什么。\lstinline{ecall} 指令从用户空间陷入内核并导致
{\tt uservec}、{\tt usertrap}，然后 {\tt syscall} 执行。

此时，请阅读 {\tt kernel/syscall.c}、{\tt kernel/sysfile.c} 中的 {\tt sys\_write()}，以及 {\tt kernel/vm.c} 中的 {\tt copyout()}、{\tt copyin()} 和 {\tt copyinstr()}。

\indexcode{syscall}
\lineref{kernel/syscall.c:/^syscall/}
从陷阱帧中保存的
\texttt{a7}
检索系统调用号
并使用它索引到
{\tt syscalls}
\lineref{kernel/syscall.c:/syscalls/}。
在我们的示例中，
\texttt{a7}
包含
\indexcode{SYS_write}
\lineref{kernel/syscall.h:/SYS_write/}，
导致调用系统调用实现函数
\lstinline{sys_write}。

当 \lstinline{sys_write} 返回时，\lstinline{syscall} 将其返回值记录在 \lstinline{p->trapframe->a0} 中。这将导致原始用户空间对 {\tt write()} 的调用返回该值，因为 RISC-V 上的 C 调用约定将返回值放在 {\tt a0} 中。系统调用通常返回负数以指示错误，返回零或正数表示成功。

\section{代码：系统调用参数}

系统调用参数最初位于用户寄存器中，然后由内核陷阱代码移动到陷阱帧中。内核函数
\lstinline{argint}、
\lstinline{argaddr}
和
\lstinline{argfd}
从陷阱帧中检索第
\textit{n}
个系统调用参数作为整数、指针或文件描述符。

一些系统调用将指针作为参数传递，内核必须使用这些指针来读取或写入用户内存。例如，{\tt write} 系统调用向内核传递一个用户空间指针，指向要写入的数据。此类指针带来两个挑战。首先，用户程序可能存在缺陷或是恶意的，可能向内核传递无效指针，或旨在欺骗内核访问内核内存而非用户内存的指针。其次，xv6 内核页表映射与用户页表映射不同，因此内核不能使用普通指令从用户提供的地址加载或存储。

内核实现函数来安全地在用户提供的地址之间传输数据。{\tt fetchstr} 是一个例子 \lineref{kernel/syscall.c:/^fetchstr/}。
文件系统调用如
{\tt exec} 使用 {\tt fetchstr} 从用户空间检索字符串文件名参数。
\lstinline{fetchstr} 调用 \lstinline{copyinstr}
来完成艰苦的工作。

\indexcode{copyinstr}
\lineref{kernel/vm.c:/^copyinstr/} 将最多 \lstinline{max} 字节从用户页表 \lstinline{pagetable} 中的虚拟地址 \lstinline{srcva} 复制到 \lstinline{dst}。由于 \lstinline{pagetable} \textit{不是}当前页表，\lstinline{copyinstr} 使用 {\tt walkaddr}（调用 {\tt walk}）在 \lstinline{pagetable} 中查找 \lstinline{srcva}，产生物理地址 \lstinline{pa0}。内核的页表将所有物理 RAM 映射到等于 RAM 物理地址的虚拟地址。这允许 {\tt copyinstr} 直接从 {\tt pa0} 复制字符串字节到 {\tt dst}。{\tt walkaddr}
\lineref{kernel/vm.c:/^walkaddr/}
检查用户提供的虚拟地址是否是进程用户地址空间的一部分，因此程序不能欺骗内核读取其他内存。类似的函数 {\tt copyout} 将数据从内核复制到用户提供的地址。

\section{来自内核空间的陷入}
\label{s:ktraps}

请阅读 {\tt kernel/kernelvec.S} 和 {\tt kernel/trap.c} 中的 {\tt kerneltrap()}。

Xv6 以不同于用户代码陷入的方式处理内核代码的陷入。进入内核时，{\tt usertrap} 将 {\tt stvec} 指向 {\tt kernelvec} 处的汇编代码
\lineref{kernel/kernelvec.S:/^kernelvec/}。
因为 {\tt kernelvec} 只有在 xv6 已经在内核中时才会执行，所以 {\tt kernelvec} 可以依赖 {\tt satp} 被设置为内核页表，并且栈指针引用有效的内核栈。{\tt kernelvec} 将所有 32 个寄存器推入当前栈，稍后将从那里恢复它们，以便被中断的内核代码可以在不受干扰的情况下恢复。

{\tt kernelvec} 将被中断内核线程的寄存器保存在其栈上，这很合理，因为寄存器值属于该线程。如果陷入导致切换到不同线程，这一点尤其重要——在这种情况下，陷入实际上将从新线程的栈返回，将被中断线程的保存寄存器安全地留在其栈上。

{\tt kernelvec} 在保存寄存器后跳转到 {\tt kerneltrap}
\lineref{kernel/trap.c:/^kerneltrap/}。
{\tt kerneltrap} 只准备处理一种类型的陷入：设备中断。它调用 {\tt devintr}
\lineref{kernel/trap.c:/^devintr/}
来处理它们。如果陷入不是设备中断，它一定是异常，例如内核代码尝试使用无效指针。这只能由内核代码中的错误引起。内核在这种情况下无法恢复，因此它调用 {\tt panic()}，打印错误消息然后停止。

如果由于定时器中断而调用 {\tt kerneltrap}，并且进程的内核线程正在运行（相对于调度程序线程），{\tt kerneltrap} 调用 {\tt yield} 给其他线程运行的机会。在某个时候，其中一个线程会让出，让我们的线程及其 {\tt kerneltrap} 再次恢复。Chapter~\ref{CH:SCHED} 解释了 {\tt yield} 中发生的事情。

当 {\tt kerneltrap} 的工作完成时，它需要返回到被陷入中断的代码。因为 {\tt yield} 可能扰乱了 {\tt sepc} 和 {\tt sstatus} 中的先前模式，{\tt kerneltrap} 在开始时保存它们。它现在恢复那些控制寄存器并返回到 {\tt kernelvec}
\lineref{kernel/kernelvec.S:/call.kerneltrap$/}。
{\tt kernelvec} 从栈中弹出保存的寄存器并执行 {\tt sret}，它将 {\tt sepc} 复制到 {\tt pc} 并恢复被中断的内核代码。

Xv6 在 CPU 从用户空间进入内核时将 CPU 的 {\tt stvec} 设置为 {\tt kernelvec}；你可以在 {\tt usertrap}
\lineref{kernel/trap.c:/DOC:.kernelvec/}
中看到这一点。但有一个时间窗口，内核已经开始执行但 {\tt stvec} 仍设置为 {\tt uservec}，关键是该窗口内不能发生设备中断。幸运的是，RISC-V 在开始陷入时总是禁用中断，并且 {\tt usertrap} 在设置 {\tt stvec} 之后才重新启用它们。

\section{现实世界}

如果内核内存映射到每个进程的用户页表中（使用 \lstinline{PTE_U} 清除），则可以消除对 trampoline 页的需求。这也会消除从用户空间陷入内核时切换页表的需要。反过来，这将允许内核中的系统调用实现利用当前进程的用户内存被映射的优势，允许内核代码直接解引用用户指针。许多操作系统使用这些想法来提高效率。Xv6 避免它们是为了减少由于无意使用用户指针而导致内核安全错误的机会，并减少确保用户和内核虚拟地址不重叠所需的一些复杂性。

\section{练习}

\begin{enumerate}

\item {\tt trampoline.S} 和 {\tt kernelvec.S} 中的部分或全部代码可以用 C 而不是汇编器编写吗？

\item 有没有办法消除每个用户地址空间中的特殊 {\tt TRAPFRAME} 页映射？例如，可以修改 {\tt uservec} 以简单地将 32 个用户寄存器推入内核栈，或将它们存储在 {\tt proc} 结构中吗？

\item 可以修改 xv6 以消除特殊的 {\tt TRAMPOLINE} 页映射吗？

\end{enumerate}

\caption{处理来自用户代码的陷入的概要。}
\label{fig:usertrap}
\end{figure}

Xv6 根据陷入发生在内核执行还是用户代码执行中而有所不同地处理陷入。以下是用户代码陷入的处理流程；Section~\ref{s:ktraps} 将介绍内核代码的陷入处理。

如果用户程序进行系统调用（{\tt ecall} 指令），或执行了非法操作，或发生设备中断，则可能在用户空间执行时发生陷入。如 Figure~\ref{fig:usertrap} 所示，来自用户空间的高级别陷入路径是
{\tt uservec}
\lineref{kernel/trampoline.S:/^uservec/}，
然后是 {\tt usertrap}
\lineref{kernel/trap.c:/^usertrap/}；
当内核准备返回时，
{\tt usertrap}
返回给
{\tt userret}
\lineref{kernel/trampoline.S:/^userret/}，
它执行 {\tt sret}
返回到用户空间。

% talk about why RISC-V doesn't switch page tables

xv6 陷入处理设计的一个主要限制是 RISC-V 硬件在强制陷入时不会切换页表。这意味着 {\tt stvec} 中的陷入处理程序地址必须在用户页表中有有效映射，因为这是在陷入处理代码开始执行时生效的页表。此外，xv6 的陷入处理代码需要切换到内核页表；为了能够在切换后继续执行，内核页表也必须具有 {\tt stvec} 指向的处理程序的映射。

Xv6 使用 \indextext{trampoline}（跳板）页来满足这些要求。此页包含 {\tt uservec}，即 {\tt stvec} 指向的 xv6 陷入处理代码。trampoline 页在每个进程的页表中映射到虚拟地址 {\tt 0x3ffffff000}（称为 \indexcode{TRAMPOLINE}），这是虚拟地址空间中的最后一页，以便它位于程序自己使用的内存之上。trampoline 页在内核页表中映射到相同的虚拟地址。参见 Figure~\ref{fig:as} 和 Figure~\ref{fig:xv6_layout}。因为 trampoline 页映射在用户页表中，所以陷入可以在监管者模式下开始在那里执行。因为 trampoline 页在内核地址空间中映射到相同的地址，所以陷入处理程序可以在切换到内核页表后继续执行。

{\tt uservec} 陷阱处理程序的代码位于 {\tt trampoline.S}
\lineref{kernel/trampoline.S:/^uservec/}。
当 {\tt uservec} 开始时，所有 32 个寄存器都包含被中断用户代码的值。这 32 个值需要保存在内存中的某个位置，以便内核在稍后返回用户空间之前恢复它们。存储到内存需要使用寄存器来保存目标地址，但此时没有通用寄存器可用！幸运的是，RISC-V 以 {\tt sscratch} 寄存器的形式提供了帮助。{\tt uservec} 开头的 {\tt csrw} 指令将 {\tt a0} 保存在 {\tt sscratch} 中。现在 {\tt uservec} 有一个寄存器（{\tt a0}）可以使用了。

{\tt uservec} 的下一个任务是保存 32 个用户寄存器。内核为每个进程分配一页内存作为 {\tt trapframe} 结构，该结构（除其他外）有空间保存 32 个用户寄存器
\lineref{kernel/proc.h:/^struct.trapframe/}。因为 {\tt satp} 仍然引用用户页表，所以 {\tt uservec} 需要 trapframe 映射在用户地址空间中。Xv6 将每个进程的 trapframe 映射到该进程用户页表中的虚拟地址 {\tt TRAPFRAME}（{\tt 0x3fffffe000}）；在 {\tt TRAMPOLINE} 下方一页。每个进程的 {\tt p->trapframe}
包含进程 trapframe 的内核虚拟地址。

{\tt uservec} 将寄存器 {\tt a0} 设置为地址 {\tt TRAPFRAME} 并将所有用户寄存器保存在那里。然后它从 {\tt sscratch} 检索用户 {\tt a0} 并将其保存在 trapframe 中。

内核先前初始化 trapframe 以包含对 {\tt uservec} 有用的一些值：当前进程的内核栈地址、当前 CPU 的 hartid、{\tt usertrap} 函数的地址以及内核页表的地址。{\tt uservec} 检索这些值，将 {\tt satp} 切换到内核页表，并跳转到 {\tt usertrap}，一个 C 函数。

{\tt usertrap} 的工作是确定陷入的原因，进行处理，然后返回
\lineref{kernel/trap.c:/^usertrap/}。
它首先将 {\tt stvec} 指向 {\tt kernelvec}，以便内核中的陷入将由 {\tt kernelvec} 而非 {\tt uservec} 处理。它保存 {\tt sepc} 寄存器（保存的用户程序计数器），以便将来返回用户空间时使用。如果陷入是系统调用，{\tt usertrap} 调用 {\tt syscall} 处理；如果是设备中断，调用 {\tt devintr}；如果是页错误，调用 {\tt vmfault}；否则就是异常（例如，使用了无效地址），内核会终止故障进程。系统调用路径将保存的用户程序计数器加 4，因为在系统调用的情况下，RISC-V 使程序指针指向 {\tt ecall} 指令，但用户代码需要从后续指令恢复执行。{\tt usertrap} 检查进程是否已被终止或应该让出 CPU（如果此陷入是定时器中断）。

返回用户空间的第一步是调用 {\tt prepare\_return}
{\lineref{kernel/trap.c:/^prepare/}}。
此函数设置 RISC-V 控制寄存器，为将来从用户空间陷入做准备：将 {\tt stvec}
设置为 {\tt uservec}，并准备 {\tt uservec} 依赖的 trapframe 字段。{\tt prepare\_return} 将 {\tt sepc} 设置为先前保存的用户程序计数器。最后，{\tt usertrap}
返回到 trampoline 页中的 {\tt userret}
\lineref{kernel/trampoline.S:/^userret/}，
在 {\tt a0} 中传递指向用户页表的指针。

{\tt userret} 将 {\tt satp} 切换到进程的用户页表。回想一下，用户页表映射了 trampoline 页和 {\tt TRAPFRAME}，但不映射内核的其他内容。trampoline 页在用户和内核页表中映射到相同的虚拟地址，这使得 {\tt userret} 在更改 {\tt satp} 后可以继续执行。从此刻起，{\tt userret} 唯一能使用的数据是寄存器内容和 trapframe 内容。{\tt userret} 将 {\tt TRAPFRAME} 地址加载到 {\tt a0} 中，通过 {\tt a0} 从 trapframe 恢复保存的用户寄存器，恢复保存的用户 {\tt a0}，然后执行 {\tt sret} 返回用户空间。

{\tt uservec} 和 {\tt userret} 用汇编语言编写，因为很难编写能够保存或恢复所有寄存器，或在切换页表后继续执行的 C 代码。

\section{代码：调用系统调用}

用户程序调用库函数以进行系统调用。例如，shell 使用此函数调用显示提示符（在 {\tt user/sh.c} 中）：

\begin{verbatim}
  write(2, "$ ", 2);
\end{verbatim}

这是库函数，在 {\tt user/usys.S} 中：

\begin{verbatim}
write:
 li a7, SYS_write
 ecall
 ret
\end{verbatim}

C 编译器为函数调用生成的代码将三个参数加载到寄存器 {\tt a0}、{\tt a1} 和 {\tt a2} 中。然后 {\tt write()} 函数将系统调用号 {\tt SYS\_write} (16) 加载到 {\tt a7} 中。内核将查看这些寄存器以找出打算进行什么系统调用以及参数是什么。\lstinline{ecall} 指令从用户空间陷入内核并导致
{\tt uservec}、{\tt usertrap}，然后 {\tt syscall} 执行。

此时，请阅读 {\tt kernel/syscall.c}、{\tt kernel/sysfile.c} 中的 {\tt sys\_write()}，以及 {\tt kernel/vm.c} 中的 {\tt copyout()}、{\tt copyin()} 和 {\tt copyinstr()}。

\indexcode{syscall}
\lineref{kernel/syscall.c:/^syscall/}
从陷阱帧中保存的
\texttt{a7}
检索系统调用号
并使用它索引到
{\tt syscalls}
\lineref{kernel/syscall.c:/syscalls/}。
在我们的示例中，
\texttt{a7}
包含
\indexcode{SYS_write}
\lineref{kernel/syscall.h:/SYS_write/}，
导致调用系统调用实现函数
\lstinline{sys_write}。

当 \lstinline{sys_write} 返回时，\lstinline{syscall} 将其返回值记录在 \lstinline{p->trapframe->a0} 中。这将导致原始用户空间对 {\tt write()} 的调用返回该值，因为 RISC-V 上的 C 调用约定将返回值放在 {\tt a0} 中。系统调用通常返回负数以指示错误，返回零或正数表示成功。

\section{代码：系统调用参数}

系统调用参数最初位于用户寄存器中，然后由内核陷阱代码移动到陷阱帧中。内核函数
\lstinline{argint}、
\lstinline{argaddr}
和
\lstinline{argfd}
从陷阱帧中检索第
\textit{n}
个系统调用参数作为整数、指针或文件描述符。

一些系统调用将指针作为参数传递，内核必须使用这些指针来读取或写入用户内存。例如，{\tt write} 系统调用向内核传递一个用户空间指针，指向要写入的数据。此类指针带来两个挑战。首先，用户程序可能存在缺陷或是恶意的，可能向内核传递无效指针，或旨在欺骗内核访问内核内存而非用户内存的指针。其次，xv6 内核页表映射与用户页表映射不同，因此内核不能使用普通指令从用户提供的地址加载或存储。

内核实现函数来安全地在用户提供的地址之间传输数据。{\tt fetchstr} 是一个例子 \lineref{kernel/syscall.c:/^fetchstr/}。
文件系统调用如
{\tt exec} 使用 {\tt fetchstr} 从用户空间检索字符串文件名参数。
\lstinline{fetchstr} 调用 \lstinline{copyinstr}
来完成艰苦的工作。

\indexcode{copyinstr}
\lineref{kernel/vm.c:/^copyinstr/} 将最多 \lstinline{max} 字节从用户页表 \lstinline{pagetable} 中的虚拟地址 \lstinline{srcva} 复制到 \lstinline{dst}。由于 \lstinline{pagetable} \textit{不是}当前页表，\lstinline{copyinstr} 使用 {\tt walkaddr}（调用 {\tt walk}）在 \lstinline{pagetable} 中查找 \lstinline{srcva}，产生物理地址 \lstinline{pa0}。内核的页表将所有物理 RAM 映射到等于 RAM 物理地址的虚拟地址。这允许 {\tt copyinstr} 直接从 {\tt pa0} 复制字符串字节到 {\tt dst}。{\tt walkaddr}
\lineref{kernel/vm.c:/^walkaddr/}
检查用户提供的虚拟地址是否是进程用户地址空间的一部分，因此程序不能欺骗内核读取其他内存。类似的函数 {\tt copyout} 将数据从内核复制到用户提供的地址。

\section{来自内核空间的陷入}
\label{s:ktraps}

请阅读 {\tt kernel/kernelvec.S} 和 {\tt kernel/trap.c} 中的 {\tt kerneltrap()}。

Xv6 以不同于用户代码陷入的方式处理内核代码的陷入。进入内核时，{\tt usertrap} 将 {\tt stvec} 指向 {\tt kernelvec} 处的汇编代码
\lineref{kernel/kernelvec.S:/^kernelvec/}。
因为 {\tt kernelvec} 只有在 xv6 已经在内核中时才会执行，所以 {\tt kernelvec} 可以依赖 {\tt satp} 被设置为内核页表，并且栈指针引用有效的内核栈。{\tt kernelvec} 将所有 32 个寄存器推入当前栈，稍后将从那里恢复它们，以便被中断的内核代码可以在不受干扰的情况下恢复。

{\tt kernelvec} 将被中断内核线程的寄存器保存在其栈上，这很合理，因为寄存器值属于该线程。如果陷入导致切换到不同线程，这一点尤其重要——在这种情况下，陷入实际上将从新线程的栈返回，将被中断线程的保存寄存器安全地留在其栈上。

{\tt kernelvec} 在保存寄存器后跳转到 {\tt kerneltrap}
\lineref{kernel/trap.c:/^kerneltrap/}。
{\tt kerneltrap} 只准备处理一种类型的陷入：设备中断。它调用 {\tt devintr}
\lineref{kernel/trap.c:/^devintr/}
来处理它们。如果陷入不是设备中断，它一定是异常，例如内核代码尝试使用无效指针。这只能由内核代码中的错误引起。内核在这种情况下无法恢复，因此它调用 {\tt panic()}，打印错误消息然后停止。

如果由于定时器中断而调用 {\tt kerneltrap}，并且进程的内核线程正在运行（相对于调度程序线程），{\tt kerneltrap} 调用 {\tt yield} 给其他线程运行的机会。在某个时候，其中一个线程会让出，让我们的线程及其 {\tt kerneltrap} 再次恢复。Chapter~\ref{CH:SCHED} 解释了 {\tt yield} 中发生的事情。

当 {\tt kerneltrap} 的工作完成时，它需要返回到被陷入中断的代码。因为 {\tt yield} 可能扰乱了 {\tt sepc} 和 {\tt sstatus} 中的先前模式，{\tt kerneltrap} 在开始时保存它们。它现在恢复那些控制寄存器并返回到 {\tt kernelvec}
\lineref{kernel/kernelvec.S:/call.kerneltrap$/}。
{\tt kernelvec} 从栈中弹出保存的寄存器并执行 {\tt sret}，它将 {\tt sepc} 复制到 {\tt pc} 并恢复被中断的内核代码。

Xv6 在 CPU 从用户空间进入内核时将 CPU 的 {\tt stvec} 设置为 {\tt kernelvec}；你可以在 {\tt usertrap}
\lineref{kernel/trap.c:/DOC:.kernelvec/}
中看到这一点。但有一个时间窗口，内核已经开始执行但 {\tt stvec} 仍设置为 {\tt uservec}，关键是该窗口内不能发生设备中断。幸运的是，RISC-V 在开始陷入时总是禁用中断，并且 {\tt usertrap} 在设置 {\tt stvec} 之后才重新启用它们。

\section{现实世界}

如果内核内存映射到每个进程的用户页表中（使用 \lstinline{PTE_U} 清除），则可以消除对 trampoline 页的需求。这也会消除从用户空间陷入内核时切换页表的需要。反过来，这将允许内核中的系统调用实现利用当前进程的用户内存被映射的优势，允许内核代码直接解引用用户指针。许多操作系统使用这些想法来提高效率。Xv6 避免它们是为了减少由于无意使用用户指针而导致内核安全错误的机会，并减少确保用户和内核虚拟地址不重叠所需的一些复杂性。

\section{练习}

\begin{enumerate}

\item {\tt trampoline.S} 和 {\tt kernelvec.S} 中的部分或全部代码可以用 C 而不是汇编器编写吗？

\item 有没有办法消除每个用户地址空间中的特殊 {\tt TRAPFRAME} 页映射？例如，可以修改 {\tt uservec} 以简单地将 32 个用户寄存器推入内核栈，或将它们存储在 {\tt proc} 结构中吗？

\item 可以修改 xv6 以消除特殊的 {\tt TRAMPOLINE} 页映射吗？

\end{enumerate}

\caption{处理来自用户代码的陷入的概要。}
\label{fig:usertrap}
\end{figure}

Xv6 根据陷入发生在内核执行还是用户代码执行中而有所不同地处理陷入。以下是用户代码陷入的处理流程；Section~\ref{s:ktraps} 将介绍内核代码的陷入处理。

如果用户程序进行系统调用（{\tt ecall} 指令），或执行了非法操作，或发生设备中断，则可能在用户空间执行时发生陷入。如 Figure~\ref{fig:usertrap} 所示，来自用户空间的高级别陷入路径是
{\tt uservec}
\lineref{kernel/trampoline.S:/^uservec/}，
然后是 {\tt usertrap}
\lineref{kernel/trap.c:/^usertrap/}；
当内核准备返回时，
{\tt usertrap}
返回给
{\tt userret}
\lineref{kernel/trampoline.S:/^userret/}，
它执行 {\tt sret}
返回到用户空间。

% talk about why RISC-V doesn't switch page tables

xv6 陷入处理设计的一个主要限制是 RISC-V 硬件在强制陷入时不会切换页表。这意味着 {\tt stvec} 中的陷入处理程序地址必须在用户页表中有有效映射，因为这是在陷入处理代码开始执行时生效的页表。此外，xv6 的陷入处理代码需要切换到内核页表；为了能够在切换后继续执行，内核页表也必须具有 {\tt stvec} 指向的处理程序的映射。

Xv6 使用 \indextext{trampoline}（跳板）页来满足这些要求。此页包含 {\tt uservec}，即 {\tt stvec} 指向的 xv6 陷入处理代码。trampoline 页在每个进程的页表中映射到虚拟地址 {\tt 0x3ffffff000}（称为 \indexcode{TRAMPOLINE}），这是虚拟地址空间中的最后一页，以便它位于程序自己使用的内存之上。trampoline 页在内核页表中映射到相同的虚拟地址。参见 Figure~\ref{fig:as} 和 Figure~\ref{fig:xv6_layout}。因为 trampoline 页映射在用户页表中，所以陷入可以在监管者模式下开始在那里执行。因为 trampoline 页在内核地址空间中映射到相同的地址，所以陷入处理程序可以在切换到内核页表后继续执行。

{\tt uservec} 陷阱处理程序的代码位于 {\tt trampoline.S}
\lineref{kernel/trampoline.S:/^uservec/}。
当 {\tt uservec} 开始时，所有 32 个寄存器都包含被中断用户代码的值。这 32 个值需要保存在内存中的某个位置，以便内核在稍后返回用户空间之前恢复它们。存储到内存需要使用寄存器来保存目标地址，但此时没有通用寄存器可用！幸运的是，RISC-V 以 {\tt sscratch} 寄存器的形式提供了帮助。{\tt uservec} 开头的 {\tt csrw} 指令将 {\tt a0} 保存在 {\tt sscratch} 中。现在 {\tt uservec} 有一个寄存器（{\tt a0}）可以使用了。

{\tt uservec} 的下一个任务是保存 32 个用户寄存器。内核为每个进程分配一页内存作为 {\tt trapframe} 结构，该结构（除其他外）有空间保存 32 个用户寄存器
\lineref{kernel/proc.h:/^struct.trapframe/}。因为 {\tt satp} 仍然引用用户页表，所以 {\tt uservec} 需要 trapframe 映射在用户地址空间中。Xv6 将每个进程的 trapframe 映射到该进程用户页表中的虚拟地址 {\tt TRAPFRAME}（{\tt 0x3fffffe000}）；在 {\tt TRAMPOLINE} 下方一页。每个进程的 {\tt p->trapframe}
包含进程 trapframe 的内核虚拟地址。

{\tt uservec} 将寄存器 {\tt a0} 设置为地址 {\tt TRAPFRAME} 并将所有用户寄存器保存在那里。然后它从 {\tt sscratch} 检索用户 {\tt a0} 并将其保存在 trapframe 中。

内核先前初始化 trapframe 以包含对 {\tt uservec} 有用的一些值：当前进程的内核栈地址、当前 CPU 的 hartid、{\tt usertrap} 函数的地址以及内核页表的地址。{\tt uservec} 检索这些值，将 {\tt satp} 切换到内核页表，并跳转到 {\tt usertrap}，一个 C 函数。

{\tt usertrap} 的工作是确定陷入的原因，进行处理，然后返回
\lineref{kernel/trap.c:/^usertrap/}。
它首先将 {\tt stvec} 指向 {\tt kernelvec}，以便内核中的陷入将由 {\tt kernelvec} 而非 {\tt uservec} 处理。它保存 {\tt sepc} 寄存器（保存的用户程序计数器），以便将来返回用户空间时使用。如果陷入是系统调用，{\tt usertrap} 调用 {\tt syscall} 处理；如果是设备中断，调用 {\tt devintr}；如果是页错误，调用 {\tt vmfault}；否则就是异常（例如，使用了无效地址），内核会终止故障进程。系统调用路径将保存的用户程序计数器加 4，因为在系统调用的情况下，RISC-V 使程序指针指向 {\tt ecall} 指令，但用户代码需要从后续指令恢复执行。{\tt usertrap} 检查进程是否已被终止或应该让出 CPU（如果此陷入是定时器中断）。

返回用户空间的第一步是调用 {\tt prepare\_return}
{\lineref{kernel/trap.c:/^prepare/}}。
此函数设置 RISC-V 控制寄存器，为将来从用户空间陷入做准备：将 {\tt stvec}
设置为 {\tt uservec}，并准备 {\tt uservec} 依赖的 trapframe 字段。{\tt prepare\_return} 将 {\tt sepc} 设置为先前保存的用户程序计数器。最后，{\tt usertrap}
返回到 trampoline 页中的 {\tt userret}
\lineref{kernel/trampoline.S:/^userret/}，
在 {\tt a0} 中传递指向用户页表的指针。

{\tt userret} 将 {\tt satp} 切换到进程的用户页表。回想一下，用户页表映射了 trampoline 页和 {\tt TRAPFRAME}，但不映射内核的其他内容。trampoline 页在用户和内核页表中映射到相同的虚拟地址，这使得 {\tt userret} 在更改 {\tt satp} 后可以继续执行。从此刻起，{\tt userret} 唯一能使用的数据是寄存器内容和 trapframe 内容。{\tt userret} 将 {\tt TRAPFRAME} 地址加载到 {\tt a0} 中，通过 {\tt a0} 从 trapframe 恢复保存的用户寄存器，恢复保存的用户 {\tt a0}，然后执行 {\tt sret} 返回用户空间。

{\tt uservec} 和 {\tt userret} 用汇编语言编写，因为很难编写能够保存或恢复所有寄存器，或在切换页表后继续执行的 C 代码。

\section{代码：调用系统调用}

用户程序调用库函数以进行系统调用。例如，shell 使用此函数调用显示提示符（在 {\tt user/sh.c} 中）：

\begin{verbatim}
  write(2, "$ ", 2);
\end{verbatim}

这是库函数，在 {\tt user/usys.S} 中：

\begin{verbatim}
write:
 li a7, SYS_write
 ecall
 ret
\end{verbatim}

C 编译器为函数调用生成的代码将三个参数加载到寄存器 {\tt a0}、{\tt a1} 和 {\tt a2} 中。然后 {\tt write()} 函数将系统调用号 {\tt SYS\_write} (16) 加载到 {\tt a7} 中。内核将查看这些寄存器以找出打算进行什么系统调用以及参数是什么。\lstinline{ecall} 指令从用户空间陷入内核并导致
{\tt uservec}、{\tt usertrap}，然后 {\tt syscall} 执行。

此时，请阅读 {\tt kernel/syscall.c}、{\tt kernel/sysfile.c} 中的 {\tt sys\_write()}，以及 {\tt kernel/vm.c} 中的 {\tt copyout()}、{\tt copyin()} 和 {\tt copyinstr()}。

\indexcode{syscall}
\lineref{kernel/syscall.c:/^syscall/}
从陷阱帧中保存的
\texttt{a7}
检索系统调用号
并使用它索引到
{\tt syscalls}
\lineref{kernel/syscall.c:/syscalls/}。
在我们的示例中，
\texttt{a7}
包含
\indexcode{SYS_write}
\lineref{kernel/syscall.h:/SYS_write/}，
导致调用系统调用实现函数
\lstinline{sys_write}。

当 \lstinline{sys_write} 返回时，\lstinline{syscall} 将其返回值记录在 \lstinline{p->trapframe->a0} 中。这将导致原始用户空间对 {\tt write()} 的调用返回该值，因为 RISC-V 上的 C 调用约定将返回值放在 {\tt a0} 中。系统调用通常返回负数以指示错误，返回零或正数表示成功。

\section{代码：系统调用参数}

系统调用参数最初位于用户寄存器中，然后由内核陷阱代码移动到陷阱帧中。内核函数
\lstinline{argint}、
\lstinline{argaddr}
和
\lstinline{argfd}
从陷阱帧中检索第
\textit{n}
个系统调用参数作为整数、指针或文件描述符。

一些系统调用将指针作为参数传递，内核必须使用这些指针来读取或写入用户内存。例如，{\tt write} 系统调用向内核传递一个用户空间指针，指向要写入的数据。此类指针带来两个挑战。首先，用户程序可能存在缺陷或是恶意的，可能向内核传递无效指针，或旨在欺骗内核访问内核内存而非用户内存的指针。其次，xv6 内核页表映射与用户页表映射不同，因此内核不能使用普通指令从用户提供的地址加载或存储。

内核实现函数来安全地在用户提供的地址之间传输数据。{\tt fetchstr} 是一个例子 \lineref{kernel/syscall.c:/^fetchstr/}。
文件系统调用如
{\tt exec} 使用 {\tt fetchstr} 从用户空间检索字符串文件名参数。
\lstinline{fetchstr} 调用 \lstinline{copyinstr}
来完成艰苦的工作。

\indexcode{copyinstr}
\lineref{kernel/vm.c:/^copyinstr/} 将最多 \lstinline{max} 字节从用户页表 \lstinline{pagetable} 中的虚拟地址 \lstinline{srcva} 复制到 \lstinline{dst}。由于 \lstinline{pagetable} \textit{不是}当前页表，\lstinline{copyinstr} 使用 {\tt walkaddr}（调用 {\tt walk}）在 \lstinline{pagetable} 中查找 \lstinline{srcva}，产生物理地址 \lstinline{pa0}。内核的页表将所有物理 RAM 映射到等于 RAM 物理地址的虚拟地址。这允许 {\tt copyinstr} 直接从 {\tt pa0} 复制字符串字节到 {\tt dst}。{\tt walkaddr}
\lineref{kernel/vm.c:/^walkaddr/}
检查用户提供的虚拟地址是否是进程用户地址空间的一部分，因此程序不能欺骗内核读取其他内存。类似的函数 {\tt copyout} 将数据从内核复制到用户提供的地址。

\section{来自内核空间的陷入}
\label{s:ktraps}

请阅读 {\tt kernel/kernelvec.S} 和 {\tt kernel/trap.c} 中的 {\tt kerneltrap()}。

Xv6 以不同于用户代码陷入的方式处理内核代码的陷入。进入内核时，{\tt usertrap} 将 {\tt stvec} 指向 {\tt kernelvec} 处的汇编代码
\lineref{kernel/kernelvec.S:/^kernelvec/}。
因为 {\tt kernelvec} 只有在 xv6 已经在内核中时才会执行，所以 {\tt kernelvec} 可以依赖 {\tt satp} 被设置为内核页表，并且栈指针引用有效的内核栈。{\tt kernelvec} 将所有 32 个寄存器推入当前栈，稍后将从那里恢复它们，以便被中断的内核代码可以在不受干扰的情况下恢复。

{\tt kernelvec} 将被中断内核线程的寄存器保存在其栈上，这很合理，因为寄存器值属于该线程。如果陷入导致切换到不同线程，这一点尤其重要——在这种情况下，陷入实际上将从新线程的栈返回，将被中断线程的保存寄存器安全地留在其栈上。

{\tt kernelvec} 在保存寄存器后跳转到 {\tt kerneltrap}
\lineref{kernel/trap.c:/^kerneltrap/}。
{\tt kerneltrap} 只准备处理一种类型的陷入：设备中断。它调用 {\tt devintr}
\lineref{kernel/trap.c:/^devintr/}
来处理它们。如果陷入不是设备中断，它一定是异常，例如内核代码尝试使用无效指针。这只能由内核代码中的错误引起。内核在这种情况下无法恢复，因此它调用 {\tt panic()}，打印错误消息然后停止。

如果由于定时器中断而调用 {\tt kerneltrap}，并且进程的内核线程正在运行（相对于调度程序线程），{\tt kerneltrap} 调用 {\tt yield} 给其他线程运行的机会。在某个时候，其中一个线程会让出，让我们的线程及其 {\tt kerneltrap} 再次恢复。Chapter~\ref{CH:SCHED} 解释了 {\tt yield} 中发生的事情。

当 {\tt kerneltrap} 的工作完成时，它需要返回到被陷入中断的代码。因为 {\tt yield} 可能扰乱了 {\tt sepc} 和 {\tt sstatus} 中的先前模式，{\tt kerneltrap} 在开始时保存它们。它现在恢复那些控制寄存器并返回到 {\tt kernelvec}
\lineref{kernel/kernelvec.S:/call.kerneltrap$/}。
{\tt kernelvec} 从栈中弹出保存的寄存器并执行 {\tt sret}，它将 {\tt sepc} 复制到 {\tt pc} 并恢复被中断的内核代码。

Xv6 在 CPU 从用户空间进入内核时将 CPU 的 {\tt stvec} 设置为 {\tt kernelvec}；你可以在 {\tt usertrap}
\lineref{kernel/trap.c:/DOC:.kernelvec/}
中看到这一点。但有一个时间窗口，内核已经开始执行但 {\tt stvec} 仍设置为 {\tt uservec}，关键是该窗口内不能发生设备中断。幸运的是，RISC-V 在开始陷入时总是禁用中断，并且 {\tt usertrap} 在设置 {\tt stvec} 之后才重新启用它们。

\section{现实世界}

如果内核内存映射到每个进程的用户页表中（使用 \lstinline{PTE_U} 清除），则可以消除对 trampoline 页的需求。这也会消除从用户空间陷入内核时切换页表的需要。反过来，这将允许内核中的系统调用实现利用当前进程的用户内存被映射的优势，允许内核代码直接解引用用户指针。许多操作系统使用这些想法来提高效率。Xv6 避免它们是为了减少由于无意使用用户指针而导致内核安全错误的机会，并减少确保用户和内核虚拟地址不重叠所需的一些复杂性。

\section{练习}

\begin{enumerate}

\item {\tt trampoline.S} 和 {\tt kernelvec.S} 中的部分或全部代码可以用 C 而不是汇编器编写吗？

\item 有没有办法消除每个用户地址空间中的特殊 {\tt TRAPFRAME} 页映射？例如，可以修改 {\tt uservec} 以简单地将 32 个用户寄存器推入内核栈，或将它们存储在 {\tt proc} 结构中吗？

\item 可以修改 xv6 以消除特殊的 {\tt TRAMPOLINE} 页映射吗？

\end{enumerate}

\chapter{页错误}
\label{CH:PGFAULTS}
%%

当访问页表中无映射的虚拟地址，或映射的 \lstinline{PTE_V} 标志已清除的地址，或映射的权限位（\lstinline{PTE_R}、\lstinline{PTE_W}、\lstinline{PTE_X}、\lstinline{PTE_U}）禁止当前操作时，RISC-V CPU 会触发页错误异常。RISC-V 区分三种页错误：加载页错误（由加载指令引起）、存储页错误（由存储指令引起）和指令页错误（由取指执行指令引起）。\lstinline{scause} 寄存器指示页错误的类型，\indexcode{stval} 寄存器包含无法翻译的地址。

页表与页错误的组合是一个强大的工具。页表在内核虚拟地址与物理地址之间提供了一层间接层，使内核能够控制地址空间的结构和内容。页错误则允许内核拦截加载和存储操作，并通过修改页表来实时指定这些引用所指向的数据。内核可以利用这些能力提高效率：例如，写时复制fork允许内核在父进程和子进程之间透明地共享内存，避免复制双方都不会写入的页面。应用程序员也能从中受益。一种可能性是内存映射文件，内核通过分页机制将文件内容映射到应用程序的地址空间中，在响应页错误时透明地读取文件块。另一种是惰性内存分配，允许程序申请巨大的虚拟地址空间，但只需为实际读写的页面支付分配物理内存的成本。xv6 仅将页错误用于一个目的：惰性分配。

在继续之前，请阅读 {\tt kernel/sysproc.c} 中的函数 {\tt sys\_sbrk()} 和 {\tt kernel/vm.c} 中的 {\tt vmfault}。在 {\tt kernel/trap.c} 和 {\tt kernel/vm.c} 中搜索对 {\tt vmfault} 的调用。

\section{惰性分配}
\label{sec:lazy}

xv6 的 \indextext{lazy allocation} 包含两部分。首先，当应用程序通过调用带有 \lstinline{SBRK_LAZY} 标志的 \lstinline{sbrk} 请求内存时，内核仅记录大小的增加，不为新范围的虚拟地址分配物理内存或创建 PTE。其次，当在这些新地址之一上发生页错误时，内核分配一页物理内存并将其映射到页表中。内核透明地实现惰性分配：应用程序无需修改即可受益。

惰性分配对应用程序很方便，因为它们不必准确预测需要多少内存。例如，应用程序可能要处理输入，但事先不知道输入有多大。使用惰性分配，应用程序可以为最坏情况申请内存，但不必为此付费：内核无需为应用程序从未使用的页面做任何工作。

此外，如果应用程序要求大幅增加地址空间，那么没有惰性分配的 \lstinline{sbrk} 会很昂贵：如果应用程序申请 1 GB 内存，内核必须分配并清零 262,144 个 4096 字节的物理页面。惰性分配允许将这种成本分散到时间上。但另一方面，惰性分配会产生页错误的额外开销，这涉及用户/内核态切换。操作系统可以通过每次页错误分配一批连续页面而非单个页面，以及为这种页错误专门优化内核进入/退出代码来降低这种成本（尽管 xv6 两者都不做）。

另一方面，当为惰性分配页面发生页错误时，内核可能发现没有空闲内存可分配。在这种情况下，内核无法向应用程序返回内存不足错误，只能杀死应用程序。对于希望在分配失败时收到错误的应用程序，xv6 允许通过使用 \lstinline{SBRK_EAGER} 标志调用 \lstinline{sbrk} 来立即分配内存。

\section{代码}

系统调用 {\tt sbrk(n)} 通过 {\tt n} 字节增加（如果 {\tt n} 为负则减少）进程的内存大小，并返回新分配区域的起始地址（即旧大小）。内核实现是 \lstinline{sys_sbrk}
\lineref{kernel/sysproc.c:/^sys_sbrk/}。

如果应用程序指定 \lstinline{SBRK_EAGER}，系统调用由函数
\lstinline{growproc}
\lineref{kernel/proc.c:/^growproc/}
实现。\lstinline{growproc} 调用 \lstinline{uvmalloc}。
\lstinline{uvmalloc}
\lineref{kernel/vm.c:/^uvmalloc/}
使用 {\tt kalloc} 分配物理内存，将分配的内存清零，并使用 {\tt mappages} 将 PTE 添加到用户页表。

如果应用程序惰性分配内存，\lstinline{sys_sbrk}
只将进程的大小
(\lstinline{myproc()->sz}) 增加 {\tt n} 并返回旧大小；它不分配物理内存或向进程页表添加 PTE。

当进程对缺少有效页表映射的虚拟地址进行加载或存储时，CPU 将引发 \indextext{page-fault exception}。
\lstinline{usertrap} 检查此情况
\lineref{kernel/trap.c:/page fault/}
并调用 \lstinline{vmfault}
\lineref{kernel/vm.c:/^vmfault/}
来处理页错误。\lstinline{vmfault}
检查出错地址是否在 {\tt sbrk} 先前授予的区域内，
使用 {\tt kalloc} 分配一页物理内存，
将分配的页面清零，并使用
{\tt mappages} 将 PTE 添加到用户页表。Xv6 在 PTE 中设置新页面的 \lstinline{PTE_W}、\lstinline{PTE_R}、\lstinline{PTE_U} 和 \lstinline{PTE_V} 标志。然后，\lstinline{usertrap} 在导致错误的指令处恢复进程。因为 PTE 现在有效，重新执行的加载或存储指令将不会出错。

如果应用程序使用 {\tt sbrk} 释放内存，\lstinline{sys_sbrk}
调用 \lstinline{shrinkproc}，后者调用 \lstinline{uvmdealloc}。实际工作由 {\tt uvmunmap} \lineref{kernel/vm.c:/^uvmunmap/} 完成，它使用 {\tt walk} 查找 PTE。由于某些页面可能从未被进程使用，因此从未由 {\tt vmfault} 分配，{\tt uvmunmap} 跳过没有 \lstinline{PTE_V} 标志的 PTE。如果 PTE 映射有效，{\tt uvmunmap} 调用 {\tt kfree} 释放它引用的物理内存。

请注意，xv6 使用进程页表不仅告诉硬件如何映射用户虚拟地址，还作为分配给该进程的物理内存页面的唯一记录。这就是在 {\tt uvmunmap} 中释放用户内存时需要检查用户页表的原因。

\section{现实世界：写时复制（COW）fork}

许多内核（尽管不是 xv6）使用页错误来实现
\indextext{copy-on-write (COW) fork}。{\tt fork} 系统调用承诺子进程看到的内存的初始内容与 fork 时父进程的内存相同。实现这一点的一种方法是将父进程的整个内存复制到为子进程新分配的物理内存中；这就是 xv6 所做的。复制可能很慢，如果子进程可以共享父进程的物理内存，效率会更高。然而，这种简单实现不会奏效，因为它会导致父进程和子进程通过写入共享的栈和堆来相互干扰执行。

写时复制fork通过适当使用页表权限和页错误使父进程和子进程安全地共享物理内存。基本计划是让父进程和子进程最初共享所有物理页面，但每个都将其映射为只读（\lstinline{PTE_W} 标志清除）。然后父进程和子进程都可以从共享物理内存中读取。如果任何一方写入共享页面，RISC-V CPU 会触发页错误异常。支持 COW 的内核会响应分配新物理内存页面并将共享页面复制到新页面中。然后内核会更改错误进程的页表中的相关 PTE 以指向副本并允许读写，然后在导致错误的指令处恢复错误进程。因为 PTE 现在允许写入，重新执行的存储指令将不会出错，并将修改页面的私有副本而不是共享页面。

写时复制需要簿记来帮助决定何时可以释放物理页面，因为每个页面可以根据 fork、页错误、exec 和 exit 的历史被不同数量的页表引用。这种记账允许一个重要的优化：如果进程产生存储页错误且物理页面仅从该进程的页表引用，则不需要复制。

写时复制使 \lstinline{fork} 更快，因为 \lstinline{fork} 不需要复制内存。某些内存稍后写入时必须复制，但通常大部分内存永远不需要复制。一个常见例子是
\lstinline{fork} 后跟 \lstinline{exec}：
\lstinline{fork} 之后可能写入几页，
但随后子进程的 \lstinline{exec} 释放从父进程继承的大部分内存。写时复制 \lstinline{fork} 消除了复制这些内存的需要。
此外，COW fork 是透明的：
应用程序无需修改即可受益。

\section{现实世界：按需分页}

另一个利用页错误的广泛使用的功能是
\indextext{demand paging}。在 \lstinline{exec} 系统调用中，xv6 在启动应用程序之前将所有应用程序的文本和数据加载到内存中。由于应用程序可能很大，从磁盘读取需要时间，这种启动成本对用户来说可能很明显。为了减少启动时间，现代内核最初不将可执行文件加载到内存中，而只是创建用户页表并将所有 PTE 标记为无效。内核启动程序运行；每次程序第一次使用页面时，发生页错误，内核响应从磁盘读取页面内容并将其映射到用户地址空间。像 COW fork 和惰性分配一样，内核可以透明地为应用程序实现此功能。

在计算机上运行的程序可能需要比计算机拥有的 RAM 更多的内存。为了优雅地应对，操作系统可以实现 \indextext{paging to disk}。想法是只在 RAM 中存储一小部分用户页面，其余存储在磁盘上的 \indextext{paging area} 中。内核将对应于存储在分页区域中（因此不在 RAM 中）的内存的 PTE 标记为无效。如果应用程序尝试使用已 \emph{分页出去} 到磁盘的页面之一，应用程序将产生页错误，页面必须被 \emph{分页进来}：内核陷阱处理程序将分配一页物理 RAM，将页面从磁盘读入 RAM，并修改相关 PTE 以指向 RAM。

如果需要分页进入页面，但没有空闲物理 RAM 怎么办？在这种情况下，内核必须首先通过将其分页出去或 \emph{驱逐} 到磁盘上的分页区域，并将引用该物理页面的 PTE 标记为无效来释放物理页面。驱逐是昂贵的，因此如果它不频繁，分页性能最好：如果应用程序只使用其内存页面的一个子集，且这些子集的并集适合 RAM。这个属性通常被称为具有良好的引用局部性。与许多虚拟内存技术一样，内核通常以实现为对应用程序透明的方式实现分页到磁盘。

计算机通常在很少有或没有 \emph{空闲} 物理内存的情况下运行，无论硬件提供多少 RAM。例如，云提供商将许多客户多路复用到单台机器上以有效使用其硬件。另一个例子是用户在少量物理内存的智能手机上运行许多应用程序。在这种情况下，分配页面可能需要首先驱逐现有页面。因此，当空闲物理内存稀缺时，分配是昂贵的。

当空闲内存稀缺且程序只积极使用其分配内存的一部分时，惰性分配和按需分页特别有利。这些技术还可以避免分配或加载页面但从未使用或在可以使用之前被驱逐时浪费的工作。

\section{现实世界：内存映射文件}

结合分页和页错误异常的其他功能包括自动扩展栈和 \indextext{memory-mapped files}，这是程序使用 \texttt{mmap} 系统调用映射到其地址空间的文件，以便程序可以使用加载和存储指令读写它们。

% "virtual" memory, eviction, page-in
% lazy allocation
% auto stack expansion
% guard pages
% mmap files
% cow fork
% shared text, shared libraries
% demand paging of text
% virtual machine migration
% distributed shared memory
% fast IPC
% zero-copy write
% DPDK
% (unified block / page cache)

%%
\section{练习}
%%

\begin{enumerate}

\item 编写一个用户程序，通过调用
\lstinline{sbrk(1)}
将其地址空间增加一个字节。
运行该程序并调查程序调用
\lstinline{sbrk}
之前和调用
\lstinline{sbrk}
之后的页表。
内核分配了多少空间？新内存的
PTE
包含什么？

\item 实现 COW fork。

\item 实现 {\tt mmap}。

\end{enumerate}

\chapter{中断和设备驱动程序}
\label{CH:INTERRUPT}

\indextext{驱动程序}
是操作系统中管理特定设备的代码：
它配置设备硬件，
指示设备执行操作，
处理产生的中断，
并与使用该设备的进程交互。
驱动程序代码可能比较棘手，
因为驱动程序与设备并发执行，
而且通常与使用该设备的进程并发执行。
此外，驱动程序必须理解设备的硬件接口，
这可能很复杂且文档不完善。

需要操作系统关注的设备通常可以
配置为生成中断（interrupt），这是陷入（trap）的一种类型。
内核陷入处理代码识别到设备
产生了中断并调用驱动程序的中断处理程序；
在 xv6 中，这个分发发生在 {\tt devintr} \lineref{kernel/trap.c:/^devintr/} 中。

许多设备驱动程序在两个上下文中执行代码：
\indextext{上半部分（top half）}
在进程的内核线程中运行，
以及 \indextext{下半部分（bottom half）}
在中断时执行。上半部分
通过系统调用（如 {\tt read} 和 {\tt write}）被调用，
这些调用希望设备执行 I/O。
这段代码可能会要求硬件开始一个
操作（例如，要求磁盘读取一个块）；然后代码等待
操作完成。最终设备完成
操作并产生中断。驱动程序的中断处理程序，
作为下半部分，
确定什么操作已经完成，
如果需要则唤醒等待的进程，
并告诉硬件开始下一个操作（如果有的话）。

\section{代码：控制台输入}

控制台驱动程序 \fileref{kernel/console.c}
是驱动程序结构的简单示例。
控制台驱动程序通过连接到 RISC-V 的 \indextext{UART（通用异步收发传输器）}
串口硬件接收人类输入的字符。
控制台驱动程序一次累积一行输入，
处理特殊输入字符如
退格字符和控制键-u。用户进程，如 shell，使用
{\tt read} 系统调用从控制台获取输入行。
当你在 QEMU 中向 xv6 输入时，你的按键通过
QEMU 模拟的 UART 硬件传递给 xv6。

驱动程序与之通信的 UART 硬件是一个 16550
芯片~\cite{ns16550a}，由 QEMU 模拟。在真实的计算机上，16550
会管理连接到终端或其他计算机的 RS232 串行链路。
当运行 QEMU 时，它连接到你的键盘和显示器。

UART 硬件对软件来说表现为一组 \indextext{内存映射（memory-mapped）}
控制寄存器。也就是说，有一些物理地址
连接到 UART 设备，因此加载和存储
与设备硬件交互而不是与 RAM 交互。
UART 的内存映射地址从 0x10000000 开始，即 {\tt UART0}
\lineref{kernel/memlayout.h:/UART0.0x/}。
UART 有一些控制寄存器，每个都是
一个字节的宽度。它们相对于 {\tt UART0} 的偏移量在
\lineref{kernel/uart.c:/define.RHR/} 中定义。例如，
{\tt LSR} 寄存器包含指示输入
字符是否等待驱动程序读取的位。这些
字符（如果有）可以从
{\tt RHR} 寄存器读取。每次读取一个字符，UART 硬件
就从等待字符的内部 FIFO 中删除它，
并在 FIFO 为空时清除 {\tt LSR} 中的``就绪''位。
要传输数据，驱动程序向 {\tt THR} 寄存器
写入一个字节，这会导致 UART 将该字节附加到
UART 将在 RS232 串行链路上发送的字节 FIFO 中。
UART 发送和接收硬件在很大程度上是
相互独立的。

Xv6 的 {\tt main} 调用 {\tt consoleinit}
\lineref{kernel/console.c:/^consoleinit/} 来初始化 UART
硬件。这段代码配置 UART 在接收到
每个输入字节时产生
接收中断，
并在 UART 完成发送每个输出字节时产生
\indextext{发送完成（transmit complete）}
中断 \lineref{kernel/uart.c:/^uartinit/}。

xv6 shell 通过 {\tt init.c} 打开的文件描述符
\lineref{user/init.c:/open..console/} 从控制台读取。
对 {\tt read} 系统调用的调用通过内核到达 {\tt
  consoleread} \lineref{kernel/console.c:/^consoleread/}。{\tt
  consoleread} 等待输入到达（通过中断）并
缓冲在 {\tt cons.buf} 中，将输入复制到用户空间，
并（在一整行到达后）返回到用户进程。如果用户
还没有输入完整的一行，任何读取进程都会在
{\tt sleep} 调用中等待
\lineref{kernel/console.c:/sleep..cons/}
（第~\ref{CH:SLEEP} 章解释了 {\tt sleep} 的细节）。

当用户键入一个字符时，UART 硬件请求 RISC-V
产生中断，这会激活
xv6 的陷入处理程序。
陷入处理程序调用 {\tt devintr}
\lineref{kernel/trap.c:/^devintr/}，
它查看 RISC-V 的 {\tt scause} 寄存器以发现
中断来自外部设备。
然后它询问一个称为 PLIC 的硬件单元
\cite{riscv:priv}
告诉它是哪个设备产生的中断
\lineref{kernel/trap.c:/plic.claim/}。
如果是 UART，{\tt devintr} 调用 {\tt uartintr}。

{\tt uartintr}
\lineref{kernel/uart.c:/^uartintr/}
从 UART 硬件读取任何等待的输入字符
并将它们传递给 {\tt consoleintr}
\lineref{kernel/console.c:/^consoleintr/}；它不会
等待字符，因为将来的输入会产生新的中断。
{\tt consoleintr} 的工作是在
{\tt cons.buf} 中累积输入字符
直到一整行到达。
{\tt consoleintr} 特殊处理退格字符和一些其他字符。
当新行到达时，{\tt consoleintr} 唤醒
等待的 {\tt consoleread}（如果有的话）。

一旦被唤醒，{\tt consoleread} 将观察到 {\tt
  cons.buf} 中有一整行，将其复制到用户空间，
并（通过系统调用机制）返回到用户空间。

\section{代码：控制台输出}

在连接到控制台的文件描述符上执行 {\tt write} 系统调用
最终会到达
{\tt uartputc}
\lineref{kernel/uart.c:/^uartputc/}。
设备驱动程序维护一个输出缓冲区（{\tt uart\_tx\_buf}）
以便写入进程不必等待 UART 完成
发送；相反，{\tt uartputc} 将每个字符附加到缓冲区，
调用 {\tt uartstart} 启动设备传输（如果它还没有
开始），然后返回。{\tt uartputc}
唯一等待的情况是缓冲区已满。

每次 UART 完成发送一个字节时，它会产生一个中断。
{\tt uartintr} 调用 {\tt uartstart}，它检查设备
确实已经完成发送，并将下一个缓冲的输出字符交给设备。
因此，如果一个进程向控制台写入多个字节，
通常第一个字节将由 {\tt uartputc} 调用
{\tt uartstart} 发送，其余缓冲的字节将由
{\tt uartintr} 调用 {\tt uartstart} 在传输完成
中断到达时发送。

需要注意的一个通用模式是通过缓冲和中断将设备活动
与进程活动解耦。即使没有进程在等待读取，
控制台驱动程序也可以处理输入；
后续的读取会看到该输入。同样，进程可以发送输出
而不必等待设备。这种解耦可以通过允许进程
与设备 I/O 并发执行来提高性能，
并且在设备较慢（如 UART）或需要立即关注（如回显
键入的字符）时尤为重要。这个想法有时称为 \indextext{I/O
  并发（I/O concurrency）}。

\section{驱动程序中的并发}

你可能已经注意到在 {\tt consoleread}
和 {\tt consoleintr} 中调用了 {\tt acquire}。这些调用获取一个锁，
保护控制台驱动程序的数据结构免受并发访问。
这里有三个并发危险：两个在不同 CPU 上的进程可能
同时调用 {\tt consoleread}；
硬件可能请求 CPU 在 CPU 已经在 {\tt consoleread} 内部
执行时传递控制台（实际上是
UART）中断；
以及硬件可能在 {\tt consoleread} 执行时
在不同的 CPU 上传递控制台中断。
第~\ref{CH:LOCK} 章解释了如何使用锁
确保这些危险不会导致不正确的结果。

驱动程序中需要注意并发的另一种方式是，一个
进程可能正在等待设备的输入，但指示输入到达的中断
可能在不同的进程（或根本没有进程）运行时到达。
因此中断处理程序不允许考虑它们中断的进程或代码。
例如，中断处理程序不能安全地使用当前进程的页表调用 {\tt copyout}。中断处理程序通常
做相对较少的工作（例如，只是将输入数据复制到缓冲区），
并唤醒上半部分代码来完成其余工作。

\section{定时器中断}

Xv6 使用定时器中断来维持其对
当前时间的概念，并在计算密集型进程之间切换。定时器
中断来自连接到每个 RISC-V CPU 的时钟硬件。Xv6
为每个 CPU 的时钟硬件编程以定期中断 CPU。

{\tt start.c} 中的代码
\lineref{kernel/start.c:/^timerinit/} 设置一些控制位
允许监督模式访问定时器控制
寄存器，然后请求第一个定时器中断。
{\tt time} 控制寄存器包含一个
硬件以稳定速率递增的计数值；这作为
当前时间的概念。{\tt stimecmp} 寄存器
包含 CPU 将产生定时器
中断的时间；将 {\tt stimecmp} 设置为 {\tt time} 的当前值
加上 {\it x} 将安排一个 {\it x} 时间单位后的中断。
对于 {\tt qemu} 的 RISC-V
仿真，1000000 个时间单位约等于十分之一秒。

定时器中断通过 {\tt usertrap} 或 {\tt kerneltrap}
和 {\tt devintr} 到达，就像其他设备中断一样。
定时器中断到达时 {\tt scause} 的低位设置为
5；{\tt trap.c} 中的 {\tt devintr} 检测这种情况
并调用 {\tt clockintr}
\lineref{kernel/trap.c:/clockintr/}。
后者递增 {\tt ticks}，
允许内核跟踪
时间的流逝。递增只发生在一个 CPU 上，以避免多个 CPU 时时间过得更快。
{\tt clockintr} 唤醒在 {\tt pause}
系统调用中等待的任何进程，
并通过写入
{\tt stimecmp} 安排下一个定时器中断。

{\tt devintr} 为定时器中断返回 2
以向 {\tt kerneltrap}
或 {\tt usertrap} 指示它们应该调用 {\tt yield} 以便
CPU 可以在可运行进程之间多路复用。

内核代码可能被强制通过 {\tt yield} 进行上下文切换的定时器中断中断，
这是 {\tt usertrap} 中早期代码在启用中断之前小心保存 {\tt
  sepc} 等状态的部分原因。这些上下文切换也意味着
内核代码必须知道它可能在毫无预警的情况下从一个 CPU 移动到另一个 CPU。

\section{现实世界}

Xv6，像许多操作系统一样，允许在执行内核时产生中断甚至上下文切换（通过 {\tt yield}）。这样做的原因是在运行很长时间的复杂系统调用期间保持快速响应时间。然而，如上所述，在内核中允许中断是一些复杂性的来源；因此，少数操作系统只允许在执行用户代码时产生中断。

在典型计算机上支持所有设备的完整功能
是一项巨大的工作，因为有很多设备，设备有很多
功能，而且设备和驱动程序之间的协议可能很复杂
且文档不完善。在许多操作系统中，驱动程序占用的代码
比核心内核还多。

UART 驱动程序通过读取 UART 控制寄存器一次读取一个字节的数据；
这种模式称为 \indextext{程序控制 I/O（programmed I/O）}，
因为软件驱动数据移动。程序控制 I/O 很简单，但
速度太慢，无法在高数据速率下使用。需要以高速移动大量
数据的设备通常使用 \indextext{直接内存访问（DMA, direct memory access）}。
DMA 设备硬件直接将传入的数据写入 RAM，并从 RAM 读取
传出数据。现代磁盘和网络设备使用 DMA。DMA
设备的驱动程序会在 RAM 中准备数据，然后使用对控制寄存器的
单次写入告诉设备处理准备好的数据。

当设备需要在不可预测的时间关注，而且频率不太高时，
中断是有意义的。但中断具有很高的 CPU 开销。因此
高速设备，如网络和磁盘控制器，使用技巧来减少对中断的需求。
一个技巧是为整批传入或传出的请求产生单个
中断。另一个技巧是让驱动程序完全禁用中断，并定期检查设备是否需要关注。这种技术
称为 \indextext{轮询（polling）}。如果设备以高频率执行
操作，轮询是有意义的，但如果设备大部分时间是空闲的，
它会浪费 CPU 时间。一些驱动程序根据当前设备负载
在轮询和中断之间动态切换。

UART 驱动程序首先将传入数据复制到内核中的缓冲区，
然后复制到用户空间。这在低数据速率下是有意义的，但这种
双重复制可能会显著降低生成或消费数据非常快的设备的性能。
一些操作系统能够直接在用户空间缓冲区和设备硬件之间
移动数据，通常使用 DMA。

如第~\ref{CH:UNIX} 章所述，控制台对
应用程序来说表现为常规文件，应用程序使用 \lstinline{read} 和 \lstinline{write} 系统调用读取输入和写入输出。
应用程序可能想要控制无法通过标准文件系统调用表达的设备方面（例如，
在控制台驱动程序中启用/禁用行缓冲）。Unix
操作系统为此类情况提供 \lstinline{ioctl} 系统调用。

计算机的某些用途需要对外部事件的``实时''响应：
即保证在有限时间内发生的响应。
例如，在安全关键系统中，错过截止时间可能导致
灾难。Xv6 不适合实时场景。除了其他问题外，
xv6 的调度器在决定接下来运行哪个进程时没有考虑实时截止时间，
而且 xv6 有很长的禁用中断的内核代码路径，因此它可能无法快速响应中断。
实时操作系统不仅要解决这些问题，还必须以一种允许分析最坏情况响应时间的方式构建。

\section{练习}

\begin{enumerate}

\item 修改 {\tt uart.c} 以完全不使用中断。你可能需要
修改 {\tt console.c} 。

\item 为以太网卡编写驱动程序。

\end{enumerate}

\chapter{锁}
\label{CH:LOCK}

大多数内核，包括 xv6，都交错执行多个活动。交错的一个来源是多处理器硬件：具有多个独立执行的 CPU 的计算机，例如 xv6 的 RISC-V。这些多个 CPU 共享物理 RAM，xv6 利用这种共享来维护所有 CPU 读写内核数据结构。这种共享带来了一种可能性：一个 CPU 正在读取数据结构，而另一个 CPU 正在更新它的中途，甚至多个 CPU 同时更新同一数据；如果没有精心设计，这种并行访问可能会产生不正确的结果或破坏的数据结构。即使在单处理器上，内核也可能在多个线程之间切换 CPU，导致它们的执行被交错。最后，修改与某些可中断代码相同数据的设备中断处理程序，如果中断发生在恰好的错误时间，可能会损坏数据。
\indextext{并发}
指的是由于多处理器并行、线程切换或中断而导致多条指令流交错的情况。

内核充满了并发访问的数据。例如，两个 CPU 可以同时调用 {\tt kalloc}，从而并发地从空闲列表头部弹出。内核设计者喜欢允许大量并发，因为它可以通过并行性提高性能，并提高响应性。然而，因此内核设计者必须在尽管存在这种并发的情况下说服自己代码的正确性。有许多方法可以实现正确的代码，有些比其他更容易推理。旨在实现并发下正确性的策略和支持它们的抽象被称为
\indextext{并发控制}
技术。

Xv6 根据情况使用多种并发控制技术；还有更多可能的。本章专注于一种广泛使用的技术：
\indextext{锁}。
锁提供互斥，确保一次只有一个 CPU 可以持有锁。如果程序员将锁与每个共享数据项关联，并且代码在使用项目时始终持有关联的锁，那么一次只有一个 CPU 会使用该项目。在这种情况下，我们说锁保护数据项。虽然锁是一种易于理解的并发控制机制，但锁的缺点是它们可能限制性能，因为它们串行化并发操作。

本章的其余部分解释 xv6 为什么需要锁、xv6 如何实现它们以及如何使用它们。


% The following is most relevant to lock-free code:
% 
% A key observation will be that if you look at some code in
% xv6, you must ask yourself if concurrent code could change
% the intended behavior of the code by modifying data (or hardware resources)
% it depends on.
% You must keep in mind that the compiler may turn a
% single C statement into several machine instructions,
% and that those instructions may execute in a way that is
% interleaved with instructions executing on other CPUs.
% That is, you cannot assume that lines of C code
% on the page are executed atomically.
% Concurrency makes reasoning about correctness difficult.

%% 
\section{竞态}
%% 

\begin{figure}[t]
\center
\includegraphics[scale=0.8]{fig/smp.pdf}
\caption{简化的 SMP 架构}
\label{fig:smp}
\end{figure}

作为我们需要锁的一个例子，考虑两个具有已退出子进程的进程在两个不同的 CPU 上调用
{\tt wait} 系统调用。{\tt wait} 释放子进程的内存。因此在每个 CPU 上，内核将调用 {\tt kfree}
来释放子进程的内存页。内核分配器维护一个空闲页面的链表：\lstinline{kalloc()} \lineref{kernel/kalloc.c:/^kalloc/} 从列表中弹出一页内存，\lstinline{kfree()}
\lineref{kernel/kalloc.c:/^kfree/} 将一页压入列表。为了获得最佳性能，我们可能希望两个父进程的 {\tt kfree} 并行执行，而无需等待对方，但鉴于 xv6 的 {\tt kfree} 实现，这是不正确的。

Figure~\ref{fig:smp} 更详细地说明了设置：空闲页面链表位于两个 CPU 共享的内存中，两个 CPU 使用加载和存储指令操作列表。（实际上，处理器有缓存，但概念上多处理器系统表现得好像有单个共享内存。）
If there were no concurrent
requests, you might implement a list \lstinline{push} operation as
follows:
\begin{lstlisting}[numbers=left,firstnumber=1]
    struct element {
      int data;
      struct element *next;
    };
    
    struct element *list = 0;
    
    void
    push(int data)
    {
      struct element *l;
   
      l = malloc(sizeof *l);
      l->data = data;
      l->next = list; (*@\label{line:next}@*)
      list = l;  (*@\label{line:list}@*)
   }
\end{lstlisting}

\begin{figure}[t]
\center
\includegraphics[scale=0.5]{fig/race.pdf}
\caption{Example race}
\label{fig:race}
\end{figure}
This implementation is correct if executed in isolation.
However, the code is not correct if more than one
copy executes concurrently.
If two CPUs execute
\lstinline{push}
at the same time,
both might execute line~\ref{line:next} as shown in Fig~\ref{fig:smp},
before either executes line~\ref{line:list},  which results
in an incorrect outcome as illustrated by Figure~\ref{fig:race}.
There would then be two
list elements with
\lstinline{next}
set to the same former value of
\lstinline{list}.
When the two assignments to
\lstinline{list}
happen at line ~\ref{line:list},
the second one will overwrite the first;
the element involved in the first assignment
will be lost.

第 \ref{line:list} 行的丢失更新是
\indextext{竞态}
的一个例子。竞态是一种内存位置被并发访问的情况，且至少一次访问是写入。竞态通常是错误的迹象，要么是丢失更新（如果访问是写入），要么是读取不完全更新的数据结构。竞态的结果取决于编译器生成的机器代码、两个涉及的 CPU 的时序，以及它们的内存操作被内存系统排序的方式，这可能使竞态引起的错误难以重现和调试。例如，在调试
\lstinline{push}
时添加打印语句可能会改变执行的时序，足以使竞态消失。

避免竞态的常用方法是使用锁。锁确保
\indextext{互斥}，
以便一次只有一个 CPU 可以执行
\lstinline{push}
的敏感代码行；这使上述情况不可能发生。上述代码的正确加锁版本只添加了几行（以黄色高亮显示）：
\begin{lstlisting}[numbers=left,firstnumber=6]
   struct element *list = 0;
   (*@\hl{struct lock listlock;}@*)
    	
   void
   push(int data)
   {
     struct element *l;
     l = malloc(sizeof *l); (*@\label{line:malloc}@*)
     l->data = data;
   
     (*@\hl{acquire(\&listlock);} @*)
     l->next = list;     (*@\label{line:next1}@*)
     list = l;           (*@\label{line:list1}@*)
     (*@\hl{release(\&listlock)}; @*)
   }
\end{lstlisting}
\lstinline{acquire}
和
\lstinline{release}
之间的指令序列通常被称为
\indextext{临界区}。
锁被认为是在保护
\lstinline{list}。

当我们说锁保护数据时，我们真正的意思是锁保护适用于数据的一组不变量。不变量是跨操作维护的数据结构的属性。通常，操作的正确行为取决于操作开始时不变量为真。操作可能暂时违反不变量，但必须在完成之前重新建立它们。例如，在链表情况下，不变量是
\lstinline{list}
指向列表中的第一个元素，且每个元素的
\lstinline{next}
字段指向下一个元素。
\lstinline{push}
的实现暂时违反此不变量：在第 \ref{line:next1} 行，
\lstinline{l}
指向
列表的下一个元素，但
\lstinline{list}
尚未指向
\lstinline{l}
（在第 \ref{line:list1} 行重新建立）。我们上面检查的竞态发生是因为第二个 CPU 执行了依赖于列表不变量的代码，而它们被（暂时）违反了。正确使用锁确保一次只有一个 CPU 可以在临界区操作数据结构，因此没有 CPU 会在数据结构不变量不成立时执行数据结构操作。

你可以将锁视为
\indextext{串行化}
并发临界区，使它们一次运行一个，从而保留不变量（假设临界区在隔离情况下是正确的）。你也可以将由相同锁保护的临界区视为彼此原子的，以便每个只看到来自早期临界区的完整更改集，而永远不会看到部分完成的更新。

Though useful for correctness, locks inherently
limit performance.  For example, if two processes call {\tt kfree}
concurrently, the locks will serialize the two critical sections,
so that there is no
benefit from running them on different CPUs.  We say that multiple
processes \indextext{conflict} if they want the same lock at the same
time, or that the lock experiences \indextext{contention}.  A major
challenge in kernel design is avoidance of lock contention
in pursuit of parallelism.  Xv6 does
little of that, but sophisticated kernels organize
data structures and algorithms specifically to avoid lock contention.  In the list
example, a kernel may maintain a separate free list per CPU and only touch
another CPU's free list if the current CPU's list is empty and it must steal
memory from another CPU.  Other use cases may require more
complicated designs.

The placement of locks is also important for performance.
For example, it would be correct to move
\lstinline{acquire}
earlier in
\lstinline{push},
before line~\ref{line:malloc}.
But this would likely reduce performance because then the calls
to
\lstinline{malloc}
would be serialized.
The section ``Using locks'' below provides some guidelines for where to insert
\lstinline{acquire}
and
\lstinline{release}
invocations.
%% 
\section{代码：锁}
%% 

请阅读 {\tt kernel/spinlock.h} 和 {\tt kernel/spinlock.c}。

Xv6 有两种类型的锁：自旋锁和睡眠锁。我们从自旋锁开始。Xv6 将自旋锁表示为
\indexcode{struct spinlock}
\lineref{kernel/spinlock.h:/struct.spinlock/}。
结构中的重要字段是
\lstinline{locked}，
当锁可用时该字段为零，当锁被持有时为非零。从逻辑上讲，xv6 应该通过执行如下代码获取锁：
\begin{lstlisting}[numbers=left,firstnumber=21]
   void
   acquire(struct spinlock *lk) // does not work!
   {
     for(;;) {
       if(lk->locked == 0) {  (*@\label{line:testlocked}@*)
         lk->locked = 1;      (*@\label{line:assign}@*)
         break;
       }
     }
   }
\end{lstlisting}
不幸的是，这种实现不能
保证多处理器上的互斥。
可能发生两个 CPU 同时
到达第 \ref{line:testlocked} 行，看到
\lstinline{lk->locked}
为零，然后都通过执行第 \ref{line:assign} 行获取锁。
此时，两个不同的 CPU 持有锁，
这违反了互斥属性。
我们需要一种方法来
使第 \ref{line:testlocked} 和 \ref{line:assign} 行作为一个
\indextext{原子}
（即不可分割的）步骤执行。

因为锁被广泛使用，
多核处理器通常提供一条指令，
可用于使
第 \ref{line:testlocked} 和 \ref{line:assign} 行
原子化。
在 RISC-V 上，这条指令是
\lstinline{amoswap r, a}。
\lstinline{amoswap}
读取内存地址 {\tt a} 处的值，
将寄存器 {\tt r} 的内容写入该地址，
并将其读取的值放入 {\tt r}。
也就是说，它交换寄存器和被寻址
内存位置的内容。
它使用特殊
硬件原子地执行此序列，以防止任何
其他 CPU 在读取和写入之间使用该内存地址。

可移植的 C 库调用
\lstinline{__sync_lock_test_and_set(addr, value)}
归结为
\lstinline{amoswap}
指令；
该函数返回 {\tt *addr} 的旧（交换的）内容。
这是在 {\tt acquire} 中编写循环的好方法：

\begin{lstlisting}[numbers=left,firstnumber=21]
  while(__sync_lock_test_and_set(&lk->locked, 1) != 0)
    ;
\end{lstlisting}

Xv6 的
\indexcode{acquire}
\lineref{kernel/spinlock.c:/^acquire/}
使用上述循环。
每次迭代将 1 交换到
\lstinline{lk->locked}
并检查先前的值；
如果先前的值为零，则我们已获取
锁，交换将设置
\lstinline{lk->locked}
为 1。
如果先前的值为 1，则其他某个 CPU
持有锁，我们将 1 原子地交换到
\lstinline{lk->locked}
的事实不会改变其值。

一旦获取锁，
\lstinline{acquire}
记录获取锁的 CPU
以供调试。
\lstinline{lk->cpu}
字段受锁保护
且只能在持有锁时更改。

函数
\indexcode{release}
\lineref{kernel/spinlock.c:/^release/}
与
\lstinline{acquire}
相反：
它清除
\lstinline{lk->cpu}
字段
然后释放锁。
从概念上讲，释放只需要将零赋给
\lstinline{lk->locked}。
C 标准允许编译器用多个存储指令实现赋值，
因此 C 赋值可能相对于
并发代码是非原子的。
Instead,
\lstinline{release}
uses the C library function
\lstinline{__sync_lock_release}
that performs an atomic assignment.
该函数也归结为 RISC-V
\lstinline{amoswap}
指令。
%% 
\section{代码：使用锁}
%% 
Xv6 在许多地方使用锁来避免竞态。
如上所述，
\lstinline{kalloc}
\lineref{kernel/kalloc.c:/^kalloc/}
和
\lstinline{kfree}
\lineref{kernel/kalloc.c:/^kfree/}
是一个很好的例子。
尝试练习 1 和 2 看看如果这些
函数省略锁会发生什么。
你可能会发现很难触发不正确的
行为，这表明很难可靠地测试代码
是否没有锁错误和竞态。
Xv6 可能还有尚未发现的竞态。

使用锁的一个难点是决定使用多少锁以及每个锁应该保护哪些数据和不变量。
有一些基本原则。
首先，任何时候一个 CPU 写入变量
而另一个 CPU 同时读取或写入它，
应该使用锁来防止两个
操作重叠。
其次，记住锁保护不变量：
如果不变量涉及多个内存位置，
通常它们都需要被单个锁保护
以确保不变量被维护。

上面的规则说明了何时需要锁，但没有说明何时不需要锁，为了效率不过度加锁很重要，
因为锁减少并行性。
如果并行性不重要，那么可以
安排只有一个线程而不用担心锁。
一个简单的内核可以通过拥有一个锁来在多处理器上做到这一点；
内核每次从用户空间进入内核时获取锁，
用于系统调用或中断；
内核在返回用户空间时释放锁。
许多单处理器操作系统已被转换以
使用这种方法在多处理器上运行，有时称为``大
内核锁''，但这种方法牺牲了并行性：一次只有一个
CPU 可以在内核中执行。
如果内核消耗
大量 CPU 时间，可以通过用不同的锁保护不同的对象或模块来获得更多并行性，
以便不同的 CPU 可以同时
在内核的不同部分执行。

作为粗粒度锁的一个例子，xv6 的 \lstinline{kalloc.c}
分配器有一个由单个锁保护的单个空闲列表。如果
不同 CPU 上的多个进程同时尝试分配页面，
每个都必须在 {\tt acquire} 中自旋等待轮次。
自旋浪费 CPU 时间，因为这不是有用的工作。
如果锁争用浪费了很大比例的 CPU 时间，
也许可以通过修改分配器
为每个 CPU 提供单独的空闲列表，每个列表有自己的锁，以允许真正
并行的分配来提高性能。

作为细粒度锁的一个例子，xv6 为每个文件有一个单独的锁，以便操作不同文件的进程通常可以
继续而不必等待对方的锁。如果希望允许
进程同时写入同一文件的不同区域，文件锁定
方案可以做得更细粒度。
最终锁粒度决策需要由性能
测量以及复杂性考虑来驱动。

随着后续章节解释 xv6 的每个部分，它们
会提到 xv6 使用锁
处理并发的例子。
作为预览，
Figure~\ref{fig:locktable}
列出了 xv6 中的所有锁。

\begin{figure}[t]
\center
\begin{tabular}{ll}
{\bf Lock} & {\bf Description} \\
\midrule
bcache.lock & Protects allocation of block buffer cache entries \\
cons.lock & Serializes read processing of console input \\
tx\_lock & Serializes access to console (uart) output hardware \\
ftable.lock & Serializes allocation of a struct file in file table \\
itable.lock & Protects allocation of in-memory inode entries \\
vdisk\_lock & Serializes access to disk hardware and queue of DMA descriptors \\
kmem.lock & Serializes allocation of memory \\
log.lock & Serializes operations on the transaction log \\
pipe's pi->lock & Serializes operations on each pipe \\
pid\_lock & Serializes increments of next\_pid \\
proc's p->lock & Serializes changes to process's state \\
wait\_lock & Helps wait avoid lost wakeups \\
tickslock & Serializes operations on the ticks counter \\
inode's ip->lock & Serializes operations on each inode and its content \\
buf's b->lock & Serializes operations on each block buffer \\
\end{tabular}
\caption{xv6 中的锁}
\label{fig:locktable}
\end{figure}

%% 
\newpage
\section{死锁和锁排序}
%% 

假设在 CPU C1 和 C2 上运行的函数都有一个点，每个都需要持有锁 A 和锁 B，并且它们以不同的顺序获取它们：

\begin{center}
  \input{fig/order.tex}
\end{center}

运气不好的话，C1 和 C2 可能都在完全相同的时刻执行它们的第一次获取；两者都可能成功，因为它们要求的是不同的锁。但然后 C1 和 C2 都必须在它们对 {\tt acquire()} 的第二次调用中等待，因为两个锁都已被另一个 CPU 持有。因为两个 CPU 都在等待对方，所以谁都不会释放锁，两者都将永远等待。这种情况被称为
\indextext{死锁}。

C1/C2 示例中的关键问题是两个 CPU 以不同的顺序获取锁。如果它们都先尝试获取 A，一个会先获取 A 然后 B 然后释放它们，然后另一个 CPU 可以继续。更一般地说，锁代码必须遵循以下规则以避免死锁：所有持有多个锁的代码路径必须以相同的顺序获取锁。这种全局锁获取顺序的需要意味着锁实际上是每个函数规范的一部分：调用者必须以导致锁按约定顺序获取的方式调用函数。

Xv6 有许多长度为二的锁顺序链，涉及每个进程的锁（每个
\lstinline{struct proc}
中的锁）由于
\lstinline{sleep}
的工作方式（参见 Chapter~\ref{CH:SLEEP}）。例如，
\lstinline{consoleintr}
\lineref{kernel/console.c:/^consoleintr/}
是处理输入字符的中断例程。当新行到达时，任何等待控制台输入的进程都应该被唤醒。为此，
\lstinline{consoleintr}
持有
\lstinline{cons.lock}
同时调用
\indexcode{wakeup}，后者获取等待进程的锁以唤醒它。因此，全局避免死锁的锁顺序包括规则
\lstinline{cons.lock}
必须在任何进程锁之前获取。文件系统代码包含 xv6 最长的锁链。例如，创建文件需要同时持有目录上的锁、新文件 inode 上的锁、磁盘块缓冲区上的锁、磁盘驱动程序的 \lstinline{vdisk_lock} 以及调用进程的 \lstinline{p->lock}。为了避免死锁，文件系统代码始终按上一句提到的顺序获取锁。

遵守全局避免死锁的顺序可能出人意料地困难。有时锁顺序与逻辑程序结构冲突，例如，也许代码模块 M1 调用模块 M2，但锁顺序要求先获取 M2 中的锁再获取 M1 中的锁。有时锁的标识事先不知道，也许是因为必须先持有一个锁才能发现要获取的下一个锁的标识。这种情况在文件系统中出现，因为它查找路径名中的连续组件，以及在 {\tt wait} 和 {\tt exit} 系统调用的代码中，因为它们搜索进程表寻找子进程。最后，死锁的危险通常是锁定方案可以做得多细粒度的约束，因为更多的锁通常意味着更多的死锁机会。避免死锁的需要通常是内核实现中的一个主要因素。

有时会出现一个问题是，如果 CPU 尝试获取同一 CPU 已持有的锁应该发生什么。一条推理线是这应该被允许：没有其他 CPU 可以持有该锁，所以不必担心 CPU 干扰彼此使用受保护的数据。允许 CPU 重新获取它已持有的锁的锁系统称为{\it 可重入}或{\it 递归}。另一方面，如果锁已被持有，即使在同一 CPU 上，也意味着操作可能已暂时违反且尚未恢复某些不变量；允许在不变量不成立时开始新操作似乎是错误的邀请。Xv6 持后一种观点，并禁止当前持有锁的 CPU 重新获取它。检测这种情况是 {\tt acquire} 中调用 {\tt holding} 的目的。

\section{锁和中断}
\label{s:lockinter}

%% 
Some xv6 spinlocks protect data that is used by
both threads and interrupt handlers.
For example, the
\lstinline{clockintr}
timer interrupt handler might increment
\indexcode{ticks} 
\lineref{kernel/trap.c:/^clockintr/}
at about the same time that a kernel
thread reads
\lstinline{ticks} 
in
\indexcode{sys_pause}
\lineref{kernel/sysproc.c:/ticks0.=.ticks/}.
The lock
\indexcode{tickslock}
serializes the two accesses.

The interaction of spinlocks and interrupts raises a potential danger.
Suppose
\indexcode{sys_pause}
holds
\indexcode{tickslock},
and its CPU is interrupted by a timer interrupt.
\lstinline{clockintr}
would try to acquire
\lstinline{tickslock},
see it was held, and wait for it to be released.
In this situation,
\lstinline{tickslock}
will never be released: only
\lstinline{sys_pause}
can release it, but
\lstinline{sys_pause}
will not continue running until
\lstinline{clockintr}
returns.
So the CPU will deadlock, and any code
that needs either lock will also freeze.

To avoid this situation, if a spinlock is used by an interrupt handler,
a CPU must never hold that lock with interrupts enabled.
Xv6 is more conservative: when a CPU acquires any
lock, xv6 always disables interrupts on that CPU.
Interrupts may still occur on other CPUs, so 
an interrupt's
\lstinline{acquire}
can wait for a thread to release a spinlock; just not on the same CPU.

Xv6 re-enables interrupts when a CPU holds no spinlocks; it must
do a little book-keeping to cope with nested critical sections.
\lstinline{acquire}
calls
\indexcode{push_off}
\lineref{kernel/spinlock.c:/^push_off/}
and
\lstinline{release}
calls
\indexcode{pop_off}
\lineref{kernel/spinlock.c:/^pop_off/}
来跟踪当前 CPU 上锁的嵌套级别。
当该计数达到零时，
\lstinline{pop_off}
恢复在最外层临界区开始时存在的中断使能状态。
\lstinline{intr_off}
和
\lstinline{intr_on}
函数分别执行 RISC-V 指令以禁用和启用中断。

重要的是
\indexcode{acquire}
严格在设置
\lstinline{lk->locked}
\lineref{kernel/spinlock.c:/sync_lock_test_and_set/}
之前调用
\lstinline{push_off}。
如果两者顺序颠倒，会有一个短暂的窗口期，锁被持有但中断是启用的，一个不幸时序的中断会使系统死锁。
类似地，重要的是
\indexcode{release}
只在
释放锁之后
\lineref{kernel/spinlock.c:/sync_lock_release/}
调用
\indexcode{pop_off}。

%% 
\section{指令和内存排序}
%% 

将程序视为按源代码语句出现的顺序执行是很自然的。对于单线程代码，这是一个合理的心智模型，但当多个线程通过共享内存交互时，这是不正确的。一个原因是编译器以不同于源代码暗示的顺序发出加载和存储指令，并且可能完全省略它们（例如通过在寄存器中缓存数据）。另一个原因是 CPU 可能为了性能而乱序执行指令。例如，CPU 可能注意到在串行指令序列中 A 和 B 彼此不依赖。CPU 可能先开始指令 B，要么因为它的输入比 A 的输入先准备好，要么为了重叠执行 A 和 B。

As an example of what could go wrong,
in this code for
\lstinline{push},
it would be a disaster if the compiler or CPU moved the
store corresponding to
line~\ref{line:next2} to a point after the
\lstinline{release}
on line~\ref{line:release}:
\begin{lstlisting}[numbers=left,firstnumber=1]
      l = malloc(sizeof *l);
      l->data = data;
      acquire(&listlock);
      l->next = list;   (*@\label{line:next2}@*)
      list = l;      
      release(&listlock);  (*@\label{line:release}@*)
\end{lstlisting}
If such a re-ordering occurred, there would be a window during
which another CPU could acquire the lock and
observe the updated
\lstinline{list},
but see an uninitialized
\lstinline{list->next}.

The good news is that compilers and CPUs help concurrent programmers
by following a set of rules called the \indextext{memory model}, and
by providing some primitives to help programmers control re-ordering.

To tell the hardware and compiler not to re-order,
xv6 uses
\lstinline{__sync_synchronize()} 
in both
\lstinline{acquire} \lineref{kernel/spinlock.c:/^acquire/}
and
\lstinline{release} \lineref{kernel/spinlock.c:/^release/}.
\lstinline{__sync_synchronize()}
is a \indextext{memory barrier}:
it tells the compiler and CPU to not reorder loads or stores across the
barrier.
The barriers in xv6's 
\lstinline{acquire}
and
\lstinline{release}
force order in almost all cases where it matters,
since xv6 uses locks around accesses to shared data.
Chapter~\ref{CH:LOCK2} discusses a few exceptions.

%% 
\section{Sleep locks}
%% 

% this material refers a lot to sleep, yield, etc, so maybe it should
% be later, after sleep/wakeup in sched.tex.
% maybe in lock2.tex

Sometimes xv6 needs to hold a lock for a long time. For example, the
file system (Chapter~\ref{CH:FS}) keeps a file locked while reading
and writing its content on the disk, and these disk operations can
take tens of milliseconds. Holding a spinlock that long would lead to
waste if another process wanted to acquire it, since the acquiring
process would waste CPU for a long time while spinning. Another
drawback of spinlocks is that a process cannot yield the CPU while
retaining a spinlock; we'd like to do this so that other processes can
use the CPU while the process with the lock waits for the disk.
Yielding while holding a spinlock is illegal because it might
lead to deadlock if a second thread then tried to acquire the spinlock;
since {\tt acquire} doesn't yield the CPU, the second thread's
spinning might
prevent the first thread from running and releasing the lock.
% it's true that timer interrupts might force a switch from the
% second back to the first thread -- but not if the second
% already holds one or more locks.
% it's also true that the danger is most obvious on a uniprocessor,
% but it can also arise with a longer chain of threads on
% a multiprocessor.
Yielding while holding a lock would also violate
the requirement that interrupts must be off while a spinlock is held.
Thus we'd like a type of lock that yields the CPU while waiting to
acquire, and allows yields (and interrupts) while the lock is held.

Xv6 provides such locks in the form of
\indextext{sleep-locks}.
\lstinline{acquiresleep}
\lineref{kernel/sleeplock.c:/^acquiresleep/}
yields the CPU while waiting,
using techniques that will be explained in
Chapter~\ref{CH:SLEEP}.
At a high level, a sleep-lock has a
\lstinline{locked}
field that is protected by a spinlock, and 
\lstinline{acquiresleep} 's
call to
\lstinline{sleep}
atomically yields the CPU and releases the spinlock.
The result is that other threads can execute while
\lstinline{acquiresleep}
waits.

Because sleep-locks leave interrupts enabled, they cannot be
used in interrupt handlers.
Because
\lstinline{acquiresleep}
may yield the CPU,
sleep-locks cannot be used inside spinlock critical
sections (though spinlocks can be used inside sleep-lock
critical sections).

Spin-locks are best suited to short critical sections,
since waiting for them wastes CPU time;
sleep-locks work well for lengthy operations.

%% 
\section{Real world}
%% 
Programming with locks remains challenging despite years of research
into concurrency primitives and parallelism.
It is often best to conceal locks within 
higher-level constructs like synchronized queues, although xv6 does not
do this.  If you program with locks, it is wise to use a tool that attempts to
identify races, because it is easy to miss an invariant that requires
a lock.

Most operating systems support POSIX threads (Pthreads), which allow a
user process to have several threads running concurrently on different
CPUs.  Pthreads has support for user-level locks, barriers, etc.
Pthreads also allows a programmer to optionally specify that a lock
should be re-entrant.

Supporting Pthreads at user level requires support from the operating
system. For example, it should be the case that if one pthread blocks
in a system call, another pthread of the same process should be able
to run on that CPU.  As another example, if a pthread changes its
process's address space (e.g., maps or unmaps memory), the kernel must
arrange that other CPUs that run threads of the same process update
their hardware page tables to reflect the change in the address space.

It is possible to implement locks without atomic
instructions~\cite{lamport:bakery}, but it is expensive, so most
operating systems use atomic instructions.

Locks can be expensive if many CPUs try to acquire the same lock
at the same time.  If one CPU has a lock
cached in its local cache, and another CPU must acquire the lock, then the
atomic instruction to update the cache line that holds the lock must move the line
from the one CPU's cache to the other CPU's cache, and perhaps
invalidate any other copies of the cache line.  Fetching a cache line from
another CPU's cache can be orders of magnitude more expensive than
fetching a line from a local cache.

To avoid the expenses associated with locks, many operating systems
use lock-free data structures and algorithms~\cite{herlihy:art,mckenney:rcuusage}.
For example, it is possible to implement a linked list like the one in
the beginning of the chapter that requires no locks during list
searches, and one atomic instruction to insert an item in a list.
Lock-free programming is more complicated, however, than programming
locks; for example, one must worry about instruction and memory
reordering.  Programming with locks is already hard, so xv6 avoids the
additional complexity of lock-free programming.

%% 
\section{Exercises}
%% 

\begin{enumerate}
  
\item Comment out the calls to
\lstinline{acquire}
and
\lstinline{release}
in
\lstinline{kalloc}
\lineref{kernel/kalloc.c:/^kalloc/}.
This seems like it should cause problems for
kernel code that calls
\lstinline{kalloc};
what symptoms do you expect to see?
When you run xv6, do you see these symptoms?
How about when running
\lstinline{usertests}?
If you don't see a problem, why not?
See if you can provoke a problem by inserting
dummy loops into the critical section of
\lstinline{kalloc}.

\item Suppose that you instead commented out the
locking in
\lstinline{kfree} 
(after restoring locking in
\lstinline{kalloc}).
What might now go wrong? Is lack of locks in
\lstinline{kfree}
less harmful than in
\lstinline{kalloc}?

\item If two CPUs call
\lstinline{kalloc}
at the same time, one will have to wait for the other,
which is bad for performance.
Modify 
\lstinline{kalloc.c}
to have more parallelism, so that simultaneous
calls to
\lstinline{kalloc}
from different CPUs can proceed without waiting for each other.

\item Write a parallel program using POSIX threads, which is supported on most
operating systems. For example, implement a parallel hash table and measure if
the number of puts/gets scales with increasing number of CPUs.

\item Implement a subset of Pthreads in xv6.  That is, implement a user-level
thread library so that a user process can have more than 1 thread and arrange
that these threads can run in parallel on different CPUs.  Come up with a
design that correctly handles a thread making a blocking system call and
changing its shared address space.

\end{enumerate}

% LocalWords:  CPUs uniprocessor interruptible allocator kalloc kfree
% LocalWords:  struct malloc sizeof listlock invariants spinlock lk
% LocalWords:  amoswap cpu bcache ftable itable inode vdisk DMA kmem
% LocalWords:  pid proc's tickslock inode's ip buf's proc SCHED sys
% LocalWords:  consoleintr wakeup clockintr intr acquiresleep POSIX
% LocalWords:  Pthreads pthread unmaps TLB IPIs SMP firstnumber int
% LocalWords:  structure's spinlocks addr tx uart pipe's wakeups A's
% LocalWords:  usertests

%
% process has one thread with two stacks
%
\chapter{调度}
\label{CH:SCHED}

任何操作系统运行的进程数都可能超过计算机的 CPU 数量，因此需要一种机制在进程之间
\indextext{多路复用（multiplexing）}
CPU。理想情况下，这种共享应对用户进程透明。一种常见的方法是为每个进程营造拥有独立虚拟 CPU 的假象，方法是将进程
多路复用到硬件 CPU 上。本章阐述 xv6 如何实现这种多路复用。

在继续阅读本章之前，请先阅读 {\tt kernel/proc.h}、
{\tt kernel/swtch.S}，
以及 {\tt kernel/proc.c} 中的
{\tt yield()}、
{\tt sched()}
和
{\tt schedule()}。

%%
\section{多路复用}
%%

Xv6 在两种情况下切换 CPU 以实现多路复用：
首先，当进程发起阻塞式系统调用（需要等待）时，例如 \lstinline{read} 或 \lstinline{wait}，xv6 会进行切换。
其次，xv6 定期强制切换，以处理长时间计算而不阻塞的进程。
前者称为自愿切换，
后者称为非自愿切换。

实现多路复用面临若干挑战。首先，如何从一个进程切换到另一个？基本思路是保存和恢复 CPU 寄存器，但这无法用 C 语言表达，因此变得棘手。其次，如何以对用户进程透明的方式强制切换？Xv6 采用标准技术，使用硬件定时器触发中断来驱动上下文切换。第三，所有 CPU 都在同一组进程之间切换，因此需要锁机制来避免错误，例如两个 CPU 同时决定运行同一个进程。第四，进程退出时必须释放其内存和其他资源，但进程自身无法完成所有这些工作。第五，多核机器的每个 CPU 必须记住它正在执行哪个进程，以便系统调用影响正确的进程内核状态。

%%
\section{上下文切换概述}
%%

术语``上下文切换（context switch）''指 CPU 停止执行一个内核线程（通常稍后恢复）并恢复执行另一个内核线程所涉及的步骤；这种切换是多路复用的核心。Xv6 不直接从一个进程的内核线程切换到另一个进程的内核线程；相反，内核线程通过切换到该 CPU 的``调度器线程（scheduler thread）''来放弃 CPU，然后调度器线程选择一个不同的进程内核线程来运行，并切换到该线程。

\begin{figure}[t]
\centering
\includegraphics[scale=0.5]{fig/switch.pdf}
\caption{从一个用户进程切换到另一个用户进程。在此示例中，xv6 运行在一个 CPU 上（因此有一个调度器线程）。}
\label{fig:switch}
\end{figure}

在更宏观的视角下，从一个用户进程切换到另一个用户进程所涉及的步骤如图~\ref{fig:switch} 所示：从旧进程的用户空间陷入（系统调用或中断）到其内核线程，切换到当前 CPU 的调度器线程，切换到新进程的内核线程，然后陷入返回到用户级进程。


%%
\section{代码：上下文切换}
%%

{\tt kernel/swtch.S} 中的函数 {\tt swtch()} 包含线程上下文切换的核心：它保存被切出线程的 CPU 寄存器，并恢复被切入线程先前保存的寄存器。这样足够的基本原因是线程状态由内存中的数据（如其栈）加上其 CPU 寄存器组成；线程内存无需保存和恢复，因为不同线程将数据保存在 RAM 的不同区域；但 CPU 只有一组寄存器，因此必须在线程之间切换（保存和恢复）它们。

% why must it be in assembler? why not C?
% why swtch(), not switch()?

每个线程的 {\tt struct proc} 包含一个 {\tt struct context}，用于在线程不运行时保存已保存的寄存器。CPU 调度器线程的 {\tt struct context} 位于该 CPU 的 {\tt struct cpu} 中。当线程 X 希望切换到线程 Y 时，线程 X 调用 {\tt swtch(\&X's context, \&Y's context)}。{\tt swtch()} 将当前 CPU 寄存器保存在 X 的上下文中，然后将 Y 的上下文内容加载到 CPU 寄存器中，然后返回。

以下是 {\tt swtch} 的简化副本：

\begin{verbatim}
swtch:
        sd ra, 0(a0)
        sd sp, 8(a0)
        sd s0, 16(a0)
        ...
        sd s11, 104(a0)

        ld ra, 0(a1)
        ld sp, 8(a1)
        ld s0, 16(a1)
        ...
        ld s11, 104(a1)

        ret
\end{verbatim}

{\tt a0} 保存第一个函数参数，{\tt a1} 保存第二个参数；在本例中，是两个 {\tt struct context} 指针。{\tt 16(a0)} 指从 {\tt a0} 指向的内存偏移 16 字节；参考 {\tt kernel/proc.h} 中 {\tt struct context} 的定义 \lineref{kernel/proc.h:/^struct.context/}，这是称为 {\tt s0} 的结构体字段。

{\tt swtch} 的 {\tt ret} 返回到哪里？它返回到 {\tt ra} 寄存器指向的指令。在线程 X 调用 {\tt swtch()} 切换到 Y 的示例中，当 {\tt ret} 执行时，{\tt ra} 刚从 Y 的 {\tt struct context} 加载。而 Y 的 {\tt struct context} 中的 {\tt ra} 最初是由 Y 调用 {\tt swtch} 时保存的，当时 Y 放弃了 CPU。因此 {\tt ret} 返回到 Y 调用 {\tt swtch()} 之后的指令；也就是说，X 调用 {\tt swtch()} 返回，就好像从 Y 最初调用 {\tt swtch()} 返回一样。而且 {\tt sp} 将是 Y 的栈指针，因为 {\tt swtch} 从 Y 的 {\tt struct context} 加载了 {\tt sp}；因此在返回时，Y 将在自己的栈上执行。{\tt swtch()} 无需直接保存或恢复程序计数器；保存和恢复 {\tt ra} 就足够了。

\lstinline{swtch}
\lineref{kernel/swtch.S:/swtch/}
保存被调用者保存的寄存器（{\tt ra,sp,s0..s11}），但不保存调用者保存的寄存器。
RISC-V 调用约定要求，如果代码需要在函数调用之间保留调用者保存寄存器中的值，
编译器必须生成指令，在函数调用之前将寄存器保存到栈中，
并在函数返回时从栈中恢复。
因此 \lstinline{swtch} 可以依赖调用它的函数已经保存了调用者保存的寄存器（要么是这样，
要么是调用函数不需要寄存器中的值）。

%%
\section{代码：调度}
%%

% why a separate thread

上一节介绍了
\indexcode{swtch}
的内部实现；
现在让我们将
\lstinline{swtch}
视为给定，并研究从一个进程的内核线程
通过调度器切换到另一个进程。
调度器以每个 CPU 一个特殊线程的形式存在，每个线程运行
\lstinline{scheduler}
函数。
该函数负责选择接下来运行哪个进程。
每个 CPU 都有自己的调度器线程，因为可能有多个 CPU 在任何给定时间寻找要运行的东西。
进程切换总是通过调度器线程进行，而不是直接从一个进程切换到另一个进程，
以避免在某些情况下调度器没有栈可以执行（例如，如果旧进程已退出，或当前没有其他进程想要运行）。

想要放弃 CPU 的进程必须
获取自己的进程锁
\indexcode{p->lock}，
释放它持有的任何其他锁，
更新自己的状态
(\lstinline{p->state})，
然后调用
\indexcode{sched}。
你可以在
\lstinline{yield}
\lineref{kernel/proc.c:/^yield/}、
\texttt{sleep}
和
\texttt{kexit}
中看到这一系列操作。
\indexcode{sched}
调用
\indexcode{swtch}
将当前上下文保存在
\lstinline{p->context}
中，并切换到
\indexcode{cpu->context}
中的调度器上下文。
\lstinline{swtch}
在调度器的栈上返回，
就好像
\indexcode{scheduler}
的
\lstinline{swtch}
已经返回一样
\lineref{kernel/proc.c:/swtch.*contex.*contex/}。

\lstinline{scheduler}
\lineref{kernel/proc.c:/^scheduler/}
运行一个循环：
找到一个要运行的进程，{\tt swtch()} 到它，最终它会 {\tt swtch()} 回调度器，然后调度器继续其循环。
调度器遍历进程表，
寻找可运行的进程，即具有
\lstinline{p->state}
\lstinline{==}
\lstinline{RUNNABLE}
的进程。
一旦找到进程，它就设置每个 CPU 的当前进程
变量
\lstinline{c->proc}，
将进程标记为
\lstinline{RUNNING}，
然后调用
\indexcode{swtch}
开始运行它
\linerefs{kernel/proc.c:/Switch.to/,/swtch/}。
在过去的某个时刻，目标进程必须调用了 {\tt swtch()}；
调度器对 {\tt swtch()} 的调用有效地从该早期调用返回。
图~\ref{fig:inout} 说明了这种模式。

\begin{figure}[t]
\input{fig/switch.tex}
\caption{{\tt swtch()} 总是以调度器线程作为源或目标，并且相关的 {\tt p->lock} 总是保持持有。}
\label{fig:inout}
\end{figure}

xv6 在调用
\lstinline{swtch}
时保持
\indexcode{p->lock}
：
\lstinline{swtch}
的调用者获取锁，
但在 {\tt swtch} 返回后在目标中释放它。
这种安排是不寻常的：更常见的是获取锁的线程也释放它。
Xv6 的上下文切换打破了这一约定，因为
\indexcode{p->state}
和
\lstinline{p->context}
必须一起原子地更新。
例如，如果在调用
\indexcode{swtch}
之前释放
\lstinline{p->lock}，
不同的 CPU $c$ 可能会决定运行该进程，因为
其状态是 \lstinline{RUNNABLE}。
CPU $c$ 将调用 \lstinline{swtch}，
它将从 \lstinline{p->context} 恢复，
而原始 CPU 仍在向
\lstinline{p->context}
保存。
结果将是该进程将在 CPU $c$ 上用部分保存的寄存器恢复，
而且两个 CPU 将使用相同的栈，这会导致混乱。
一旦 \lstinline{yield} 开始修改运行进程的状态
使其成为
\lstinline{RUNNABLE}，
\lstinline{p->lock} 必须保持持有，直到
进程已保存所有寄存器且
调度器在其栈上运行。
最早正确的释放点是在
\lstinline{scheduler}
（在其自己的栈上运行）
清除
\lstinline{c->proc}
之后。
同样，一旦
\lstinline{scheduler}
开始将 \lstinline{RUNNABLE} 进程转换为
\lstinline{RUNNING}，
在该进程的内核线程完全运行之前（在
\lstinline{swtch}
之后，例如在
\lstinline{yield}
中）不能释放锁。

有一种情况，调度器对
\indexcode{swtch}
的调用不会最终进入
\indexcode{sched}。
\lstinline{allocproc} 将新进程的上下文 {\tt ra}
寄存器设置为
\indexcode{forkret}
\lineref{kernel/proc.c:/^forkret/}，
因此其第一次 \lstinline{swtch} ``返回''到该函数的开头。
\lstinline{forkret}
存在是为了释放
\indexcode{p->lock}
并设置一些控制寄存器和陷入帧字段，
这些是返回用户空间所需的。
最后，{\tt forkret} 模拟从系统调用返回到用户空间的正常返回路径。


%%
\section{代码：mycpu 和 myproc}
%%

 xv6 经常需要获取指向当前进程的 \lstinline{proc} 结构体的指针。
在单核处理器上，可以使用全局变量指向当前 \lstinline{proc}。
这在多核机器上行不通，因为每个 CPU 执行不同的进程。
解决这个问题的方法是利用每个 CPU 都有自己的寄存器组这一事实。

当给定 CPU 在内核中执行时，xv6 确保 CPU 的 \lstinline{tp}
寄存器始终保存 CPU 的 hartid。
RISC-V 对其 CPU 进行编号，给每个一个唯一的 \indextext{hartid}。
\indexcode{mycpu}
\lineref{kernel/proc.c:/^mycpu/}
使用 \lstinline{tp} 索引一个
\lstinline{cpu} 结构体数组，并返回当前 CPU 的那个结构体。
\indexcode{struct cpu}
\lineref{kernel/proc.h:/^struct.cpu/}
保存一个指向正在该 CPU 上运行的
进程的 \lstinline{proc} 结构体的指针（如果有），
CPU 调度器线程的已保存寄存器，
以及管理中断禁用所需的嵌套自旋锁计数。

确保 CPU 的 \lstinline{tp} 保存 CPU 的 hartid 有点复杂，因为用户代码可以自由修改 \lstinline{tp}。
\lstinline{start} 在 CPU 启动序列的早期设置 \lstinline{tp}
寄存器，同时仍在机器模式中
\lineref{kernel/start.c:/w_tp/}。
在准备返回到用户空间时，
\lstinline{prepare_return} 将 \lstinline{tp} 保存在跳板页中，以防用户代码修改它。
最后，\lstinline{uservec} 在从用户空间进入内核时恢复保存的 \lstinline{tp}
\lineref{kernel/trampoline.S:/make tp hold/}。
编译器保证在内核代码中从不修改 \lstinline{tp}。
如果 xv6 可以在需要时询问 RISC-V 硬件获取当前 hartid 会更方便，
但 RISC-V 只允许在机器模式中这样做，而在监管者模式中不行。

\lstinline{cpuid}
和
\lstinline{mycpu}
的返回值是脆弱的：
如果定时器中断并导致线程让出，然后稍后在不同的 CPU 上恢复执行，
先前返回的值将不再正确。
为了避免这个问题，xv6 要求代码在调用 {\tt cpuid()} 或 {\tt mycpu()} 之前禁用中断，
并且只在使用完返回值后才启用中断。

函数
\indexcode{myproc}
\lineref{kernel/proc.c:/^myproc/}
返回当前 CPU 上运行的进程的
\lstinline{struct proc}
指针。
\lstinline{myproc}
禁用中断，调用
\lstinline{mycpu}，
从
\lstinline{struct cpu}
中获取当前进程指针
(\lstinline{c->proc})，
然后启用中断。
即使启用了中断，
\lstinline{myproc}
的返回值也是安全可用的：
如果定时器中断将调用进程移动到不同的 CPU，其
\lstinline{struct proc}
指针将保持不变。


%%
\section{现实世界}
%%

xv6 调度器实现了一个简单的调度策略，轮流执行每个进程。
这种策略叫做
\indextext{轮询（round robin）}。
真实操作系统实现了更复杂的策略，例如，允许进程具有优先级。
其思想是可运行的高优先级进程将被调度器优先于可运行的低优先级进程。
这些策略可能变得复杂，因为通常有相互竞争的目标：
例如，操作系统可能还想保证公平性和高吞吐量。

%%
\section{练习}
%%

\begin{enumerate}

\item 修改 xv6，使得从一个进程的内核线程切换到另一个进程时只使用一次上下文切换，
  而不是通过调度器线程切换。让出的线程需要自己选择下一个线程并调用 \texttt{swtch}。
  挑战将是防止多个 CPU 意外执行同一线程；
  正确处理锁定；以及避免死锁。

\end{enumerate}

\chapter{睡眠和唤醒}
\label{CH:SLEEP}
%%

调度和锁机制隐藏了一个线程的操作，
使其免受其他线程的影响，
但我们还需要支持线程有意交互的抽象机制。
例如，xv6 中管道的读取者可能需要等待
写入进程生成数据；
父进程对 \lstinline{wait} 的调用可能需要等待子进程退出；
且读取磁盘的进程需要等待
磁盘硬件完成读取。
xv6 内核在这些情况（以及许多其他情况）下使用睡眠和唤醒机制。
睡眠机制允许内核线程等待特定条件成立；
另一线程或中断处理程序可使条件成立（通常通过修改某些变量），
随后调用唤醒，通知等待该条件的线程恢复执行。
睡眠和唤醒通常被称为
\indextext{序列协调（sequence coordination）}
或
\indextext{条件同步（conditional synchronization）}
机制。

在继续之前，请先阅读 {\tt kernel/proc.c} 中的函数 {\tt sleep()} 和 {\tt
  wakeup()}，以及 {\tt kernel/pipe.c} 的全部内容。

\section{概述}

睡眠/唤醒接口如下所示：

\begin{lstlisting}
void sleep(void *chan, struct spinlock *lk)
void wakeup(void *chan)
\end{lstlisting}

\lstinline{sleep()} 将调用进程置为 \lstinline{SLEEPING}
（而非 \lstinline{RUNNABLE}），并通过上下文切换至调度器来释放 CPU，
以便其他进程可以运行。
\lstinline{chan} 参数称为 \indextext{等待通道（wait channel）}。
\lstinline{wakeup(chan)} 唤醒所有调用 \lstinline{sleep(chan, ...)} 且使用相同 \lstinline{chan}
值的进程（如果有的话）。\lstinline{sleep} 和 \lstinline{wakeup} 将 \lstinline{chan}
视为不透明的 64 位值；它们对 \lstinline{chan} 所做的唯一操作是相等性比较。
通常的模式是调用者传递某个方便的对象的地址作为 \lstinline{chan} 参数。

内核代码调用 \lstinline{sleep} 来等待某个
\indextext{条件（condition）}
变为真。例如，如果从管道读取的内核代码发现管道缓冲区
当前为空，则调用 \lstinline{sleep}；
这种情况下的条件是管道缓冲区变为非空（由于另一个进程写入管道）。
\lstinline{sleep} 和 \lstinline{wakeup} 不知道条件是什么：只有调用代码知道。
通常的模式是调用者首先检查条件，如果条件不为真则调用 \lstinline{sleep}；
稍后使条件为真的代码调用 \lstinline{wakeup}。

以下是 xv6 内核管道代码如何使用
\lstinline{sleep} 和 \lstinline{wakeup} 的草图：

\begin{lstlisting}
piperead(pipe){
  acquire(&pipe->lock);
  while(there's no data in pipe->buffer){
    (*@\textcolor{blue}{// ZZZ}@*)
    sleep(&pipe, &pipe->lock);
  }
  remove the data from the pipe;
  release(&pipe->lock);
}

pipewrite(pipe){
  acquire(&pipe->lock);
  append data to pipe->buffer;
  wakeup(&pipe);
  release(&pipe->lock);
}
\end{lstlisting}

这段代码使用管道数据结构的地址作为等待通道。

\lstinline{lk} 参数对 \lstinline{sleep} 有什么作用？
在所有 \lstinline{sleep}/\lstinline{wakeup} 的使用中，
条件涉及共享数据，由睡眠线程和调用 \lstinline{wakeup} 的线程使用，
因此总有一个保护条件的锁。该锁称为 \indextext{条件锁（condition lock）}。
在上面的管道代码中，两个函数在持有管道锁的同时使用管道及其缓冲区，
在这种情况下，管道锁也是条件锁。
规则是任何调用 \lstinline{sleep} 或 \lstinline{wakeup} 的代码必须持有条件锁，
而且该锁必须作为第二个参数传递给 \lstinline{sleep}。

在调用 \lstinline{sleep} 时必须持有条件锁，
并且必须将其传递给 \lstinline{sleep} 的原因是防止另一个线程可能在条件检查
和调用 \lstinline{sleep} 之间调用 \lstinline{wakeup} 的可能性。
在该点调用 \lstinline{wakeup} 将找不到要唤醒的睡眠进程；
\lstinline{wakeup} 将直接返回。
但随后对 \lstinline{sleep} 的调用可能永远不会醒来，
因为原本打算唤醒它的 \lstinline{wakeup} 已经发生。
这种不希望发生的情况称为 \indextext{丢失唤醒（lost wake-up）}。

在上面的管道示例中，避免的丢失唤醒的可能性是
另一个 CPU 上的线程可能在标记为 {\tt ZZZ} 的点调用 \lstinline{pipewrite}，
在 \lstinline{piperead} 检查条件和调用 \lstinline{sleep} 之间。
\lstinline{piperead} 在检查条件和调用 \lstinline{sleep} 之间持有管道锁这一事实阻止了 \lstinline{pipewrite} 执行，
从而防止了丢失唤醒。

{\tt sleep()} 释放条件锁以便调用 {\tt wakeup()} 的代码可以继续执行。
{\tt sleep()} 还上下文切换到调度器，以便在等待时让其他线程运行。
实现以相对于 {\tt wakeup()} 原子的（不可分割的）方式执行这两个步骤，以防止丢失唤醒。

%%
\section{代码：睡眠和唤醒}
%%

Xv6 的
\indexcode{sleep}
\lineref{kernel/proc.c:/^sleep/}
和
\indexcode{wakeup}
\lineref{kernel/proc.c:/^wakeup/}
实现了上面示例中使用的接口。
基本思想是让
\lstinline{sleep}
将当前进程标记为
\indexcode{SLEEPING}，
然后调用
\indexcode{sched}
释放 CPU；
\lstinline{wakeup}
查找在给定等待通道上睡眠的进程，
并将其标记为
\indexcode{RUNNABLE}。

\lstinline{sleep}
获取
\indexcode{p->lock}
\lineref{kernel/proc.c:/DOC: sleeplock1/}，
然后{\it 才}释放
条件锁
\lstinline{lk}。
\lstinline{sleep} 在任何时候都持有这两个锁中的一个这一事实
阻止了并发 \lstinline{wakeup}
（必须获取并持有两个锁）起作用，
从而防止了丢失唤醒。
现在
\lstinline{sleep}
只持有
\lstinline{p->lock}，
它可以通过记录等待通道、
将进程状态更改为 \texttt{SLEEPING}、
并调用
\lstinline{sched}
\linerefs{kernel/proc.c:/chan.=.chan/,/sched/}
来使进程睡眠。
稍后就会清楚为什么在进程被标记为 \texttt{SLEEPING} 之前不释放 \lstinline{p->lock}（由 \lstinline{scheduler}）是至关重要的。

在某个时刻，一个进程将获取条件锁，
设置睡眠者等待的条件，
并调用 \lstinline{wakeup(chan)}。
在持有条件锁时调用 \lstinline{wakeup} 很重要\footnote{%
%
严格来说，只要
\lstinline{wakeup}
在
\lstinline{acquire}
之后（即，可以在
\lstinline{release}
之后调用
\lstinline{wakeup}）就足够了。%
%
}。
\lstinline{wakeup}
遍历进程表
\lineref{kernel/proc.c:/^wakeup\(/}。
它获取所检查的每个进程的
\lstinline{p->lock}。
当 \lstinline{wakeup} 找到状态为
\indexcode{SLEEPING}
且
\indexcode{chan}
匹配的进程时，
它将该进程的状态更改为
\indexcode{RUNNABLE}。
下次 \lstinline{scheduler} 运行时，它将看到该进程已准备好运行。

\begin{figure}[t]
  \chapter{睡眠和唤醒}
\label{CH:SLEEP}
%%

调度和锁机制隐藏了一个线程的操作，
使其免受其他线程的影响，
但我们还需要支持线程有意交互的抽象机制。
例如，xv6 中管道的读取者可能需要等待
写入进程生成数据；
父进程对 \lstinline{wait} 的调用可能需要等待子进程退出；
且读取磁盘的进程需要等待
磁盘硬件完成读取。
xv6 内核在这些情况（以及许多其他情况）下使用睡眠和唤醒机制。
睡眠机制允许内核线程等待特定条件成立；
另一线程或中断处理程序可使条件成立（通常通过修改某些变量），
随后调用唤醒，通知等待该条件的线程恢复执行。
睡眠和唤醒通常被称为
\indextext{序列协调（sequence coordination）}
或
\indextext{条件同步（conditional synchronization）}
机制。

在继续之前，请先阅读 {\tt kernel/proc.c} 中的函数 {\tt sleep()} 和 {\tt
  wakeup()}，以及 {\tt kernel/pipe.c} 的全部内容。

\section{概述}

睡眠/唤醒接口如下所示：

\begin{lstlisting}
void sleep(void *chan, struct spinlock *lk)
void wakeup(void *chan)
\end{lstlisting}

\lstinline{sleep()} 将调用进程置为 \lstinline{SLEEPING}
（而非 \lstinline{RUNNABLE}），并通过上下文切换至调度器来释放 CPU，
以便其他进程可以运行。
\lstinline{chan} 参数称为 \indextext{等待通道（wait channel）}。
\lstinline{wakeup(chan)} 唤醒所有调用 \lstinline{sleep(chan, ...)} 且使用相同 \lstinline{chan}
值的进程（如果有的话）。\lstinline{sleep} 和 \lstinline{wakeup} 将 \lstinline{chan}
视为不透明的 64 位值；它们对 \lstinline{chan} 所做的唯一操作是相等性比较。
通常的模式是调用者传递某个方便的对象的地址作为 \lstinline{chan} 参数。

内核代码调用 \lstinline{sleep} 来等待某个
\indextext{条件（condition）}
变为真。例如，如果从管道读取的内核代码发现管道缓冲区
当前为空，则调用 \lstinline{sleep}；
这种情况下的条件是管道缓冲区变为非空（由于另一个进程写入管道）。
\lstinline{sleep} 和 \lstinline{wakeup} 不知道条件是什么：只有调用代码知道。
通常的模式是调用者首先检查条件，如果条件不为真则调用 \lstinline{sleep}；
稍后使条件为真的代码调用 \lstinline{wakeup}。

以下是 xv6 内核管道代码如何使用
\lstinline{sleep} 和 \lstinline{wakeup} 的草图：

\begin{lstlisting}
piperead(pipe){
  acquire(&pipe->lock);
  while(there's no data in pipe->buffer){
    (*@\textcolor{blue}{// ZZZ}@*)
    sleep(&pipe, &pipe->lock);
  }
  remove the data from the pipe;
  release(&pipe->lock);
}

pipewrite(pipe){
  acquire(&pipe->lock);
  append data to pipe->buffer;
  wakeup(&pipe);
  release(&pipe->lock);
}
\end{lstlisting}

这段代码使用管道数据结构的地址作为等待通道。

\lstinline{lk} 参数对 \lstinline{sleep} 有什么作用？
在所有 \lstinline{sleep}/\lstinline{wakeup} 的使用中，
条件涉及共享数据，由睡眠线程和调用 \lstinline{wakeup} 的线程使用，
因此总有一个保护条件的锁。该锁称为 \indextext{条件锁（condition lock）}。
在上面的管道代码中，两个函数在持有管道锁的同时使用管道及其缓冲区，
在这种情况下，管道锁也是条件锁。
规则是任何调用 \lstinline{sleep} 或 \lstinline{wakeup} 的代码必须持有条件锁，
而且该锁必须作为第二个参数传递给 \lstinline{sleep}。

在调用 \lstinline{sleep} 时必须持有条件锁，
并且必须将其传递给 \lstinline{sleep} 的原因是防止另一个线程可能在条件检查
和调用 \lstinline{sleep} 之间调用 \lstinline{wakeup} 的可能性。
在该点调用 \lstinline{wakeup} 将找不到要唤醒的睡眠进程；
\lstinline{wakeup} 将直接返回。
但随后对 \lstinline{sleep} 的调用可能永远不会醒来，
因为原本打算唤醒它的 \lstinline{wakeup} 已经发生。
这种不希望发生的情况称为 \indextext{丢失唤醒（lost wake-up）}。

在上面的管道示例中，避免的丢失唤醒的可能性是
另一个 CPU 上的线程可能在标记为 {\tt ZZZ} 的点调用 \lstinline{pipewrite}，
在 \lstinline{piperead} 检查条件和调用 \lstinline{sleep} 之间。
\lstinline{piperead} 在检查条件和调用 \lstinline{sleep} 之间持有管道锁这一事实阻止了 \lstinline{pipewrite} 执行，
从而防止了丢失唤醒。

{\tt sleep()} 释放条件锁以便调用 {\tt wakeup()} 的代码可以继续执行。
{\tt sleep()} 还上下文切换到调度器，以便在等待时让其他线程运行。
实现以相对于 {\tt wakeup()} 原子的（不可分割的）方式执行这两个步骤，以防止丢失唤醒。

%%
\section{代码：睡眠和唤醒}
%%

Xv6 的
\indexcode{sleep}
\lineref{kernel/proc.c:/^sleep/}
和
\indexcode{wakeup}
\lineref{kernel/proc.c:/^wakeup/}
实现了上面示例中使用的接口。
基本思想是让
\lstinline{sleep}
将当前进程标记为
\indexcode{SLEEPING}，
然后调用
\indexcode{sched}
释放 CPU；
\lstinline{wakeup}
查找在给定等待通道上睡眠的进程，
并将其标记为
\indexcode{RUNNABLE}。

\lstinline{sleep}
获取
\indexcode{p->lock}
\lineref{kernel/proc.c:/DOC: sleeplock1/}，
然后{\it 才}释放
条件锁
\lstinline{lk}。
\lstinline{sleep} 在任何时候都持有这两个锁中的一个这一事实
阻止了并发 \lstinline{wakeup}
（必须获取并持有两个锁）起作用，
从而防止了丢失唤醒。
现在
\lstinline{sleep}
只持有
\lstinline{p->lock}，
它可以通过记录等待通道、
将进程状态更改为 \texttt{SLEEPING}、
并调用
\lstinline{sched}
\linerefs{kernel/proc.c:/chan.=.chan/,/sched/}
来使进程睡眠。
稍后就会清楚为什么在进程被标记为 \texttt{SLEEPING} 之前不释放 \lstinline{p->lock}（由 \lstinline{scheduler}）是至关重要的。

在某个时刻，一个进程将获取条件锁，
设置睡眠者等待的条件，
并调用 \lstinline{wakeup(chan)}。
在持有条件锁时调用 \lstinline{wakeup} 很重要\footnote{%
%
严格来说，只要
\lstinline{wakeup}
在
\lstinline{acquire}
之后（即，可以在
\lstinline{release}
之后调用
\lstinline{wakeup}）就足够了。%
%
}。
\lstinline{wakeup}
遍历进程表
\lineref{kernel/proc.c:/^wakeup\(/}。
它获取所检查的每个进程的
\lstinline{p->lock}。
当 \lstinline{wakeup} 找到状态为
\indexcode{SLEEPING}
且
\indexcode{chan}
匹配的进程时，
它将该进程的状态更改为
\indexcode{RUNNABLE}。
下次 \lstinline{scheduler} 运行时，它将看到该进程已准备好运行。

\begin{figure}[t]
  \chapter{睡眠和唤醒}
\label{CH:SLEEP}
%%

调度和锁机制隐藏了一个线程的操作，
使其免受其他线程的影响，
但我们还需要支持线程有意交互的抽象机制。
例如，xv6 中管道的读取者可能需要等待
写入进程生成数据；
父进程对 \lstinline{wait} 的调用可能需要等待子进程退出；
且读取磁盘的进程需要等待
磁盘硬件完成读取。
xv6 内核在这些情况（以及许多其他情况）下使用睡眠和唤醒机制。
睡眠机制允许内核线程等待特定条件成立；
另一线程或中断处理程序可使条件成立（通常通过修改某些变量），
随后调用唤醒，通知等待该条件的线程恢复执行。
睡眠和唤醒通常被称为
\indextext{序列协调（sequence coordination）}
或
\indextext{条件同步（conditional synchronization）}
机制。

在继续之前，请先阅读 {\tt kernel/proc.c} 中的函数 {\tt sleep()} 和 {\tt
  wakeup()}，以及 {\tt kernel/pipe.c} 的全部内容。

\section{概述}

睡眠/唤醒接口如下所示：

\begin{lstlisting}
void sleep(void *chan, struct spinlock *lk)
void wakeup(void *chan)
\end{lstlisting}

\lstinline{sleep()} 将调用进程置为 \lstinline{SLEEPING}
（而非 \lstinline{RUNNABLE}），并通过上下文切换至调度器来释放 CPU，
以便其他进程可以运行。
\lstinline{chan} 参数称为 \indextext{等待通道（wait channel）}。
\lstinline{wakeup(chan)} 唤醒所有调用 \lstinline{sleep(chan, ...)} 且使用相同 \lstinline{chan}
值的进程（如果有的话）。\lstinline{sleep} 和 \lstinline{wakeup} 将 \lstinline{chan}
视为不透明的 64 位值；它们对 \lstinline{chan} 所做的唯一操作是相等性比较。
通常的模式是调用者传递某个方便的对象的地址作为 \lstinline{chan} 参数。

内核代码调用 \lstinline{sleep} 来等待某个
\indextext{条件（condition）}
变为真。例如，如果从管道读取的内核代码发现管道缓冲区
当前为空，则调用 \lstinline{sleep}；
这种情况下的条件是管道缓冲区变为非空（由于另一个进程写入管道）。
\lstinline{sleep} 和 \lstinline{wakeup} 不知道条件是什么：只有调用代码知道。
通常的模式是调用者首先检查条件，如果条件不为真则调用 \lstinline{sleep}；
稍后使条件为真的代码调用 \lstinline{wakeup}。

以下是 xv6 内核管道代码如何使用
\lstinline{sleep} 和 \lstinline{wakeup} 的草图：

\begin{lstlisting}
piperead(pipe){
  acquire(&pipe->lock);
  while(there's no data in pipe->buffer){
    (*@\textcolor{blue}{// ZZZ}@*)
    sleep(&pipe, &pipe->lock);
  }
  remove the data from the pipe;
  release(&pipe->lock);
}

pipewrite(pipe){
  acquire(&pipe->lock);
  append data to pipe->buffer;
  wakeup(&pipe);
  release(&pipe->lock);
}
\end{lstlisting}

这段代码使用管道数据结构的地址作为等待通道。

\lstinline{lk} 参数对 \lstinline{sleep} 有什么作用？
在所有 \lstinline{sleep}/\lstinline{wakeup} 的使用中，
条件涉及共享数据，由睡眠线程和调用 \lstinline{wakeup} 的线程使用，
因此总有一个保护条件的锁。该锁称为 \indextext{条件锁（condition lock）}。
在上面的管道代码中，两个函数在持有管道锁的同时使用管道及其缓冲区，
在这种情况下，管道锁也是条件锁。
规则是任何调用 \lstinline{sleep} 或 \lstinline{wakeup} 的代码必须持有条件锁，
而且该锁必须作为第二个参数传递给 \lstinline{sleep}。

在调用 \lstinline{sleep} 时必须持有条件锁，
并且必须将其传递给 \lstinline{sleep} 的原因是防止另一个线程可能在条件检查
和调用 \lstinline{sleep} 之间调用 \lstinline{wakeup} 的可能性。
在该点调用 \lstinline{wakeup} 将找不到要唤醒的睡眠进程；
\lstinline{wakeup} 将直接返回。
但随后对 \lstinline{sleep} 的调用可能永远不会醒来，
因为原本打算唤醒它的 \lstinline{wakeup} 已经发生。
这种不希望发生的情况称为 \indextext{丢失唤醒（lost wake-up）}。

在上面的管道示例中，避免的丢失唤醒的可能性是
另一个 CPU 上的线程可能在标记为 {\tt ZZZ} 的点调用 \lstinline{pipewrite}，
在 \lstinline{piperead} 检查条件和调用 \lstinline{sleep} 之间。
\lstinline{piperead} 在检查条件和调用 \lstinline{sleep} 之间持有管道锁这一事实阻止了 \lstinline{pipewrite} 执行，
从而防止了丢失唤醒。

{\tt sleep()} 释放条件锁以便调用 {\tt wakeup()} 的代码可以继续执行。
{\tt sleep()} 还上下文切换到调度器，以便在等待时让其他线程运行。
实现以相对于 {\tt wakeup()} 原子的（不可分割的）方式执行这两个步骤，以防止丢失唤醒。

%%
\section{代码：睡眠和唤醒}
%%

Xv6 的
\indexcode{sleep}
\lineref{kernel/proc.c:/^sleep/}
和
\indexcode{wakeup}
\lineref{kernel/proc.c:/^wakeup/}
实现了上面示例中使用的接口。
基本思想是让
\lstinline{sleep}
将当前进程标记为
\indexcode{SLEEPING}，
然后调用
\indexcode{sched}
释放 CPU；
\lstinline{wakeup}
查找在给定等待通道上睡眠的进程，
并将其标记为
\indexcode{RUNNABLE}。

\lstinline{sleep}
获取
\indexcode{p->lock}
\lineref{kernel/proc.c:/DOC: sleeplock1/}，
然后{\it 才}释放
条件锁
\lstinline{lk}。
\lstinline{sleep} 在任何时候都持有这两个锁中的一个这一事实
阻止了并发 \lstinline{wakeup}
（必须获取并持有两个锁）起作用，
从而防止了丢失唤醒。
现在
\lstinline{sleep}
只持有
\lstinline{p->lock}，
它可以通过记录等待通道、
将进程状态更改为 \texttt{SLEEPING}、
并调用
\lstinline{sched}
\linerefs{kernel/proc.c:/chan.=.chan/,/sched/}
来使进程睡眠。
稍后就会清楚为什么在进程被标记为 \texttt{SLEEPING} 之前不释放 \lstinline{p->lock}（由 \lstinline{scheduler}）是至关重要的。

在某个时刻，一个进程将获取条件锁，
设置睡眠者等待的条件，
并调用 \lstinline{wakeup(chan)}。
在持有条件锁时调用 \lstinline{wakeup} 很重要\footnote{%
%
严格来说，只要
\lstinline{wakeup}
在
\lstinline{acquire}
之后（即，可以在
\lstinline{release}
之后调用
\lstinline{wakeup}）就足够了。%
%
}。
\lstinline{wakeup}
遍历进程表
\lineref{kernel/proc.c:/^wakeup\(/}。
它获取所检查的每个进程的
\lstinline{p->lock}。
当 \lstinline{wakeup} 找到状态为
\indexcode{SLEEPING}
且
\indexcode{chan}
匹配的进程时，
它将该进程的状态更改为
\indexcode{RUNNABLE}。
下次 \lstinline{scheduler} 运行时，它将看到该进程已准备好运行。

\begin{figure}[t]
  \input{fig/sleep.tex}
  \caption{重叠锁以避免丢失唤醒}
    \label{fig:overlap}
\end{figure}

\lstinline{sleep}
和
\lstinline{wakeup}
的锁定规则为什么能确保即将睡眠的进程不会错过并发的唤醒？
即将睡眠的进程在检查条件之前直到标记自己为 \texttt{SLEEPING} 之后，持有
条件锁或它自己的
\lstinline{p->lock}
或两者；
参见图~\ref{fig:overlap}。
调用 \texttt{wakeup} 的进程需要获取{\it 两个}锁。
唤醒者可能首先获取锁，这意味着它将在消费线程检查条件之前使条件为真，
而且消费线程不需要调用 {\tt sleep()}；
或者唤醒者的 {\tt acquire()} 可能需要等待，直到消费线程完全完成睡眠并
释放锁，在这种情况下唤醒者将看到消费线程被标记为 {\tt SLEEPING}
并将唤醒它。

有时多个进程在同一个通道上睡眠；
例如，多个进程从管道读取。
对
\lstinline{wakeup}
的单个调用将唤醒它们所有。
其中一个将首先运行并获取 \lstinline{sleep}
被调用时使用的锁，并（在管道的情况下）读取等待的数据。
其他进程将发现，尽管被唤醒了，
但没有数据可读取。
从它们的角度来看，唤醒是``虚假的（spurious）''，
它们必须再次睡眠。
因此 \lstinline{sleep} 总是在重新检查条件的循环内调用，如上面的 \lstinline{P} 所示。

如果睡眠/唤醒 的两个使用意外选择了相同的通道：
它们会看到虚假的唤醒，
但如上所述的循环将容忍这个问题。
睡眠/唤醒 的大部分魅力在于它既轻量级（不需要创建特殊的数据结构作为等待通道）
又提供了一层间接（调用者不需要知道它们正在与哪个特定进程交互）。
%%
\section{代码：管道}
%%
xv6 的管道是使用
\lstinline{sleep}
和 \lstinline{wakeup} 来同步生产者和消费者的代码示例。
我们在第~\ref{CH:UNIX} 章看到了管道的接口：
写入管道一端的数据被复制到内核缓冲区，
然后可以从管道的另一端读取。
后续章节将检查围绕管道的文件描述符支持，
但现在让我们看看
\indexcode{pipewrite}
和
\indexcode{piperead}
的实现。

每个管道由
\indexcode{struct pipe}
表示，它包含一个
\lstinline{lock}
和一个
\lstinline{data}
缓冲区。
字段
\lstinline{nread}
和
\lstinline{nwrite}
分别计数从缓冲区读取和写入的字节总数。
缓冲区环绕：
在
\lstinline{buf[PIPESIZE-1]}
之后写入的下一个字节是
\lstinline{buf[0]}。
计数不环绕。
这种约定让实现能够区分满缓冲区
(\lstinline{nwrite}
\lstinline{==}
\lstinline{nread+PIPESIZE})
和空缓冲区
(\lstinline{nwrite}
\lstinline{==}
\lstinline{nread})，
但它意味着索引到缓冲区时必须使用
\lstinline{buf[nread}
\lstinline{%}
\lstinline{PIPESIZE]}
而不是仅仅
\lstinline{buf[nread]}
（对
\lstinline{nwrite}
同样如此）。

假设对
\lstinline{piperead}
和
\lstinline{pipewrite}
的调用同时在两个不同的 CPU 上发生。
\lstinline{pipewrite}
\lineref{kernel/pipe.c:/^pipewrite/}
首先获取管道的锁，该锁保护计数、数据及其相关的
不变量。
\lstinline{piperead}
\lineref{kernel/pipe.c:/^piperead/}
然后也尝试获取锁，但无法获取。
它在
\lstinline{acquire}
\lineref{kernel/spinlock.c:/^acquire/}
中自旋等待锁。
在
\lstinline{piperead}
等待时，
\lstinline{pipewrite}
循环遍历要写入的字节
(\lstinline{addr[0..n-1]})，
依次将每个字节添加到管道中
\lineref{kernel/pipe.c:/nwrite\+\+/}。
在这个循环期间，缓冲区可能填满
\lineref{kernel/pipe.c:/DOC: pipewrite-full/}。
在这种情况下，
\lstinline{pipewrite}
调用
\lstinline{wakeup}
来提醒任何睡眠的读取者缓冲区中有数据等待，
然后在
\lstinline{&pi->nwrite}
上睡眠以等待读取者从缓冲区中取出一些字节。
\lstinline{sleep}
释放管道锁作为将
\lstinline{pipewrite}
的进程置于睡眠的一部分。

\lstinline{piperead}
现在获取管道锁并进入其临界区：
它发现
\lstinline{pi->nread}
\lstinline{!=}
\lstinline{pi->nwrite}
\lineref{kernel/pipe.c:/DOC: pipe-empty/}
(\lstinline{pipewrite}
进入睡眠是因为
\lstinline{pi->nwrite}
\lstinline{==}
\lstinline{pi->nread}
\lstinline{+}
\lstinline{PIPESIZE}
\lineref{kernel/pipe.c:/pipewrite-full/})，
因此它进入
\lstinline{for}
循环，从管道中复制数据
\lineref{kernel/pipe.c:/DOC: piperead-copy/}，
并将
\lstinline{nread}
增加复制的字节数。
缓冲区中的那么多空间现在可用于写入，因此
\lstinline{piperead}
调用
\lstinline{wakeup}
\lineref{kernel/pipe.c:/DOC: piperead-wakeup/}
来唤醒任何睡眠的写入者，然后返回。
\lstinline{wakeup}
找到一个在
\lstinline{&pi->nwrite}
上睡眠的进程，
该进程正在运行
\lstinline{pipewrite}
但在缓冲区填满时停止了。
它将该进程标记为
\indexcode{RUNNABLE}。

管道代码为读取者和写入者使用单独的等待通道
(\lstinline{pi->nread}
和
\lstinline{pi->nwrite})；
这在不太可能有大量读取者和写入者等待同一管道的情况下可能使系统更高效。
管道代码在检查睡眠条件的循环内睡眠；
如果有多个读取者或写入者，除第一个唤醒的进程外，所有其他进程将看到条件为假并再次睡眠。
%%
\section{代码：Wait、exit 和 kill}

请阅读 {\tt kernel/proc.c} 中函数 {\tt kwait()}、{\tt kexit()}
和 {\tt kkill()} 的代码；这些是对应系统调用的内部实现。

%%
\lstinline{sleep}
和
\lstinline{wakeup}
可用于许多类型的等待。
第~\ref{CH:UNIX} 章介绍了一个有趣的例子，
即子进程 \indexcode{exit}
与其父进程 \indexcode{wait} 之间的交互。
在子进程死亡时，父进程可能已经在 {\tt wait} 中睡眠，
或者可能在做其他事情；
在后一种情况下，后续对 {\tt wait} 的调用必须观察子进程的死亡，
可能在调用 {\tt exit} 很久之后。
Xv6 记录子进程死亡直到 {\tt wait} 观察到它的方式是 {\tt exit} 将调用者放入 \indexcode{ZOMBIE}
状态，在那里保持直到父进程的 {\tt wait} 注意到它，将子进程的状态更改为 {\tt UNUSED}，
复制子进程的退出状态，并将子进程的进程 ID 返回给父进程。
如果父进程在子进程之前退出，
父进程将子进程交给
\lstinline{init}
进程，该进程永久调用 {\tt wait}；
因此每个子进程都有一个父进程来清理它。
一个挑战是避免同时父进程
\lstinline{wait}
和子进程
\lstinline{exit}
之间的竞争和死锁，
以及同时
\lstinline{exit} 和 \lstinline{exit} 之间的竞争。

\lstinline{kwait}，{\tt wait} 的内核实现，首先获取
\lstinline{wait_lock}
\lineref{kernel/proc.c:/^kwait/}，
它作为条件锁有助于确保 \lstinline{kwait} 不会错过来自退出子进程的 \lstinline{wakeup}。
然后 \lstinline{kwait} 扫描进程表。
如果它找到处于 \texttt{ZOMBIE} 状态的子进程，
它释放该子进程的资源和其 \lstinline{proc} 结构体，
将子进程的退出状态复制到提供给 \lstinline{wait} 的地址
（如果不是 0），
并返回子进程的进程 ID。
如果
\lstinline{kwait}
找到子进程但没有一个退出，
它调用
\lstinline{sleep}
等待它们中的任何一个退出
\lineref{kernel/proc.c:/DOC: wait-sleep/}，
然后再次扫描。
\lstinline{kwait} 经常持有两个锁，
\lstinline{wait_lock} 和某个进程的 \lstinline{pp->lock}；
避免死锁的顺序是先 \lstinline{wait_lock} 然后 \lstinline{pp->lock}。

\lstinline{kexit} \lineref{kernel/proc.c:/^kexit/} 记录退出状态，
释放一些资源，调用 \lstinline{reparent} 将任何子进程交给 \lstinline{init} 进程，
唤醒父进程以防它在 \lstinline{wait} 中睡眠，将调用者标记为僵尸，并永久让出 CPU。
\lstinline{kexit} 在此过程中持有
\lstinline{wait_lock} 和 \lstinline{p->lock}。
它持有 \lstinline{wait_lock} 因为它是
\lstinline{wakeup(p->parent)} 的条件锁，防止在 \lstinline{wait} 中的父进程丢失唤醒。
\lstinline{kexit} 还必须为此序列持有
\lstinline{p->lock}，以防止在 \lstinline{wait} 中的父进程在子进程最终调用
\lstinline{swtch} 之前看到子进程处于状态
\lstinline{ZOMBIE}。
\lstinline{kexit} 以与 \lstinline{kwait} 相同的顺序获取这些锁以避免死锁。

\lstinline{kexit} 在将其状态设置为 \lstinline{ZOMBIE} 之前唤醒父进程
可能看起来不正确，
但这是安全的：
尽管
\lstinline{wakeup}
可能导致父进程运行，
但父进程的
\lstinline{kwait}
中的循环在子进程的
\lstinline{p->lock}
被 {\tt scheduler} 释放之前无法检查子进程，
因此
\lstinline{kwait}
不能在
\lstinline{kexit}
将其状态设置为
\lstinline{ZOMBIE}
之前查看退出进程
\lineref{kernel/proc.c:/state.=.ZOMBIE/}。

虽然
\lstinline{exit}
允许进程终止自己，
但 \lstinline{kill} 系统调用
\lineref{kernel/proc.c:/^kkill/}
让一个进程请求另一个进程终止。
\lstinline{kill}
直接销毁受害者进程会太复杂，因为受害者
可能在另一个 CPU 上执行，
可能正处于更新内核数据结构的关键序列的中间。
因此
\lstinline{kkill}
做得很少：它只是设置受害者的
\indexcode{p->killed}
并且，如果它正在睡眠，就唤醒它。
最终受害者将进入或离开内核，
此时
\lstinline{usertrap}
中的代码将调用
\lstinline{kexit}
如果
\lstinline{p->killed}
被设置
（它通过调用
\lstinline{killed}
\lineref{kernel/proc.c:/^killed/}
来检查）。
如果受害者在用户空间运行，它将在下次进入内核时看到它已被终止，
通过进行系统调用或因为定时器（或其他设备）中断。

如果受害者进程在
\lstinline{sleep}
中，
\lstinline{kkill} 对
\lstinline{wakeup}
的调用将导致受害者从
\lstinline{sleep}
返回。
这有潜在危险，因为正在等待的条件可能不为真。
然而，xv6 对
\lstinline{sleep}
的调用总是包装在一个在
\lstinline{sleep}
返回后重新测试条件的
\lstinline{while}
循环中。
一些对
\lstinline{sleep}
的调用也在循环中测试
\lstinline{p->killed}，
如果设置则放弃当前活动。
这仅在放弃是正确的情况下才会这样做。
例如，管道读取和写入代码
\lineref{kernel/pipe.c:/killed.pr/}
如果设置了终止标志则返回；最终代码将返回到 trap，trap 将再次检查 \lstinline{p->killed} 并退出。

一些 xv6
\lstinline{sleep}
循环不检查
\lstinline{p->killed}，
因为代码处于多步系统调用的中间，应该是原子的（即，如果中途放弃会不正确）。
Virtio 驱动程序
\lineref{kernel/virtio\_disk.c:/sleep.b/}
是一个例子：它不检查
\lstinline{p->killed}
因为磁盘操作可能是一组写入中的一部分，所有这些写入都需要才能使文件系统保持正确状态。
在磁盘 I/O 上等待时被终止的进程直到完成当前系统调用并且
\lstinline{usertrap} 看到终止标志才会退出。

\section{进程锁}

与每个进程关联的锁（\lstinline{p->lock}）是 xv6 中最复杂的锁。
理解 \lstinline{p->lock} 的一个简单方法是，
在读取或写入以下任何 \lstinline{struct proc} 字段时必须持有它：
\lstinline{p->state}、
\lstinline{p->chan}、
\lstinline{p->killed}、
\lstinline{p->xstate}、
和
\lstinline{p->pid}。
这些字段可以被其他进程使用，或被其他 CPU 上的调度器线程使用，
因此它们自然必须由锁保护。

然而，\lstinline{p->lock} 的大多数使用是保护 xv6 进程数据结构和算法的更高级别不变量。
以下是 \lstinline{p->lock} 所做的全部事情：

% 很难知道如何以统一的方式表述这些事情。
% 讨论应该围绕不变量？保护数据？避免竞争？确保原子性？避免特定危险？

% "原子"表述通常不能真正说明危险是什么。

\begin{itemize}

\item 与 \lstinline{p->state} 一起，它防止在为新进程分配
  \lstinline{proc[]} 槽时的竞争。
% 使对 UNUSED 的检查与分配（或其他）原子化。

\item 它在进程被创建或销毁时将其隐藏起来。
% 使分配的所有步骤和销毁原子化。

\item 它防止父进程的 \lstinline{wait} 收集一个已将其状态设置为 \lstinline{ZOMBIE} 但尚未让出 CPU 的进程。
% 它将设置状态为 ZOMIE 和让出 CPU 原子化。

\item 它防止另一个 CPU 的调度器决定在进程将其状态设置为 \lstinline{RUNNABLE} 之后但在完成 \lstinline{swtch} 之前运行一个让出的进程。
% 它将设置 p->state 和 swtch() 原子化。

\item 它确保只有一个 CPU 的调度器决定运行一个
  \lstinline{RUNNABLE} 进程。
% 它将调度器对 RUNNABLE 的检查与实际运行进程（或设置状态为 RUNNING）原子化。

\item 它防止定时器中断导致进程在 \lstinline{swtch} 中时让出。
% 它使 swtch（或实际上是整个序列）相对于定时器中断原子化

\item 与条件锁一起，它有助于防止 \lstinline{wakeup}
忽略一个正在调用 \lstinline{sleep} 但尚未完成让出 CPU 的进程。
% 避免 conditionCheck+sleep 和 wakeup 之间的竞争？
% 使条件检查与让出原子化？

\item 它防止 \lstinline{kill} 的受害者进程在 \lstinline{kkill} 检查
  \lstinline{p->pid} 和设置 \lstinline{p->killed} 之间退出并可能被重新分配。
% 它使 pid 检查与设置 killed 原子化。

\item 它使 \lstinline{kkill} 对 \lstinline{p->state} 的检查和写入原子化。

\end{itemize}


\lstinline{p->parent} 字段由全局锁
\lstinline{wait_lock} 而不是 \lstinline{p->lock} 保护。
只有进程的父进程修改 \lstinline{p->parent}，
尽管该字段被进程本身和其他寻找其子进程的进程读取。
\lstinline{wait_lock} 的目的是作为条件锁，当
\lstinline{wait} 睡眠等待任何子进程退出时。
退出的子进程在将其状态设置为 \lstinline{ZOMBIE}、唤醒其父进程并让出 CPU 之后持有
\lstinline{wait_lock} 或 \lstinline{p->lock}。
\lstinline{wait_lock} 还序列化父进程和子进程的并发
\lstinline{exit}，
以便 \lstinline{init} 进程（继承子进程）
保证从其 \lstinline{wait} 中被唤醒。
\lstinline{wait_lock} 是全局锁而不是每个父进程的锁，
因为，直到进程获取它之前，它不知道其父进程是谁。

%%
\section{现实世界}
%%

\lstinline{sleep}
和
\lstinline{wakeup}
是一种简单有效的同步方法，
但还有许多其他方法；
信号量~
\cite{dijkstra65}
是一个例子。
所有这些方法的第一个挑战是避免我们在本章开头看到的``丢失唤醒''问题。
原始 Unix 内核的
\lstinline{sleep}
只是禁用中断，
这就足够了，因为 Unix 运行在单 CPU 系统上。
因为 xv6 运行在多处理器上，
它向
\lstinline{sleep}
添加了一个显式锁。
FreeBSD 的
\lstinline{msleep}
采用相同的方法。
Plan 9 的
\lstinline{sleep}
使用一个在即将睡眠之前持有调度锁运行的回调函数；
该函数作为睡眠条件的最后检查，以避免丢失唤醒。
Linux 内核的
\lstinline{sleep}
使用显式的进程队列，称为等待队列，而不是等待通道；该队列有自己的内部锁。

在
\lstinline{wakeup}
中扫描整个进程集是低效的。
更好的解决方案是将
\lstinline{sleep}
和
\lstinline{wakeup}
中的
\lstinline{chan}
替换为一个保存该结构上睡眠进程列表的数据结构，
如 Linux 的等待队列。
Plan 9 的
\lstinline{sleep}
和
\lstinline{wakeup}
将该结构称为会合点（rendezvous point）。
许多线程库将相同的结构称为条件变量（condition variable）；
在这种情况下，操作
\lstinline{sleep}
和
\lstinline{wakeup}
称为
\lstinline{wait}
和
\lstinline{signal}。
所有这些机制都有相同的特点：睡眠条件由某种在睡眠期间原子地释放的锁保护。

xv6 的
\lstinline{wakeup}
唤醒所有在特定等待通道上等待的进程。
如果有多个，它们都将尝试获取条件锁并重新检查条件；
在许多情况下只有一个能够做有用的事情（例如，读取管道中等待的所有数据）。
其余的将发现条件不再为真并返回睡眠；
唤醒它们是浪费 CPU 时间。
因此，
大多数条件变量设计提供两个原语：
\lstinline{signal}，
唤醒等待条件变量的进程之一，和
\lstinline{broadcast}，
唤醒它们所有。

强制终止进程带来一些问题。
例如，被终止的进程可能在内核深处睡眠，
展开其堆栈需要小心，因为调用堆栈上的每个函数可能需要做一些清理。
一些语言通过提供异常机制来帮助，但 C 没有。
此外，还有其他事件可以导致睡眠进程被唤醒，
即使它等待的事件尚未发生。
例如，当 Unix 进程睡眠时，另一个进程可能向它发送
\indexcode{signal}。
在这种情况下，进程将从被中断的系统调用返回，值为 -1 且错误代码设置为 EINTR。
应用程序可以检查这些值并决定做什么。
Xv6 不支持信号，这种复杂性不会出现。

Xv6 对
\lstinline{kill}
的支持不完全令人满意：有些睡眠循环可能应该检查
\lstinline{p->killed}。
一个相关的问题是，即使对于检查
\lstinline{p->killed}
的
\lstinline{sleep}
循环，
\lstinline{sleep}
和
\lstinline{kill}
之间也存在竞争；
后者可能设置
\lstinline{p->killed}
并尝试唤醒受害者，就在受害者的循环检查
\lstinline{p->killed}
之后但在其调用
\lstinline{sleep}
之前。
如果发生这个问题，受害者直到它等待的条件发生才会注意到
\lstinline{p->killed}。
这可能很久之后或永远不会
（例如，如果受害者正在等待控制台的输入，但用户没有键入任何输入）。

%%
\section{练习}
%%

\begin{enumerate}

\item 在 xv6 中实现计数信号量。
选择 xv6 的一些睡眠和唤醒用法并用信号量替换它们。
评判结果。

\item 你能实现一个只带一个参数的 {\tt sleep()} 变体吗，即通道，不需要锁参数？

\item 修复上面提到的
\lstinline{kill}
和
\lstinline{sleep}
之间的竞争，
使得在受害者的睡眠循环检查
\lstinline{p->killed}
之后但在调用
\lstinline{sleep}
之前发生的
\lstinline{kill}
导致受害者放弃当前系统调用。
% 答案：一个解决方案是在设置进程状态为睡眠之前在 sleep 中检查 p->killed 是否已设置。

\item 设计一个计划，使每个睡眠循环都检查
\lstinline{p->killed}
，以便例如，在 virtio 驱动程序中的进程如果被另一个进程终止，可以从 while 循环快速返回。
% 答案：这很困难。现代 Unix 使用 setjmp 和 longjmp 以及非常仔细的编程来清理被中断系统调用可能已建立的任何部分状态。

\end{enumerate}

  \caption{重叠锁以避免丢失唤醒}
    \label{fig:overlap}
\end{figure}

\lstinline{sleep}
和
\lstinline{wakeup}
的锁定规则为什么能确保即将睡眠的进程不会错过并发的唤醒？
即将睡眠的进程在检查条件之前直到标记自己为 \texttt{SLEEPING} 之后，持有
条件锁或它自己的
\lstinline{p->lock}
或两者；
参见图~\ref{fig:overlap}。
调用 \texttt{wakeup} 的进程需要获取{\it 两个}锁。
唤醒者可能首先获取锁，这意味着它将在消费线程检查条件之前使条件为真，
而且消费线程不需要调用 {\tt sleep()}；
或者唤醒者的 {\tt acquire()} 可能需要等待，直到消费线程完全完成睡眠并
释放锁，在这种情况下唤醒者将看到消费线程被标记为 {\tt SLEEPING}
并将唤醒它。

有时多个进程在同一个通道上睡眠；
例如，多个进程从管道读取。
对
\lstinline{wakeup}
的单个调用将唤醒它们所有。
其中一个将首先运行并获取 \lstinline{sleep}
被调用时使用的锁，并（在管道的情况下）读取等待的数据。
其他进程将发现，尽管被唤醒了，
但没有数据可读取。
从它们的角度来看，唤醒是``虚假的（spurious）''，
它们必须再次睡眠。
因此 \lstinline{sleep} 总是在重新检查条件的循环内调用，如上面的 \lstinline{P} 所示。

如果睡眠/唤醒 的两个使用意外选择了相同的通道：
它们会看到虚假的唤醒，
但如上所述的循环将容忍这个问题。
睡眠/唤醒 的大部分魅力在于它既轻量级（不需要创建特殊的数据结构作为等待通道）
又提供了一层间接（调用者不需要知道它们正在与哪个特定进程交互）。
%%
\section{代码：管道}
%%
xv6 的管道是使用
\lstinline{sleep}
和 \lstinline{wakeup} 来同步生产者和消费者的代码示例。
我们在第~\ref{CH:UNIX} 章看到了管道的接口：
写入管道一端的数据被复制到内核缓冲区，
然后可以从管道的另一端读取。
后续章节将检查围绕管道的文件描述符支持，
但现在让我们看看
\indexcode{pipewrite}
和
\indexcode{piperead}
的实现。

每个管道由
\indexcode{struct pipe}
表示，它包含一个
\lstinline{lock}
和一个
\lstinline{data}
缓冲区。
字段
\lstinline{nread}
和
\lstinline{nwrite}
分别计数从缓冲区读取和写入的字节总数。
缓冲区环绕：
在
\lstinline{buf[PIPESIZE-1]}
之后写入的下一个字节是
\lstinline{buf[0]}。
计数不环绕。
这种约定让实现能够区分满缓冲区
(\lstinline{nwrite}
\lstinline{==}
\lstinline{nread+PIPESIZE})
和空缓冲区
(\lstinline{nwrite}
\lstinline{==}
\lstinline{nread})，
但它意味着索引到缓冲区时必须使用
\lstinline{buf[nread}
\lstinline{%}
\lstinline{PIPESIZE]}
而不是仅仅
\lstinline{buf[nread]}
（对
\lstinline{nwrite}
同样如此）。

假设对
\lstinline{piperead}
和
\lstinline{pipewrite}
的调用同时在两个不同的 CPU 上发生。
\lstinline{pipewrite}
\lineref{kernel/pipe.c:/^pipewrite/}
首先获取管道的锁，该锁保护计数、数据及其相关的
不变量。
\lstinline{piperead}
\lineref{kernel/pipe.c:/^piperead/}
然后也尝试获取锁，但无法获取。
它在
\lstinline{acquire}
\lineref{kernel/spinlock.c:/^acquire/}
中自旋等待锁。
在
\lstinline{piperead}
等待时，
\lstinline{pipewrite}
循环遍历要写入的字节
(\lstinline{addr[0..n-1]})，
依次将每个字节添加到管道中
\lineref{kernel/pipe.c:/nwrite\+\+/}。
在这个循环期间，缓冲区可能填满
\lineref{kernel/pipe.c:/DOC: pipewrite-full/}。
在这种情况下，
\lstinline{pipewrite}
调用
\lstinline{wakeup}
来提醒任何睡眠的读取者缓冲区中有数据等待，
然后在
\lstinline{&pi->nwrite}
上睡眠以等待读取者从缓冲区中取出一些字节。
\lstinline{sleep}
释放管道锁作为将
\lstinline{pipewrite}
的进程置于睡眠的一部分。

\lstinline{piperead}
现在获取管道锁并进入其临界区：
它发现
\lstinline{pi->nread}
\lstinline{!=}
\lstinline{pi->nwrite}
\lineref{kernel/pipe.c:/DOC: pipe-empty/}
(\lstinline{pipewrite}
进入睡眠是因为
\lstinline{pi->nwrite}
\lstinline{==}
\lstinline{pi->nread}
\lstinline{+}
\lstinline{PIPESIZE}
\lineref{kernel/pipe.c:/pipewrite-full/})，
因此它进入
\lstinline{for}
循环，从管道中复制数据
\lineref{kernel/pipe.c:/DOC: piperead-copy/}，
并将
\lstinline{nread}
增加复制的字节数。
缓冲区中的那么多空间现在可用于写入，因此
\lstinline{piperead}
调用
\lstinline{wakeup}
\lineref{kernel/pipe.c:/DOC: piperead-wakeup/}
来唤醒任何睡眠的写入者，然后返回。
\lstinline{wakeup}
找到一个在
\lstinline{&pi->nwrite}
上睡眠的进程，
该进程正在运行
\lstinline{pipewrite}
但在缓冲区填满时停止了。
它将该进程标记为
\indexcode{RUNNABLE}。

管道代码为读取者和写入者使用单独的等待通道
(\lstinline{pi->nread}
和
\lstinline{pi->nwrite})；
这在不太可能有大量读取者和写入者等待同一管道的情况下可能使系统更高效。
管道代码在检查睡眠条件的循环内睡眠；
如果有多个读取者或写入者，除第一个唤醒的进程外，所有其他进程将看到条件为假并再次睡眠。
%%
\section{代码：Wait、exit 和 kill}

请阅读 {\tt kernel/proc.c} 中函数 {\tt kwait()}、{\tt kexit()}
和 {\tt kkill()} 的代码；这些是对应系统调用的内部实现。

%%
\lstinline{sleep}
和
\lstinline{wakeup}
可用于许多类型的等待。
第~\ref{CH:UNIX} 章介绍了一个有趣的例子，
即子进程 \indexcode{exit}
与其父进程 \indexcode{wait} 之间的交互。
在子进程死亡时，父进程可能已经在 {\tt wait} 中睡眠，
或者可能在做其他事情；
在后一种情况下，后续对 {\tt wait} 的调用必须观察子进程的死亡，
可能在调用 {\tt exit} 很久之后。
Xv6 记录子进程死亡直到 {\tt wait} 观察到它的方式是 {\tt exit} 将调用者放入 \indexcode{ZOMBIE}
状态，在那里保持直到父进程的 {\tt wait} 注意到它，将子进程的状态更改为 {\tt UNUSED}，
复制子进程的退出状态，并将子进程的进程 ID 返回给父进程。
如果父进程在子进程之前退出，
父进程将子进程交给
\lstinline{init}
进程，该进程永久调用 {\tt wait}；
因此每个子进程都有一个父进程来清理它。
一个挑战是避免同时父进程
\lstinline{wait}
和子进程
\lstinline{exit}
之间的竞争和死锁，
以及同时
\lstinline{exit} 和 \lstinline{exit} 之间的竞争。

\lstinline{kwait}，{\tt wait} 的内核实现，首先获取
\lstinline{wait_lock}
\lineref{kernel/proc.c:/^kwait/}，
它作为条件锁有助于确保 \lstinline{kwait} 不会错过来自退出子进程的 \lstinline{wakeup}。
然后 \lstinline{kwait} 扫描进程表。
如果它找到处于 \texttt{ZOMBIE} 状态的子进程，
它释放该子进程的资源和其 \lstinline{proc} 结构体，
将子进程的退出状态复制到提供给 \lstinline{wait} 的地址
（如果不是 0），
并返回子进程的进程 ID。
如果
\lstinline{kwait}
找到子进程但没有一个退出，
它调用
\lstinline{sleep}
等待它们中的任何一个退出
\lineref{kernel/proc.c:/DOC: wait-sleep/}，
然后再次扫描。
\lstinline{kwait} 经常持有两个锁，
\lstinline{wait_lock} 和某个进程的 \lstinline{pp->lock}；
避免死锁的顺序是先 \lstinline{wait_lock} 然后 \lstinline{pp->lock}。

\lstinline{kexit} \lineref{kernel/proc.c:/^kexit/} 记录退出状态，
释放一些资源，调用 \lstinline{reparent} 将任何子进程交给 \lstinline{init} 进程，
唤醒父进程以防它在 \lstinline{wait} 中睡眠，将调用者标记为僵尸，并永久让出 CPU。
\lstinline{kexit} 在此过程中持有
\lstinline{wait_lock} 和 \lstinline{p->lock}。
它持有 \lstinline{wait_lock} 因为它是
\lstinline{wakeup(p->parent)} 的条件锁，防止在 \lstinline{wait} 中的父进程丢失唤醒。
\lstinline{kexit} 还必须为此序列持有
\lstinline{p->lock}，以防止在 \lstinline{wait} 中的父进程在子进程最终调用
\lstinline{swtch} 之前看到子进程处于状态
\lstinline{ZOMBIE}。
\lstinline{kexit} 以与 \lstinline{kwait} 相同的顺序获取这些锁以避免死锁。

\lstinline{kexit} 在将其状态设置为 \lstinline{ZOMBIE} 之前唤醒父进程
可能看起来不正确，
但这是安全的：
尽管
\lstinline{wakeup}
可能导致父进程运行，
但父进程的
\lstinline{kwait}
中的循环在子进程的
\lstinline{p->lock}
被 {\tt scheduler} 释放之前无法检查子进程，
因此
\lstinline{kwait}
不能在
\lstinline{kexit}
将其状态设置为
\lstinline{ZOMBIE}
之前查看退出进程
\lineref{kernel/proc.c:/state.=.ZOMBIE/}。

虽然
\lstinline{exit}
允许进程终止自己，
但 \lstinline{kill} 系统调用
\lineref{kernel/proc.c:/^kkill/}
让一个进程请求另一个进程终止。
\lstinline{kill}
直接销毁受害者进程会太复杂，因为受害者
可能在另一个 CPU 上执行，
可能正处于更新内核数据结构的关键序列的中间。
因此
\lstinline{kkill}
做得很少：它只是设置受害者的
\indexcode{p->killed}
并且，如果它正在睡眠，就唤醒它。
最终受害者将进入或离开内核，
此时
\lstinline{usertrap}
中的代码将调用
\lstinline{kexit}
如果
\lstinline{p->killed}
被设置
（它通过调用
\lstinline{killed}
\lineref{kernel/proc.c:/^killed/}
来检查）。
如果受害者在用户空间运行，它将在下次进入内核时看到它已被终止，
通过进行系统调用或因为定时器（或其他设备）中断。

如果受害者进程在
\lstinline{sleep}
中，
\lstinline{kkill} 对
\lstinline{wakeup}
的调用将导致受害者从
\lstinline{sleep}
返回。
这有潜在危险，因为正在等待的条件可能不为真。
然而，xv6 对
\lstinline{sleep}
的调用总是包装在一个在
\lstinline{sleep}
返回后重新测试条件的
\lstinline{while}
循环中。
一些对
\lstinline{sleep}
的调用也在循环中测试
\lstinline{p->killed}，
如果设置则放弃当前活动。
这仅在放弃是正确的情况下才会这样做。
例如，管道读取和写入代码
\lineref{kernel/pipe.c:/killed.pr/}
如果设置了终止标志则返回；最终代码将返回到 trap，trap 将再次检查 \lstinline{p->killed} 并退出。

一些 xv6
\lstinline{sleep}
循环不检查
\lstinline{p->killed}，
因为代码处于多步系统调用的中间，应该是原子的（即，如果中途放弃会不正确）。
Virtio 驱动程序
\lineref{kernel/virtio\_disk.c:/sleep.b/}
是一个例子：它不检查
\lstinline{p->killed}
因为磁盘操作可能是一组写入中的一部分，所有这些写入都需要才能使文件系统保持正确状态。
在磁盘 I/O 上等待时被终止的进程直到完成当前系统调用并且
\lstinline{usertrap} 看到终止标志才会退出。

\section{进程锁}

与每个进程关联的锁（\lstinline{p->lock}）是 xv6 中最复杂的锁。
理解 \lstinline{p->lock} 的一个简单方法是，
在读取或写入以下任何 \lstinline{struct proc} 字段时必须持有它：
\lstinline{p->state}、
\lstinline{p->chan}、
\lstinline{p->killed}、
\lstinline{p->xstate}、
和
\lstinline{p->pid}。
这些字段可以被其他进程使用，或被其他 CPU 上的调度器线程使用，
因此它们自然必须由锁保护。

然而，\lstinline{p->lock} 的大多数使用是保护 xv6 进程数据结构和算法的更高级别不变量。
以下是 \lstinline{p->lock} 所做的全部事情：

% 很难知道如何以统一的方式表述这些事情。
% 讨论应该围绕不变量？保护数据？避免竞争？确保原子性？避免特定危险？

% "原子"表述通常不能真正说明危险是什么。

\begin{itemize}

\item 与 \lstinline{p->state} 一起，它防止在为新进程分配
  \lstinline{proc[]} 槽时的竞争。
% 使对 UNUSED 的检查与分配（或其他）原子化。

\item 它在进程被创建或销毁时将其隐藏起来。
% 使分配的所有步骤和销毁原子化。

\item 它防止父进程的 \lstinline{wait} 收集一个已将其状态设置为 \lstinline{ZOMBIE} 但尚未让出 CPU 的进程。
% 它将设置状态为 ZOMIE 和让出 CPU 原子化。

\item 它防止另一个 CPU 的调度器决定在进程将其状态设置为 \lstinline{RUNNABLE} 之后但在完成 \lstinline{swtch} 之前运行一个让出的进程。
% 它将设置 p->state 和 swtch() 原子化。

\item 它确保只有一个 CPU 的调度器决定运行一个
  \lstinline{RUNNABLE} 进程。
% 它将调度器对 RUNNABLE 的检查与实际运行进程（或设置状态为 RUNNING）原子化。

\item 它防止定时器中断导致进程在 \lstinline{swtch} 中时让出。
% 它使 swtch（或实际上是整个序列）相对于定时器中断原子化

\item 与条件锁一起，它有助于防止 \lstinline{wakeup}
忽略一个正在调用 \lstinline{sleep} 但尚未完成让出 CPU 的进程。
% 避免 conditionCheck+sleep 和 wakeup 之间的竞争？
% 使条件检查与让出原子化？

\item 它防止 \lstinline{kill} 的受害者进程在 \lstinline{kkill} 检查
  \lstinline{p->pid} 和设置 \lstinline{p->killed} 之间退出并可能被重新分配。
% 它使 pid 检查与设置 killed 原子化。

\item 它使 \lstinline{kkill} 对 \lstinline{p->state} 的检查和写入原子化。

\end{itemize}


\lstinline{p->parent} 字段由全局锁
\lstinline{wait_lock} 而不是 \lstinline{p->lock} 保护。
只有进程的父进程修改 \lstinline{p->parent}，
尽管该字段被进程本身和其他寻找其子进程的进程读取。
\lstinline{wait_lock} 的目的是作为条件锁，当
\lstinline{wait} 睡眠等待任何子进程退出时。
退出的子进程在将其状态设置为 \lstinline{ZOMBIE}、唤醒其父进程并让出 CPU 之后持有
\lstinline{wait_lock} 或 \lstinline{p->lock}。
\lstinline{wait_lock} 还序列化父进程和子进程的并发
\lstinline{exit}，
以便 \lstinline{init} 进程（继承子进程）
保证从其 \lstinline{wait} 中被唤醒。
\lstinline{wait_lock} 是全局锁而不是每个父进程的锁，
因为，直到进程获取它之前，它不知道其父进程是谁。

%%
\section{现实世界}
%%

\lstinline{sleep}
和
\lstinline{wakeup}
是一种简单有效的同步方法，
但还有许多其他方法；
信号量~
\cite{dijkstra65}
是一个例子。
所有这些方法的第一个挑战是避免我们在本章开头看到的``丢失唤醒''问题。
原始 Unix 内核的
\lstinline{sleep}
只是禁用中断，
这就足够了，因为 Unix 运行在单 CPU 系统上。
因为 xv6 运行在多处理器上，
它向
\lstinline{sleep}
添加了一个显式锁。
FreeBSD 的
\lstinline{msleep}
采用相同的方法。
Plan 9 的
\lstinline{sleep}
使用一个在即将睡眠之前持有调度锁运行的回调函数；
该函数作为睡眠条件的最后检查，以避免丢失唤醒。
Linux 内核的
\lstinline{sleep}
使用显式的进程队列，称为等待队列，而不是等待通道；该队列有自己的内部锁。

在
\lstinline{wakeup}
中扫描整个进程集是低效的。
更好的解决方案是将
\lstinline{sleep}
和
\lstinline{wakeup}
中的
\lstinline{chan}
替换为一个保存该结构上睡眠进程列表的数据结构，
如 Linux 的等待队列。
Plan 9 的
\lstinline{sleep}
和
\lstinline{wakeup}
将该结构称为会合点（rendezvous point）。
许多线程库将相同的结构称为条件变量（condition variable）；
在这种情况下，操作
\lstinline{sleep}
和
\lstinline{wakeup}
称为
\lstinline{wait}
和
\lstinline{signal}。
所有这些机制都有相同的特点：睡眠条件由某种在睡眠期间原子地释放的锁保护。

xv6 的
\lstinline{wakeup}
唤醒所有在特定等待通道上等待的进程。
如果有多个，它们都将尝试获取条件锁并重新检查条件；
在许多情况下只有一个能够做有用的事情（例如，读取管道中等待的所有数据）。
其余的将发现条件不再为真并返回睡眠；
唤醒它们是浪费 CPU 时间。
因此，
大多数条件变量设计提供两个原语：
\lstinline{signal}，
唤醒等待条件变量的进程之一，和
\lstinline{broadcast}，
唤醒它们所有。

强制终止进程带来一些问题。
例如，被终止的进程可能在内核深处睡眠，
展开其堆栈需要小心，因为调用堆栈上的每个函数可能需要做一些清理。
一些语言通过提供异常机制来帮助，但 C 没有。
此外，还有其他事件可以导致睡眠进程被唤醒，
即使它等待的事件尚未发生。
例如，当 Unix 进程睡眠时，另一个进程可能向它发送
\indexcode{signal}。
在这种情况下，进程将从被中断的系统调用返回，值为 -1 且错误代码设置为 EINTR。
应用程序可以检查这些值并决定做什么。
Xv6 不支持信号，这种复杂性不会出现。

Xv6 对
\lstinline{kill}
的支持不完全令人满意：有些睡眠循环可能应该检查
\lstinline{p->killed}。
一个相关的问题是，即使对于检查
\lstinline{p->killed}
的
\lstinline{sleep}
循环，
\lstinline{sleep}
和
\lstinline{kill}
之间也存在竞争；
后者可能设置
\lstinline{p->killed}
并尝试唤醒受害者，就在受害者的循环检查
\lstinline{p->killed}
之后但在其调用
\lstinline{sleep}
之前。
如果发生这个问题，受害者直到它等待的条件发生才会注意到
\lstinline{p->killed}。
这可能很久之后或永远不会
（例如，如果受害者正在等待控制台的输入，但用户没有键入任何输入）。

%%
\section{练习}
%%

\begin{enumerate}

\item 在 xv6 中实现计数信号量。
选择 xv6 的一些睡眠和唤醒用法并用信号量替换它们。
评判结果。

\item 你能实现一个只带一个参数的 {\tt sleep()} 变体吗，即通道，不需要锁参数？

\item 修复上面提到的
\lstinline{kill}
和
\lstinline{sleep}
之间的竞争，
使得在受害者的睡眠循环检查
\lstinline{p->killed}
之后但在调用
\lstinline{sleep}
之前发生的
\lstinline{kill}
导致受害者放弃当前系统调用。
% 答案：一个解决方案是在设置进程状态为睡眠之前在 sleep 中检查 p->killed 是否已设置。

\item 设计一个计划，使每个睡眠循环都检查
\lstinline{p->killed}
，以便例如，在 virtio 驱动程序中的进程如果被另一个进程终止，可以从 while 循环快速返回。
% 答案：这很困难。现代 Unix 使用 setjmp 和 longjmp 以及非常仔细的编程来清理被中断系统调用可能已建立的任何部分状态。

\end{enumerate}

  \caption{重叠锁以避免丢失唤醒}
    \label{fig:overlap}
\end{figure}

\lstinline{sleep}
和
\lstinline{wakeup}
的锁定规则为什么能确保即将睡眠的进程不会错过并发的唤醒？
即将睡眠的进程在检查条件之前直到标记自己为 \texttt{SLEEPING} 之后，持有
条件锁或它自己的
\lstinline{p->lock}
或两者；
参见图~\ref{fig:overlap}。
调用 \texttt{wakeup} 的进程需要获取{\it 两个}锁。
唤醒者可能首先获取锁，这意味着它将在消费线程检查条件之前使条件为真，
而且消费线程不需要调用 {\tt sleep()}；
或者唤醒者的 {\tt acquire()} 可能需要等待，直到消费线程完全完成睡眠并
释放锁，在这种情况下唤醒者将看到消费线程被标记为 {\tt SLEEPING}
并将唤醒它。

有时多个进程在同一个通道上睡眠；
例如，多个进程从管道读取。
对
\lstinline{wakeup}
的单个调用将唤醒它们所有。
其中一个将首先运行并获取 \lstinline{sleep}
被调用时使用的锁，并（在管道的情况下）读取等待的数据。
其他进程将发现，尽管被唤醒了，
但没有数据可读取。
从它们的角度来看，唤醒是``虚假的（spurious）''，
它们必须再次睡眠。
因此 \lstinline{sleep} 总是在重新检查条件的循环内调用，如上面的 \lstinline{P} 所示。

如果睡眠/唤醒 的两个使用意外选择了相同的通道：
它们会看到虚假的唤醒，
但如上所述的循环将容忍这个问题。
睡眠/唤醒 的大部分魅力在于它既轻量级（不需要创建特殊的数据结构作为等待通道）
又提供了一层间接（调用者不需要知道它们正在与哪个特定进程交互）。
%%
\section{代码：管道}
%%
xv6 的管道是使用
\lstinline{sleep}
和 \lstinline{wakeup} 来同步生产者和消费者的代码示例。
我们在第~\ref{CH:UNIX} 章看到了管道的接口：
写入管道一端的数据被复制到内核缓冲区，
然后可以从管道的另一端读取。
后续章节将检查围绕管道的文件描述符支持，
但现在让我们看看
\indexcode{pipewrite}
和
\indexcode{piperead}
的实现。

每个管道由
\indexcode{struct pipe}
表示，它包含一个
\lstinline{lock}
和一个
\lstinline{data}
缓冲区。
字段
\lstinline{nread}
和
\lstinline{nwrite}
分别计数从缓冲区读取和写入的字节总数。
缓冲区环绕：
在
\lstinline{buf[PIPESIZE-1]}
之后写入的下一个字节是
\lstinline{buf[0]}。
计数不环绕。
这种约定让实现能够区分满缓冲区
(\lstinline{nwrite}
\lstinline{==}
\lstinline{nread+PIPESIZE})
和空缓冲区
(\lstinline{nwrite}
\lstinline{==}
\lstinline{nread})，
但它意味着索引到缓冲区时必须使用
\lstinline{buf[nread}
\lstinline{%}
\lstinline{PIPESIZE]}
而不是仅仅
\lstinline{buf[nread]}
（对
\lstinline{nwrite}
同样如此）。

假设对
\lstinline{piperead}
和
\lstinline{pipewrite}
的调用同时在两个不同的 CPU 上发生。
\lstinline{pipewrite}
\lineref{kernel/pipe.c:/^pipewrite/}
首先获取管道的锁，该锁保护计数、数据及其相关的
不变量。
\lstinline{piperead}
\lineref{kernel/pipe.c:/^piperead/}
然后也尝试获取锁，但无法获取。
它在
\lstinline{acquire}
\lineref{kernel/spinlock.c:/^acquire/}
中自旋等待锁。
在
\lstinline{piperead}
等待时，
\lstinline{pipewrite}
循环遍历要写入的字节
(\lstinline{addr[0..n-1]})，
依次将每个字节添加到管道中
\lineref{kernel/pipe.c:/nwrite\+\+/}。
在这个循环期间，缓冲区可能填满
\lineref{kernel/pipe.c:/DOC: pipewrite-full/}。
在这种情况下，
\lstinline{pipewrite}
调用
\lstinline{wakeup}
来提醒任何睡眠的读取者缓冲区中有数据等待，
然后在
\lstinline{&pi->nwrite}
上睡眠以等待读取者从缓冲区中取出一些字节。
\lstinline{sleep}
释放管道锁作为将
\lstinline{pipewrite}
的进程置于睡眠的一部分。

\lstinline{piperead}
现在获取管道锁并进入其临界区：
它发现
\lstinline{pi->nread}
\lstinline{!=}
\lstinline{pi->nwrite}
\lineref{kernel/pipe.c:/DOC: pipe-empty/}
(\lstinline{pipewrite}
进入睡眠是因为
\lstinline{pi->nwrite}
\lstinline{==}
\lstinline{pi->nread}
\lstinline{+}
\lstinline{PIPESIZE}
\lineref{kernel/pipe.c:/pipewrite-full/})，
因此它进入
\lstinline{for}
循环，从管道中复制数据
\lineref{kernel/pipe.c:/DOC: piperead-copy/}，
并将
\lstinline{nread}
增加复制的字节数。
缓冲区中的那么多空间现在可用于写入，因此
\lstinline{piperead}
调用
\lstinline{wakeup}
\lineref{kernel/pipe.c:/DOC: piperead-wakeup/}
来唤醒任何睡眠的写入者，然后返回。
\lstinline{wakeup}
找到一个在
\lstinline{&pi->nwrite}
上睡眠的进程，
该进程正在运行
\lstinline{pipewrite}
但在缓冲区填满时停止了。
它将该进程标记为
\indexcode{RUNNABLE}。

管道代码为读取者和写入者使用单独的等待通道
(\lstinline{pi->nread}
和
\lstinline{pi->nwrite})；
这在不太可能有大量读取者和写入者等待同一管道的情况下可能使系统更高效。
管道代码在检查睡眠条件的循环内睡眠；
如果有多个读取者或写入者，除第一个唤醒的进程外，所有其他进程将看到条件为假并再次睡眠。
%%
\section{代码：Wait、exit 和 kill}

请阅读 {\tt kernel/proc.c} 中函数 {\tt kwait()}、{\tt kexit()}
和 {\tt kkill()} 的代码；这些是对应系统调用的内部实现。

%%
\lstinline{sleep}
和
\lstinline{wakeup}
可用于许多类型的等待。
第~\ref{CH:UNIX} 章介绍了一个有趣的例子，
即子进程 \indexcode{exit}
与其父进程 \indexcode{wait} 之间的交互。
在子进程死亡时，父进程可能已经在 {\tt wait} 中睡眠，
或者可能在做其他事情；
在后一种情况下，后续对 {\tt wait} 的调用必须观察子进程的死亡，
可能在调用 {\tt exit} 很久之后。
Xv6 记录子进程死亡直到 {\tt wait} 观察到它的方式是 {\tt exit} 将调用者放入 \indexcode{ZOMBIE}
状态，在那里保持直到父进程的 {\tt wait} 注意到它，将子进程的状态更改为 {\tt UNUSED}，
复制子进程的退出状态，并将子进程的进程 ID 返回给父进程。
如果父进程在子进程之前退出，
父进程将子进程交给
\lstinline{init}
进程，该进程永久调用 {\tt wait}；
因此每个子进程都有一个父进程来清理它。
一个挑战是避免同时父进程
\lstinline{wait}
和子进程
\lstinline{exit}
之间的竞争和死锁，
以及同时
\lstinline{exit} 和 \lstinline{exit} 之间的竞争。

\lstinline{kwait}，{\tt wait} 的内核实现，首先获取
\lstinline{wait_lock}
\lineref{kernel/proc.c:/^kwait/}，
它作为条件锁有助于确保 \lstinline{kwait} 不会错过来自退出子进程的 \lstinline{wakeup}。
然后 \lstinline{kwait} 扫描进程表。
如果它找到处于 \texttt{ZOMBIE} 状态的子进程，
它释放该子进程的资源和其 \lstinline{proc} 结构体，
将子进程的退出状态复制到提供给 \lstinline{wait} 的地址
（如果不是 0），
并返回子进程的进程 ID。
如果
\lstinline{kwait}
找到子进程但没有一个退出，
它调用
\lstinline{sleep}
等待它们中的任何一个退出
\lineref{kernel/proc.c:/DOC: wait-sleep/}，
然后再次扫描。
\lstinline{kwait} 经常持有两个锁，
\lstinline{wait_lock} 和某个进程的 \lstinline{pp->lock}；
避免死锁的顺序是先 \lstinline{wait_lock} 然后 \lstinline{pp->lock}。

\lstinline{kexit} \lineref{kernel/proc.c:/^kexit/} 记录退出状态，
释放一些资源，调用 \lstinline{reparent} 将任何子进程交给 \lstinline{init} 进程，
唤醒父进程以防它在 \lstinline{wait} 中睡眠，将调用者标记为僵尸，并永久让出 CPU。
\lstinline{kexit} 在此过程中持有
\lstinline{wait_lock} 和 \lstinline{p->lock}。
它持有 \lstinline{wait_lock} 因为它是
\lstinline{wakeup(p->parent)} 的条件锁，防止在 \lstinline{wait} 中的父进程丢失唤醒。
\lstinline{kexit} 还必须为此序列持有
\lstinline{p->lock}，以防止在 \lstinline{wait} 中的父进程在子进程最终调用
\lstinline{swtch} 之前看到子进程处于状态
\lstinline{ZOMBIE}。
\lstinline{kexit} 以与 \lstinline{kwait} 相同的顺序获取这些锁以避免死锁。

\lstinline{kexit} 在将其状态设置为 \lstinline{ZOMBIE} 之前唤醒父进程
可能看起来不正确，
但这是安全的：
尽管
\lstinline{wakeup}
可能导致父进程运行，
但父进程的
\lstinline{kwait}
中的循环在子进程的
\lstinline{p->lock}
被 {\tt scheduler} 释放之前无法检查子进程，
因此
\lstinline{kwait}
不能在
\lstinline{kexit}
将其状态设置为
\lstinline{ZOMBIE}
之前查看退出进程
\lineref{kernel/proc.c:/state.=.ZOMBIE/}。

虽然
\lstinline{exit}
允许进程终止自己，
但 \lstinline{kill} 系统调用
\lineref{kernel/proc.c:/^kkill/}
让一个进程请求另一个进程终止。
\lstinline{kill}
直接销毁受害者进程会太复杂，因为受害者
可能在另一个 CPU 上执行，
可能正处于更新内核数据结构的关键序列的中间。
因此
\lstinline{kkill}
做得很少：它只是设置受害者的
\indexcode{p->killed}
并且，如果它正在睡眠，就唤醒它。
最终受害者将进入或离开内核，
此时
\lstinline{usertrap}
中的代码将调用
\lstinline{kexit}
如果
\lstinline{p->killed}
被设置
（它通过调用
\lstinline{killed}
\lineref{kernel/proc.c:/^killed/}
来检查）。
如果受害者在用户空间运行，它将在下次进入内核时看到它已被终止，
通过进行系统调用或因为定时器（或其他设备）中断。

如果受害者进程在
\lstinline{sleep}
中，
\lstinline{kkill} 对
\lstinline{wakeup}
的调用将导致受害者从
\lstinline{sleep}
返回。
这有潜在危险，因为正在等待的条件可能不为真。
然而，xv6 对
\lstinline{sleep}
的调用总是包装在一个在
\lstinline{sleep}
返回后重新测试条件的
\lstinline{while}
循环中。
一些对
\lstinline{sleep}
的调用也在循环中测试
\lstinline{p->killed}，
如果设置则放弃当前活动。
这仅在放弃是正确的情况下才会这样做。
例如，管道读取和写入代码
\lineref{kernel/pipe.c:/killed.pr/}
如果设置了终止标志则返回；最终代码将返回到 trap，trap 将再次检查 \lstinline{p->killed} 并退出。

一些 xv6
\lstinline{sleep}
循环不检查
\lstinline{p->killed}，
因为代码处于多步系统调用的中间，应该是原子的（即，如果中途放弃会不正确）。
Virtio 驱动程序
\lineref{kernel/virtio\_disk.c:/sleep.b/}
是一个例子：它不检查
\lstinline{p->killed}
因为磁盘操作可能是一组写入中的一部分，所有这些写入都需要才能使文件系统保持正确状态。
在磁盘 I/O 上等待时被终止的进程直到完成当前系统调用并且
\lstinline{usertrap} 看到终止标志才会退出。

\section{进程锁}

与每个进程关联的锁（\lstinline{p->lock}）是 xv6 中最复杂的锁。
理解 \lstinline{p->lock} 的一个简单方法是，
在读取或写入以下任何 \lstinline{struct proc} 字段时必须持有它：
\lstinline{p->state}、
\lstinline{p->chan}、
\lstinline{p->killed}、
\lstinline{p->xstate}、
和
\lstinline{p->pid}。
这些字段可以被其他进程使用，或被其他 CPU 上的调度器线程使用，
因此它们自然必须由锁保护。

然而，\lstinline{p->lock} 的大多数使用是保护 xv6 进程数据结构和算法的更高级别不变量。
以下是 \lstinline{p->lock} 所做的全部事情：

% 很难知道如何以统一的方式表述这些事情。
% 讨论应该围绕不变量？保护数据？避免竞争？确保原子性？避免特定危险？

% "原子"表述通常不能真正说明危险是什么。

\begin{itemize}

\item 与 \lstinline{p->state} 一起，它防止在为新进程分配
  \lstinline{proc[]} 槽时的竞争。
% 使对 UNUSED 的检查与分配（或其他）原子化。

\item 它在进程被创建或销毁时将其隐藏起来。
% 使分配的所有步骤和销毁原子化。

\item 它防止父进程的 \lstinline{wait} 收集一个已将其状态设置为 \lstinline{ZOMBIE} 但尚未让出 CPU 的进程。
% 它将设置状态为 ZOMIE 和让出 CPU 原子化。

\item 它防止另一个 CPU 的调度器决定在进程将其状态设置为 \lstinline{RUNNABLE} 之后但在完成 \lstinline{swtch} 之前运行一个让出的进程。
% 它将设置 p->state 和 swtch() 原子化。

\item 它确保只有一个 CPU 的调度器决定运行一个
  \lstinline{RUNNABLE} 进程。
% 它将调度器对 RUNNABLE 的检查与实际运行进程（或设置状态为 RUNNING）原子化。

\item 它防止定时器中断导致进程在 \lstinline{swtch} 中时让出。
% 它使 swtch（或实际上是整个序列）相对于定时器中断原子化

\item 与条件锁一起，它有助于防止 \lstinline{wakeup}
忽略一个正在调用 \lstinline{sleep} 但尚未完成让出 CPU 的进程。
% 避免 conditionCheck+sleep 和 wakeup 之间的竞争？
% 使条件检查与让出原子化？

\item 它防止 \lstinline{kill} 的受害者进程在 \lstinline{kkill} 检查
  \lstinline{p->pid} 和设置 \lstinline{p->killed} 之间退出并可能被重新分配。
% 它使 pid 检查与设置 killed 原子化。

\item 它使 \lstinline{kkill} 对 \lstinline{p->state} 的检查和写入原子化。

\end{itemize}


\lstinline{p->parent} 字段由全局锁
\lstinline{wait_lock} 而不是 \lstinline{p->lock} 保护。
只有进程的父进程修改 \lstinline{p->parent}，
尽管该字段被进程本身和其他寻找其子进程的进程读取。
\lstinline{wait_lock} 的目的是作为条件锁，当
\lstinline{wait} 睡眠等待任何子进程退出时。
退出的子进程在将其状态设置为 \lstinline{ZOMBIE}、唤醒其父进程并让出 CPU 之后持有
\lstinline{wait_lock} 或 \lstinline{p->lock}。
\lstinline{wait_lock} 还序列化父进程和子进程的并发
\lstinline{exit}，
以便 \lstinline{init} 进程（继承子进程）
保证从其 \lstinline{wait} 中被唤醒。
\lstinline{wait_lock} 是全局锁而不是每个父进程的锁，
因为，直到进程获取它之前，它不知道其父进程是谁。

%%
\section{现实世界}
%%

\lstinline{sleep}
和
\lstinline{wakeup}
是一种简单有效的同步方法，
但还有许多其他方法；
信号量~
\cite{dijkstra65}
是一个例子。
所有这些方法的第一个挑战是避免我们在本章开头看到的``丢失唤醒''问题。
原始 Unix 内核的
\lstinline{sleep}
只是禁用中断，
这就足够了，因为 Unix 运行在单 CPU 系统上。
因为 xv6 运行在多处理器上，
它向
\lstinline{sleep}
添加了一个显式锁。
FreeBSD 的
\lstinline{msleep}
采用相同的方法。
Plan 9 的
\lstinline{sleep}
使用一个在即将睡眠之前持有调度锁运行的回调函数；
该函数作为睡眠条件的最后检查，以避免丢失唤醒。
Linux 内核的
\lstinline{sleep}
使用显式的进程队列，称为等待队列，而不是等待通道；该队列有自己的内部锁。

在
\lstinline{wakeup}
中扫描整个进程集是低效的。
更好的解决方案是将
\lstinline{sleep}
和
\lstinline{wakeup}
中的
\lstinline{chan}
替换为一个保存该结构上睡眠进程列表的数据结构，
如 Linux 的等待队列。
Plan 9 的
\lstinline{sleep}
和
\lstinline{wakeup}
将该结构称为会合点（rendezvous point）。
许多线程库将相同的结构称为条件变量（condition variable）；
在这种情况下，操作
\lstinline{sleep}
和
\lstinline{wakeup}
称为
\lstinline{wait}
和
\lstinline{signal}。
所有这些机制都有相同的特点：睡眠条件由某种在睡眠期间原子地释放的锁保护。

xv6 的
\lstinline{wakeup}
唤醒所有在特定等待通道上等待的进程。
如果有多个，它们都将尝试获取条件锁并重新检查条件；
在许多情况下只有一个能够做有用的事情（例如，读取管道中等待的所有数据）。
其余的将发现条件不再为真并返回睡眠；
唤醒它们是浪费 CPU 时间。
因此，
大多数条件变量设计提供两个原语：
\lstinline{signal}，
唤醒等待条件变量的进程之一，和
\lstinline{broadcast}，
唤醒它们所有。

强制终止进程带来一些问题。
例如，被终止的进程可能在内核深处睡眠，
展开其堆栈需要小心，因为调用堆栈上的每个函数可能需要做一些清理。
一些语言通过提供异常机制来帮助，但 C 没有。
此外，还有其他事件可以导致睡眠进程被唤醒，
即使它等待的事件尚未发生。
例如，当 Unix 进程睡眠时，另一个进程可能向它发送
\indexcode{signal}。
在这种情况下，进程将从被中断的系统调用返回，值为 -1 且错误代码设置为 EINTR。
应用程序可以检查这些值并决定做什么。
Xv6 不支持信号，这种复杂性不会出现。

Xv6 对
\lstinline{kill}
的支持不完全令人满意：有些睡眠循环可能应该检查
\lstinline{p->killed}。
一个相关的问题是，即使对于检查
\lstinline{p->killed}
的
\lstinline{sleep}
循环，
\lstinline{sleep}
和
\lstinline{kill}
之间也存在竞争；
后者可能设置
\lstinline{p->killed}
并尝试唤醒受害者，就在受害者的循环检查
\lstinline{p->killed}
之后但在其调用
\lstinline{sleep}
之前。
如果发生这个问题，受害者直到它等待的条件发生才会注意到
\lstinline{p->killed}。
这可能很久之后或永远不会
（例如，如果受害者正在等待控制台的输入，但用户没有键入任何输入）。

%%
\section{练习}
%%

\begin{enumerate}

\item 在 xv6 中实现计数信号量。
选择 xv6 的一些睡眠和唤醒用法并用信号量替换它们。
评判结果。

\item 你能实现一个只带一个参数的 {\tt sleep()} 变体吗，即通道，不需要锁参数？

\item 修复上面提到的
\lstinline{kill}
和
\lstinline{sleep}
之间的竞争，
使得在受害者的睡眠循环检查
\lstinline{p->killed}
之后但在调用
\lstinline{sleep}
之前发生的
\lstinline{kill}
导致受害者放弃当前系统调用。
% 答案：一个解决方案是在设置进程状态为睡眠之前在 sleep 中检查 p->killed 是否已设置。

\item 设计一个计划，使每个睡眠循环都检查
\lstinline{p->killed}
，以便例如，在 virtio 驱动程序中的进程如果被另一个进程终止，可以从 while 循环快速返回。
% 答案：这很困难。现代 Unix 使用 setjmp 和 longjmp 以及非常仔细的编程来清理被中断系统调用可能已建立的任何部分状态。

\end{enumerate}

\chapter{文件系统}
\label{CH:FS}
%

文件系统的目的是组织和存储数据。文件系统通常支持用户和应用程序之间的数据共享，以及
\indextext{持久性（persistence）}，
以便数据在重启后仍然可用。

xv6 文件系统提供类似 Unix 的文件、目录和路径名
（参见第~\ref{CH:UNIX} 章），并将其数据存储在 virtio 磁盘上以实现持久性。文件系统解决了几个挑战：
\begin{itemize}

\item 文件系统需要磁盘上的数据结构来表示命名的目录和文件树，
记录保存每个文件内容的块的标识，并记录磁盘的哪些区域是空闲的。
\item 文件系统必须支持
\indextext{崩溃恢复（crash recovery）}。
也就是说，如果发生崩溃（例如，电源故障），文件系统必须在重启后仍然正常工作。
风险在于崩溃可能会中断一系列更新，导致磁盘上的数据结构不一致（例如，一个块既在文件中使用又被标记为空闲）。
\item 不同的进程可能同时对文件系统进行操作，
因此文件系统代码必须协调以保持不变量。
\item 访问磁盘比访问内存慢几个数量级，
因此文件系统必须维护一个热门块的内存缓存。

\end{itemize}

本章的其余部分解释了 xv6 如何解决这些挑战。
%%
%%  -------------------------------------------
%%
\section{概述}

xv6 文件系统实现分为七个层次，如图~\ref{fig:fslayer} 所示。
磁盘层在 virtio 硬盘上读写块。
缓冲区缓存层缓存磁盘块并同步对它们的访问，
确保一次只有一个内核进程可以修改存储在任何特定块中的数据。
日志层允许更高层将对多个块的更新包装在一个
\indextext{事务（transaction）}中，
并确保在崩溃面前原子地更新块（即，全部更新或都不更新）。
索引节点层提供单个文件，每个文件由一个具有唯一 i-number 的
\indextext{索引节点（inode）}
和一些保存文件数据的块表示。
目录层将每个目录实现为一种特殊的 inode，其内容是一系列目录项，每个目录项包含文件的名称和 i-number。
路径名层提供分层路径名，如
\lstinline{/usr/rtm/xv6/fs.c}，
并通过递归查找来解析它们。
文件描述符层使用文件系统接口抽象许多 Unix 资源（例如，管道、设备、文件等），简化应用程序程序员的工作。

\begin{figure}[t]
\centering
\includegraphics[scale=0.5]{fig/fslayer.pdf}
\caption{xv6 文件系统的层次。}
\label{fig:fslayer}
\end{figure}

磁盘硬件传统上将磁盘上的数据表示为编号序列的 512 字节
\textit{块}
\index{块（block）}
（也称为
\textit{扇区}）：
\index{扇区（sector）}
扇区 0 是前 512 字节，扇区 1 是接下来的 512 字节，依此类推。操作系统用于其文件系统的块大小可能与磁盘使用的扇区大小不同，但通常块大小是扇区大小的倍数。Xv6 将已读入内存的块副本保存在类型为
\lstinline{struct buf}
\lineref{kernel/buf.h:/^struct.buf/}
的对象中。
该结构中存储的数据有时与磁盘不同步：它可能尚未从磁盘读入（磁盘正在处理但尚未返回扇区的内容），或者它可能已被软件更新但尚未写入磁盘。

文件系统必须有一个计划来确定它在磁盘上存储 inode 和内容块的位置。
为此，xv6 将磁盘分成几个部分，如图~\ref{fig:fslayout} 所示。
文件系统不使用块 0（它保存引导扇区）。
块 1 称为
\indextext{超级块（superblock）}；
它包含文件系统的元数据（文件系统大小（以块为单位）、数据块数量、inode 数量和日志中的块数）。
从块 2 开始的块保存日志。日志之后是 inode，每个块包含多个 inode。
之后是位图块，跟踪哪些数据块正在使用。
剩余的块是数据块；每个块要么在位图块中标记为空闲，要么保存文件或目录的内容。
超级块由一个称为
\indexcode{mkfs}
的单独程序填充，该程序构建初始文件系统。

本章的其余部分讨论每一层，从缓冲区缓存开始。
注意低层中精心选择的抽象如何简化高层设计的情况。
%%
%%  -------------------------------------------
%%
\section{缓冲区缓存层}
\label{s:bcache}

缓冲区缓存有两个工作：(1) 同步对磁盘块的访问，以确保内存中只有一个块的副本，并且一次只有一个内核线程使用该副本；
(2) 缓存热门块，以便不需要从慢速磁盘重新读取。
代码在
\lstinline{bio.c} 中。

缓冲区缓存导出的主要接口包括
\indexcode{bread}
和
\indexcode{bwrite}；
前者获取一个包含块副本的
\indextext{缓冲区（buf）}，可以在内存中读取或修改，后者将修改后的缓冲区写入磁盘上的适当块。
内核线程在使用完缓冲区后必须通过调用
\indexcode{brelse}
释放它。
缓冲区缓存使用每个缓冲区的睡眠锁来确保一次只有一个线程使用每个缓冲区
（因此每个磁盘块）；
\lstinline{bread}
返回一个锁定的缓冲区，
\lstinline{brelse}
释放锁。

让我们回到缓冲区缓存。
缓冲区缓存有固定数量的缓冲区来保存磁盘块，
这意味着如果文件系统请求一个尚未在缓存中的块，缓冲区缓存必须回收当前保存其他块的缓冲区。
缓冲区缓存为新块回收最近最少使用的缓冲区。
假设是最近最少使用的缓冲区是最不可能很快再次使用的缓冲区。

\begin{figure}[t]
\centering
\includegraphics[scale=0.5]{fig/fslayout.pdf}
\caption{xv6 文件系统的结构。}
\label{fig:fslayout}
\end{figure}
%%
%%  -------------------------------------------
%%
\section{代码：缓冲区缓存}

缓冲区缓存是缓冲区的双向链表。
函数
\indexcode{binit}，
由
\indexcode{main}
\lineref{kernel/main.c:/binit/}
调用，
用静态数组
\lstinline{buf}
中的
\indexcode{NBUF}
个缓冲区初始化列表
\linerefs{kernel/bio.c:/Create.linked.list/,/^..}/}。
所有其他对缓冲区缓存的访问都通过
\indexcode{bcache.head}
引用链表，而不是
\lstinline{buf}
数组。

缓冲区有两个与之关联的状态字段。
字段
\indexcode{valid}
表示缓冲区包含块的副本。
字段 \indexcode{disk}
表示缓冲区内容已交给磁盘，磁盘可能会更改缓冲区（例如，将数据从磁盘写入 \lstinline{data}）。

\lstinline{bread}
\lineref{kernel/bio.c:/^bread/}
调用
\indexcode{bget}
获取给定扇区的缓冲区
\lineref{kernel/bio.c:/b.=.bget/}。
如果需要从磁盘读取缓冲区，
\lstinline{bread}
调用
\indexcode{virtio_disk_rw}
在返回缓冲区之前执行此操作。

\lstinline{bget}
\lineref{kernel/bio.c:/^bget/}
扫描缓冲区列表以查找具有给定设备和扇区号的缓冲区
\linerefs{kernel/bio.c:/Is.the.block.already/,/^..}/}。
如果有这样的缓冲区，
\indexcode{bget}
获取缓冲区的睡眠锁。
\lstinline{bget}
然后返回锁定的缓冲区。

如果没有给定扇区的缓存缓冲区，
\indexcode{bget}
必须创建一个，可能重用保存不同扇区的缓冲区。
它第二次扫描缓冲区列表，查找未使用的缓冲区（\lstinline{b->refcnt = 0}）；
任何这样的缓冲区都可以使用。
\lstinline{bget}
编辑缓冲区元数据以记录新的设备和扇区号，并获取其睡眠锁。
注意赋值
\lstinline{b->valid = 0}
确保
\lstinline{bread}
将从磁盘读取块数据，
而不是错误地使用缓冲区的先前内容。

每个磁盘扇区最多有一个缓存缓冲区很重要，以确保读取者看到写入，
而且因为文件系统使用缓冲区上的锁进行同步。
\lstinline{bget}
通过从第一个循环检查块是否已缓存到第二个循环声明块现在已缓存（通过设置
\lstinline{dev}、
\lstinline{blockno}
和
\lstinline{refcnt}）持续持有
\lstinline{bcache.lock}
来确保这个不变量。
这使得检查块是否存在以及（如果不存在）指定一个缓冲区来保存块是原子的。

\lstinline{bget}
在
\lstinline{bcache.lock}
临界区外获取缓冲区的睡眠锁是安全的，
因为非零的
\lstinline{b->refcnt}
阻止缓冲区被重用于不同的磁盘块。
睡眠锁保护块缓冲内容的读写，而
\lstinline{bcache.lock}
保护关于哪些块被缓存的信息。

如果所有缓冲区都忙，那么同时执行文件系统调用的进程太多了；
\lstinline{bget}
发生 panic。
更优雅的响应可能是睡眠直到缓冲区空闲，
尽管那样会有死锁的可能性。

一旦
\indexcode{bread}
读取了磁盘（如果需要）并将缓冲区返回给其调用者，调用者就拥有缓冲区的独占使用权，可以读取或写入数据字节。
如果调用者确实修改了缓冲区，它必须在释放缓冲区之前调用
\indexcode{bwrite}
将更改的数据写入磁盘。
\lstinline{bwrite}
\lineref{kernel/bio.c:/^bwrite/}
调用
\indexcode{virtio_disk_rw}
与磁盘硬件通信。

当调用者使用完缓冲区时，它必须调用
\indexcode{brelse}
释放它。
（名称
\lstinline{brelse}，
b-release 的缩写，
是神秘的但值得学习的：
它起源于 Unix，也用于 BSD、Linux 和 Solaris。）
\lstinline{brelse}
\lineref{kernel/bio.c:/^brelse/}
释放睡眠锁并将缓冲区移到链表的前端
\linerefs{kernel/bio.c:/b->next->prev.=.b->prev/,/bcache.head.next.=.b/}。
移动缓冲区使列表按缓冲区最近使用的时间（意味着释放）排序：
列表中的第一个缓冲区是最近使用的，
最后一个是最近最少使用的。
\lstinline{bget}
中的两个循环利用这一点：
扫描现有缓冲区的最坏情况下必须处理整个列表，
但先检查最近使用的缓冲区（从
\lstinline{bcache.head}
开始并跟随
\lstinline{next}
指针）将在引用局部性良好时减少扫描时间。
选择要重用的缓冲区的扫描通过后向扫描（跟随
\lstinline{prev}
指针）选择最近最少使用的缓冲区。
%%
%%  -------------------------------------------
%%
\section{日志层}

文件系统设计中最有趣的问题之一是崩溃恢复。
问题出现是因为许多文件系统操作涉及对磁盘的多次写入，而在部分写入后发生崩溃可能会使磁盘上的文件系统处于不一致状态。
例如，假设在文件截断期间发生崩溃（将文件长度设置为零并释放其内容块）。
根据磁盘写入的顺序，崩溃可能留下一个引用标记为空闲的内容块的 inode，
或者可能留下一个已分配但未引用的内容块。

后者相对无害，但引用已释放块的 inode 可能在重启后导致严重问题。
重启后，内核可能将该块分配给另一个文件，现在我们有两个不同的文件无意中指向同一个块。
如果 xv6 支持多用户，这种情况可能是安全问题，因为旧文件的所有者将能够读写新文件中的块，而新文件由不同的用户拥有。

Xv6 使用一种简单的日志形式来解决文件系统操作期间崩溃的问题。
xv6 系统调用不直接写入磁盘上的文件系统数据结构。
相反，它将它希望进行的所有磁盘写入的描述放在磁盘上的
\indextext{日志（log）}
中。
一旦系统调用记录了其所有写入，它就向磁盘写入一个特殊的
\indextext{提交（commit）}
记录，指示日志包含一个完整的操作。
此时，系统调用将写入复制到磁盘上的文件系统数据结构。
在这些写入完成后，系统调用擦除磁盘上的日志。

如果系统崩溃并重启，文件系统代码在运行任何进程之前按如下方式从崩溃中恢复。
如果日志被标记为包含完整操作，则恢复代码将写入复制到磁盘上文件系统中的适当位置。
如果日志未被标记为包含完整操作，恢复代码将忽略日志。
恢复代码通过擦除日志完成恢复。

为什么 xv6 的日志能解决文件系统操作期间崩溃的问题？
如果崩溃发生在操作提交之前，则磁盘上的日志不会被标记为完整，恢复代码将忽略它，磁盘的状态就像操作甚至没有开始一样。
如果崩溃发生在操作提交之后，则恢复将重放操作的所有写入，如果操作已经开始将它们写入磁盘数据结构，则可能重复它们。
无论哪种情况，日志都使操作相对于崩溃是原子的：恢复后，操作的所有写入要么都出现在磁盘上，要么都不出现。
%%
%%
%%
\section{日志设计}

日志位于已知的固定位置，在超级块中指定。
它由一个头块和一系列更新的块副本（``日志块''）组成。
头块包含一个扇区号数组，每个日志块一个，以及日志块的数量。
磁盘上头块中的计数要么为零，表示日志中没有事务，
要么为非零，表示日志包含具有指示数量的日志块的完整提交事务。
Xv6 在事务提交时写入头块，但之前不写入，并在将日志块复制到文件系统后将计数设置为零。
因此，事务中途崩溃将导致日志头块中的计数为零；
提交后崩溃将导致计数为非零。

每个系统调用的代码指示必须相对于崩溃原子的写入序列的开始和结束。
为了允许不同进程并发执行文件系统操作，
日志系统可以将多个系统调用的写入累积到一个事务中。
因此，单个提交可能涉及多个完整系统调用的写入。
为了避免将系统调用跨事务分割，日志系统仅在没有文件系统系统调用正在进行时才提交。

将多个事务一起提交的想法称为
\indextext{组提交（group commit）}。
组提交减少了磁盘操作的数量，因为它将提交的固定成本分摊到多个操作上。
组提交还同时向磁盘系统提交更多并发写入，
可能允许磁盘在单次磁盘旋转期间写入它们所有。
Xv6 的 virtio 驱动程序不支持这种
\indextext{批处理（batching）}，
但 xv6 的文件系统设计允许这样做。

Xv6 在磁盘上分配固定量的空间来保存日志。
事务中系统调用写入的块总数必须适合该空间。
这有两个后果。
不能允许单个系统调用写入比日志空间更多的不同块。
这对大多数系统调用来说不是问题，但有两个可能写入许多块：
\indexcode{write}
和
\indexcode{unlink}。
大文件写入可能写入许多数据块和许多位图块以及一个 inode 块；
取消链接大文件可能写入许多位图块和一个 inode。
Xv6 的写入系统调用将大写入拆分为多个适合日志的较小写入，
而
\lstinline{unlink}
不会引起问题，因为实际上 xv6 文件系统只使用一个位图块。
有限日志空间的另一个后果是日志系统不能允许系统调用开始，除非它确定系统调用的写入将适合日志中剩余的空间。
%%
%%
%%
\section{代码：日志}

系统调用中日志的典型使用如下所示：
\begin{lstlisting}[]
  begin_op();
  ...
  bp = bread(...);
  bp->data[...] = ...;
  log_write(bp);
  ...
  end_op();
\end{lstlisting}

\indexcode{begin_op}
\lineref{kernel/log.c:/^begin.op/}
等待直到日志系统当前未提交，并且有足够的未预留日志空间来容纳此调用的写入。
\lstinline{log.outstanding}
计数已预留日志空间的系统调用数量；
总预留空间是
\lstinline{log.outstanding}
乘以
\lstinline{MAXOPBLOCKS}。
递增
\lstinline{log.outstanding}
既预留空间又防止在此系统调用期间发生提交。
代码保守地假设每个系统调用可能写入多达
\lstinline{MAXOPBLOCKS}
个不同的块。

\indexcode{log_write}
\lineref{kernel/log.c:/^log.write/}
充当
\indexcode{bwrite}
的代理。
它在内存中记录块的扇区号，
在磁盘上的日志中为它预留一个槽位，
并将缓冲区固定在块缓存中以防止块缓存逐出它。
该块必须保持在缓存中直到提交：
在此之前，缓存的副本是修改的唯一记录；
它不能在提交后写入其在磁盘上的位置；
而且同一事务中的其他读取必须看到修改。
\lstinline{log_write}
注意到当一个块在单个事务中被多次写入时，并为该块在日志中分配相同的槽位。
这种优化通常称为
\indextext{吸收（absorption）}。
常见的情况是，例如，包含多个文件的 inode 的磁盘块在事务内被多次写入。
通过将多个磁盘写入吸收到一个中，文件系统可以节省日志空间，并且可以实现更好的性能，因为只有一个磁盘块副本必须写入磁盘。

\indexcode{end_op}
\lineref{kernel/log.c:/^end.op/}
首先递减未完成系统调用的计数。
如果计数现在为零，它通过调用
\lstinline{commit()}
提交当前事务。
此过程有四个阶段。
\lstinline{write_log()}
\lineref{kernel/log.c:/^write.log/}
将事务中修改的每个块从缓冲区缓存复制到磁盘上日志中的其槽位。
\lstinline{write_head()}
\lineref{kernel/log.c:/^write.head/}
将头块写入磁盘：这是提交点，写入后崩溃将导致恢复从事务的日志重放写入。
\indexcode{install_trans}
\lineref{kernel/log.c:/^install_trans/}
从日志中读取每个块并将其写入文件系统中的适当位置。
最后
\lstinline{end_op}
写入计数为零的日志头；
这必须在下一个事务开始写入日志块之前发生，以便崩溃不会导致恢复使用一个事务的头与后续事务的日志块。

\indexcode{recover_from_log}
\lineref{kernel/log.c:/^recover_from_log/}
由
\indexcode{initlog}
\lineref{kernel/log.c:/^initlog/}
调用，
后者在第一个用户进程运行之前从 \indexcode{fsinit}\lineref{kernel/fs.c:/^fsinit/} 在引导期间调用
\lineref{kernel/proc.c:/fsinit/}。
它读取日志头，如果头指示日志包含已提交的事务，则模拟
\lstinline{end_op}
的操作。

日志的一个示例使用发生在
\indexcode{filewrite}
\lineref{kernel/file.c:/^filewrite/}。
事务如下所示：
\begin{lstlisting}[]
      begin_op();
      ilock(f->ip);
      r = writei(f->ip, ...);
      iunlock(f->ip);
      end_op();
\end{lstlisting}
这段代码包装在一个循环中，将大写入拆分为每次只有几个扇区的单独事务，以避免溢出日志。
对
\indexcode{writei}
的调用作为此事务的一部分写入许多块：文件的 inode、一个或多个位图块和一些数据块。
%%
%%
%%
\section{代码：块分配器}

文件和目录内容存储在磁盘块中，必须从空闲池中分配。
Xv6 的块分配器在磁盘上维护一个空闲位图，每个块一位。
零位表示相应的块是空闲的；
一位表示它正在使用。
程序
\lstinline{mkfs}
设置对应于引导扇区、超级块、日志块、inode 块和位图块的位。

块分配器提供两个函数：
\indexcode{balloc}
分配一个新的磁盘块，
\indexcode{bfree}
释放一个块。
\lstinline{balloc}
中的循环在
\lineref{kernel/fs.c:/^..for.b.=.0/}
考虑从块 0 到
\lstinline{sb.size}
（文件系统中的块数）的每个块。
它寻找位图为零的块，表示它是空闲的。
如果
\lstinline{balloc}
找到这样的块，它更新位图并返回该块。
为了提高效率，循环分为两部分。
外循环读取每个位图块。
内循环检查单个位图块中的所有
Bits-Per-Block (\lstinline{BPB})
位。
如果两个进程试图同时分配一个块可能发生的竞争被缓冲区缓存只允许一个进程一次使用任何一个位图块的事实阻止。

\lstinline{bfree}
\lineref{kernel/fs.c:/^bfree/}
找到正确的位图块并清除正确的位。
同样，
\lstinline{bread}
和
\lstinline{brelse}
暗示的独占使用避免了对显式锁定的需要。

与本章其余部分描述的大部分代码一样，
\lstinline{balloc}
和
\lstinline{bfree}
必须在事务内调用。
%%
%%  -------------------------------------------
%%
\section{索引节点层}

术语
\indextext{索引节点（inode）}
可以有两种相关含义之一。
它可能指包含文件大小和数据块号列表的磁盘上的数据结构。
或者``inode''可能指内存中的 inode，它包含磁盘上 inode 的副本以及内核内需要的额外信息。

磁盘上的 inode 被打包到磁盘上称为 inode 块的连续区域中。
每个 inode 大小相同，因此给定一个数字 n，很容易在磁盘上找到第 n 个 inode。
事实上，这个数字 n，称为 inode 号或 i-number，是 inode 在实现中的标识方式。

磁盘上的 inode 由
\indexcode{struct dinode}
\lineref{kernel/fs.h:/^struct.dinode/}
定义。
\lstinline{type}
字段区分文件、目录和特殊文件（设备）。
类型为零表示磁盘上的 inode 是空闲的。
\lstinline{nlink}
字段计算引用此 inode 的目录条目数，以识别何时应释放磁盘上的 inode 及其数据块。
\lstinline{size}
字段记录文件中内容的字节数。
\lstinline{addrs}
数组记录保存文件内容的磁盘块的块号。

内核在称为 \indexcode{itable} 的表中将活动 inode 集保存在内存中；
\indexcode{struct inode}
\lineref{kernel/file.h:/^struct.inode/}
是磁盘上
\lstinline{struct}
\lstinline{dinode}
的内存副本。
内核仅在内存中有 C 指针引用该 inode 时才将 inode 存储在内存中。
\lstinline{ref}
字段计算引用内存中 inode 的 C 指针数量，如果引用计数降为零，内核将从内存中丢弃该 inode。
\indexcode{iget}
和
\indexcode{iput}
函数获取和释放指向 inode 的指针，修改引用计数。
指向 inode 的指针可能来自文件描述符、当前工作目录和瞬态内核代码，如
\lstinline{kexec}。

xv6 的 inode 代码中有四种锁或类似锁的机制。
\lstinline{itable.lock}
保护 inode 在 inode 表中最多出现一次的不变量，以及内存中 inode 的
\lstinline{ref}
字段计算内存中指向该 inode 的指针数量的不变量。
每个内存中的 inode 有一个包含睡眠锁的
\lstinline{lock}
字段，它确保对 inode 字段（如文件长度）以及 inode 的文件或目录内容块的独占访问。
如果 inode 的
\lstinline{ref}
大于零，它会导致系统维护表中的 inode，而不是为不同的 inode 重用表条目。
最后，每个 inode 包含一个
\lstinline{nlink}
字段（在磁盘上，如果在内存中则复制在内存中），计算引用文件的目录条目数；
如果其链接计数大于零，xv6 不会释放 inode。

\lstinline{iget()}
返回的
\lstinline{struct}
\lstinline{inode}
指针保证在相应调用
\lstinline{iput()}
之前有效；
inode 不会被删除，指针引用的内存不会为不同的 inode 重用。
\lstinline{iget()}
提供对 inode 的非独占访问，因此可以有多个指针指向同一个 inode。
文件系统代码的许多部分依赖于
\lstinline{iget()}
的这种行为，
既为了持有对 inode 的长期引用（作为打开的文件和当前目录），也为了在操作多个 inode 的代码中防止竞争同时避免死锁（如路径名查找）。

\lstinline{iget}
返回的
\lstinline{struct}
\lstinline{inode}
可能没有任何有用的内容。
为了确保它持有磁盘上 inode 的副本，代码必须调用
\indexcode{ilock}。
这会锁定 inode（以便没有其他进程可以
\lstinline{ilock}
它）并从磁盘读取 inode，如果尚未读取的话。
\lstinline{iunlock}
释放 inode 上的锁。
将 inode 指针的获取与锁定分开有助于在某些情况下避免死锁，例如在目录查找期间。
多个进程可以持有
\lstinline{iget}
返回的 inode 的 C 指针，
但一次只有一个进程可以锁定 inode。

inode 表仅存储内核代码或数据结构持有 C 指针的 inode。
它的主要工作是同步多个进程的访问。
inode 表也恰好缓存经常使用的 inode，但缓存是次要的；
如果 inode 被频繁使用，缓冲区缓存可能会将其保留在内存中。
修改内存中 inode 的代码使用
\lstinline{iupdate}
将其写入磁盘。
%%
%%  -------------------------------------------
%%
\section{代码：索引节点}

要分配一个新的 inode（例如，创建文件时），
xv6 调用
\indexcode{ialloc}
\lineref{kernel/fs.c:/^ialloc/}。
\lstinline{ialloc}
类似于
\indexcode{balloc}：
它遍历磁盘上的 inode 结构，一次一个块，寻找一个标记为空闲的 inode。
当它找到一个时，它通过将新的
\lstinline{type}
写入磁盘来声明它，然后通过尾调用返回 inode 表中的一个条目
\indexcode{iget}
\lineref{kernel/fs.c:/return.iget\(dev..inum\)/}。
\lstinline{ialloc}
的正确操作依赖于这样一个事实：一次只有一个进程可以持有对
\lstinline{bp}
的引用：
\lstinline{ialloc}
可以确信没有其他进程同时看到该 inode 可用并试图声明它。

\lstinline{iget}
\lineref{kernel/fs.c:/^iget/}
查看 inode 表以查找具有所需设备和 inode 号的活动条目（\lstinline{ip->ref}
\lstinline{>}
\lstinline{0}）。
如果找到一个，它返回对该 inode 的新引用
\linerefs{kernel/fs.c:/^....if.ip->ref.>.0/,/^....}/}。
当
\indexcode{iget}
扫描时，它记录第一个空槽的位置
\linerefs{kernel/fs.c:/^....if.empty.==.0/,/empty.=.ip/}，
如果需要分配表条目，则使用它。

代码必须在读取或写入其元数据或内容之前使用
\indexcode{ilock}
锁定 inode。
\lstinline{ilock}
\lineref{kernel/fs.c:/^ilock/}
为此目的使用睡眠锁。
一旦
\indexcode{ilock}
获得对 inode 的独占访问，如果需要，它就从磁盘（更可能是缓冲区缓存）读取 inode。
函数
\indexcode{iunlock}
\lineref{kernel/fs.c:/^iunlock/}
释放睡眠锁，
这可能导致任何睡眠的进程被唤醒。

\lstinline{iput}
\lineref{kernel/fs.c:/^iput/}
通过递减引用计数来释放指向 inode 的 C 指针
\lineref{kernel/fs.c:/^..ip->ref--/}。
如果这是最后一个引用，inode 在 inode 表中的槽位现在空闲，可以重用于不同的 inode。

如果
\indexcode{iput}
看到没有 C 指针引用某个 inode 且该 inode 没有链接到它（不出现在任何目录中），则必须释放该 inode 及其数据块。
\lstinline{iput}
调用
\indexcode{itrunc}
将文件截断为零字节，释放数据块；
将 inode 类型设置为 0（未分配）；
并将 inode 写入磁盘
\lineref{kernel/fs.c:/inode.has.no.links.and/}。

\indexcode{iput}
在释放 inode 的情况下的锁定协议值得仔细研究。
一个危险是并发线程可能在
\lstinline{ilock}
中等待使用此 inode（例如，读取文件或列出目录），
并且不会准备发现该 inode 不再被分配。
这不可能发生，因为如果系统调用无法获得指向内存中 inode 的指针，如果它没有链接到它且
\lstinline{ip->ref}
为一，则无法获得指向内存中 inode 的指针。
该引用是调用
\lstinline{iput}
的线程拥有的引用。
另一个主要危险是对
\lstinline{ialloc}
的并发调用可能选择
\lstinline{iput}
正在释放的相同 inode。
这可能发生在
\lstinline{iupdate}
写入磁盘使 inode 具有零类型之后。
这种竞争是良性的；分配线程将礼貌地等待获取 inode 的睡眠锁，然后再读取或写入 inode，此时
\lstinline{iput}
已完成。

\lstinline{iput()}
可以写入磁盘。这意味着使用文件系统的任何系统调用都可能写入磁盘，因为系统调用可能是最后一个引用文件的系统调用。
即使是像
\lstinline{read()}
这样看似只读的调用，也可能最终调用
\lstinline{iput()}。
反过来，这意味着即使是只读的系统调用，如果使用文件系统，也必须包装在事务中。

\lstinline{iput()}
和崩溃之间存在一个具有挑战性的交互。
\lstinline{iput()}
不会立即截断文件，当文件的链接计数降为零时，因为某些进程可能仍持有对内存中 inode 的引用：进程可能仍在读写文件，因为它成功打开了它。
但是，如果在最后一个进程关闭文件的文件描述符之前发生崩溃，则文件将在磁盘上标记为已分配，但没有目录条目指向它。

文件系统以两种方式之一处理这种情况。简单的解决方案是在恢复期间，重启后，文件系统扫描整个文件系统以查找标记为已分配但没有目录条目指向它们的文件。
如果存在任何此类文件，则可以释放它们。

第二种解决方案不需要扫描文件系统。在这种解决方案中，文件系统记录链接计数降为零但引用计数不为零的文件的 inode inumber（例如，在超级块中）。
如果文件系统在其引用计数达到 0 时删除该文件，则通过从列表中删除该 inode 来更新磁盘上的列表。在恢复时，文件系统释放列表中的任何文件。

Xv6 没有实现任一解决方案，这意味着 inode 可能在磁盘上标记为已分配，即使它们不再使用。
这意味着随着时间的推移，xv6 有磁盘空间耗尽的风险。
%%
%%
%%
\section{代码：索引节点内容}

\begin{figure}[t]
\centering
\includegraphics[scale=0.5]{fig/inode.pdf}
\caption{磁盘上文件的表示。}
\label{fig:inode}
\end{figure}

磁盘上的 inode 结构，
\indexcode{struct dinode}，
包含一个大小和一个块号数组（参见图~\ref{fig:inode}）。
inode 数据位于
\lstinline{dinode}
的
\lstinline{addrs}
数组中列出的块中。
前
\indexcode{NDIRECT}
个数据块列在数组的前
\lstinline{NDIRECT}
个条目中；这些块称为
\indextext{直接块（direct blocks）}。
接下来的
\indexcode{NINDIRECT}
个数据块不在 inode 中列出，而是在一个称为
\indextext{间接块（indirect block）}
的数据块中列出。
\lstinline{addrs}
数组中的最后一个条目给出间接块的地址。
因此，文件的前 12 kB（
\lstinline{NDIRECT}
\lstinline{x}
\indexcode{BSIZE}）
字节可以从 inode 中列出的块加载，
而接下来的
\lstinline{256} kB（
\lstinline{NINDIRECT}
\lstinline{x}
\lstinline{BSIZE}）
字节只能在查阅间接块后加载。
这是一种很好的磁盘上表示，但对客户端来说很复杂。
函数
\indexcode{bmap}
管理这种表示，以便更高级别的例程，如
\indexcode{readi}
和
\indexcode{writei}，
我们稍后会看到，不需要管理这种复杂性。
\lstinline{bmap}
返回 inode
\lstinline{ip}
的第
\lstinline{bn}
个数据块的磁盘块号。
如果
\lstinline{ip}
还没有这样的块，
\lstinline{bmap}
分配一个。

函数
\indexcode{bmap}
\lineref{kernel/fs.c:/^bmap/}
从简单的情况开始：前
\indexcode{NDIRECT}
个块在 inode 本身中列出
\linerefs{kernel/fs.c:/^..if.bn.<.NDIRECT/,/^..}/}。
接下来的
\indexcode{NINDIRECT}
个块列在
\lstinline{ip->addrs[NDIRECT]}
的间接块中。
\lstinline{bmap}
读取间接块
\lineref{kernel/fs.c:/bp.=.bread.ip->dev..addr/}
然后从块中的正确位置读取块号
\lineref{kernel/fs.c:/a.=..uint\*.bp->data/}。
如果块号超过
\lstinline{NDIRECT+NINDIRECT}，
\lstinline{bmap}
发生 panic；
\lstinline{writei}
包含防止这种情况的检查
\lineref{kernel/fs.c:/off...n...MAXFILE.BSIZE/}。

\lstinline{bmap}
根据需要分配块。
\lstinline{ip->addrs[]}
或间接条目为零表示未分配块。
当
\lstinline{bmap}
遇到零时，它用新分配的块号替换它们
\linerefs{kernel/fs.c:/^....if..addr.=.*==.0/,/./}
\linerefs{kernel/fs.c:/^....if..addr.*NDIRECT.*==.0/,/./}。

\indexcode{itrunc}
释放文件的块，将 inode 的大小重置为零。
\lstinline{itrunc}
\lineref{kernel/fs.c:/^itrunc/}
从释放直接块开始
\linerefs{kernel/fs.c:/^..for.i.=.0.*NDIRECT/,/^..}/}，
然后是间接块中列出的那些
\linerefs{kernel/fs.c:/^....for.j.=.0.*NINDIRECT/,/^....}/}，
最后是间接块本身
\linerefs{kernel/fs.c:/^....bfree.*NDIRECT/,/./}。

\lstinline{bmap}
使
\indexcode{readi}
和
\indexcode{writei}
很容易访问 inode 的数据。
\lstinline{readi}
\lineref{kernel/fs.c:/^readi/}
首先确保偏移量和计数不超过文件末尾。
从文件末尾之后开始的读取返回错误
\linerefs{kernel/fs.c:/^..if.off.>.ip->size/,/./}
而从文件末尾开始或跨越文件末尾的读取返回比请求更少的字节
\linerefs{kernel/fs.c:/^..if.off.\+.n.>.ip->size/,/./}。
主循环处理文件的每个块，将数据从缓冲区复制到
\lstinline{dst}
\linerefs{kernel/fs.c:/^..for.tot=0/,/^..}/}。
%%  注意：很难为 writei 写行引用，因为很多行与 readi 中的相同。幸运的是，相同的行可能不需要注释。
\indexcode{writei}
\lineref{kernel/fs.c:/^writei/}
与
\indexcode{readi}
相同，但有三个例外：
从文件末尾开始或跨越文件末尾的写入会增长文件，直到最大文件大小
\linerefs{kernel/fs.c:/^..if.off.\+.n.>.MAXFILE/,/./}；
循环将数据复制到缓冲区中而不是从缓冲区复制出来
\lineref{kernel/fs.c:/either.copyin.*bp->data/}；
而且如果写入扩展了文件，
\indexcode{writei}
必须更新其大小
\linerefs{kernel/fs.c:/^..if.off.>.ip->size\)/,/./}。

函数
\indexcode{stati}
\lineref{kernel/fs.c:/^stati\(/}
将 inode 元数据复制到
\lstinline{stat}
结构体中，该结构体通过
\indexcode{stat}
系统调用暴露给用户程序。
%%
%%
%%
\section{代码：目录层}

目录在内部实现上很像文件。
它的 inode 具有类型
\indexcode{T_DIR}
其数据是一系列目录条目。
每个条目是一个
\indexcode{struct dirent}
\lineref{kernel/fs.h:/^struct.dirent/}，
包含一个名称和一个 inode 号。
名称最多为
\indexcode{DIRSIZ}
（14）个字符；
如果更短，则以 NULL（0）字节终止。
inode 号为零的目录条目是空闲的。

函数
\indexcode{dirlookup}
\lineref{kernel/fs.c:/^dirlookup/}
在目录中搜索具有给定名称的条目。
如果找到一个，它返回指向相应 inode 的指针（未锁定），
并设置
\lstinline{*poff}
为目录中条目的字节偏移量，
以防调用者希望编辑它。
如果
\lstinline{dirlookup}
找到具有正确名称的条目，
它更新
\lstinline{*poff}
并通过
\indexcode{iget}
返回一个未锁定的 inode。
\lstinline{dirlookup}
是
\lstinline{iget}
返回未锁定 inode 的原因。
调用者已锁定
\lstinline{dp}，
因此如果查找是针对
\indexcode{.}，
当前目录的别名，
尝试在返回之前锁定 inode 将试图重新锁定
\lstinline{dp}
并导致死锁。
（涉及多个进程和
\indexcode{..}，
父目录的别名；有更复杂的死锁场景；
\lstinline{.}
不是唯一的问题。）
调用者可以解锁
\lstinline{dp}
然后锁定
\lstinline{ip}，
确保它一次只持有一个锁。

函数
\indexcode{dirlink}
\lineref{kernel/fs.c:/^dirlink/}
将给定的名称和 inode 号的新目录条目写入目录
\lstinline{dp}。
如果名称已存在，
\lstinline{dirlink}
返回错误
\linerefs{kernel/fs.c:/Check.that.name.is.not.present/,/^..}/}。
主循环读取目录条目以查找未分配的条目。
当找到一个时，它提前停止循环
\linerefs{kernel/fs.c:/Look for an empty dirent/,/.*break/}，
将
\lstinline{off}
设置为可用条目的偏移量。
否则，循环以
\lstinline{off}
设置为
\lstinline{dp->size}
结束。
无论哪种方式，
\lstinline{dirlink}
然后通过写入偏移量
\lstinline{off}
将新条目添加到目录
\linerefs{kernel/fs.c:/if.writei/,/return -1/}。
%%
%%
%%
\section{代码：路径名}

路径名查找涉及对
\indexcode{dirlookup}
的连续调用，每个路径组件一个。
\lstinline{namei}
\lineref{kernel/fs.c:/^namei/}
评估
\lstinline{path}
并返回相应的
\lstinline{inode}。
函数
\indexcode{nameiparent}
是一个变体：它在最后一个元素之前停止，返回父目录的 inode 并将最终元素复制到
\lstinline{name}。
两者都调用通用函数
\indexcode{namex}
来完成实际工作。

\lstinline{namex}
\lineref{kernel/fs.c:/^namex/}
首先决定路径评估从哪里开始。
如果路径以斜杠开头，评估从根目录开始；
否则，从当前目录开始
\linerefs{kernel/fs.c:/..if.\*path.==....\)/,/idup/}。
然后它使用
\indexcode{skipelem}
依次考虑路径的每个元素
\lineref{kernel/fs.c:/while.*skipelem/}。
每次循环迭代必须在当前 inode
\lstinline{ip}
中查找
\lstinline{name}。
迭代从锁定
\lstinline{ip}
并检查它是否是目录开始。
如果不是，查找失败
\linerefs{kernel/fs.c:/^....ilock.ip/,/^....}/}。
（锁定
\lstinline{ip}
是必要的，不是因为
\lstinline{ip->type}
可能在脚下改变——它不能——而是因为直到
\indexcode{ilock}
运行，
\lstinline{ip->type}
不能保证已从磁盘加载。）
如果调用是
\indexcode{nameiparent}
且这是最后一个路径元素，循环提前停止，
根据
\lstinline{nameiparent}
的定义；
最终路径元素已复制到
\lstinline{name}，
因此
\indexcode{namex}
只需返回未锁定的
\lstinline{ip}
\linerefs{kernel/fs.c:/^....if.nameiparent/,/^....}/}。
最后，循环使用
\indexcode{dirlookup}
查找路径元素并通过设置
\lstinline{ip = next}
准备下一次迭代
\linerefs{kernel/fs.c:/^....if..next.*dirlookup/,/^....ip.=.next/}。
当循环用完路径元素时，它返回
\lstinline{ip}。

过程
\lstinline{namex}
可能需要很长时间才能完成：它可能涉及几个磁盘操作来读取路径中遍历的目录的 inode 和目录块（如果它们不在缓冲区缓存中）。Xv6 经过精心设计，以便如果一个内核对
\lstinline{namex}
的调用被磁盘 I/O 阻塞，另一个内核查找不同路径名可以并发进行。
\lstinline{namex}
单独锁定路径中的每个目录，以便不同目录中的查找可以并行进行。

这种并发引入了一些挑战。例如，当一个内核线程正在查找路径名时，另一个内核线程可能正在通过取消链接目录来更改目录树。一个潜在的风险是查找可能正在搜索已被另一个内核线程删除的目录，其块已被重用于另一个目录或文件。

Xv6 避免了这种竞争。例如，当在
\lstinline{namex}
中执行
\lstinline{dirlookup}
时，查找线程持有目录的锁，
\lstinline{dirlookup}
返回使用
\lstinline{iget}
获得的 inode。
\lstinline{iget}
增加 inode 的引用计数。只有在从
\lstinline{dirlookup}
接收到 inode 后，
\lstinline{namex}
才释放目录上的锁。现在另一个线程可以从目录中取消链接 inode，但 xv6 还不会删除 inode，因为 inode 的引用计数仍然大于零。

另一个风险是死锁。例如，
\lstinline{next}
在查找 "。" 时指向与
\lstinline{ip}
相同的 inode。
在释放
\lstinline{ip}
上的锁之前锁定
\lstinline{next}
将导致死锁。
为了避免这种死锁，
\lstinline{namex}
在获取
\lstinline{next}
上的锁之前解锁目录。
这里我们再次看到为什么
\lstinline{iget}
和
\lstinline{ilock}
之间的分离很重要。
%%
%%
%%
\section{文件描述符层}

Unix 接口的一个很酷的方面是 Unix 中的大多数资源都表示为文件，包括控制台、管道等设备，当然还有真正的文件。文件描述符层是实现这种统一性的层。

Xv6 给每个进程自己的打开文件表，或文件描述符，如我们在第~\ref{CH:UNIX} 章看到的。
每个打开的文件由
\indexcode{struct file}
\lineref{kernel/file.h:/^struct.file/}
表示，它是一个围绕 inode 或管道的包装器，加上一个 I/O 偏移量。
每次对
\indexcode{open}
的调用创建一个新的打开文件（一个新的
\lstinline{struct}
\lstinline{file}）：
如果多个进程独立打开同一个文件，不同的实例将具有不同的 I/O 偏移量。
另一方面，单个打开文件（相同的
\lstinline{struct}
\lstinline{file}）
可以在一个进程的文件表中出现多次，也可以在多个进程的文件表中出现。
如果一个进程使用
\lstinline{open}
打开文件，然后使用
\indexcode{dup}
创建别名或使用
\indexcode{fork}
与子进程共享，就会发生这种情况。
引用计数跟踪对特定打开文件的引用数量。
文件可以为读或写或两者打开。
\lstinline{readable}
和
\lstinline{writable}
字段跟踪这一点。

系统中的所有打开文件都保存在全局文件表
\indexcode{ftable}
中。
文件表具有分配文件
（\indexcode{filealloc}）、
创建重复引用
（\indexcode{filedup}）、
释放引用
（\indexcode{fileclose}）、
以及读写数据
（\indexcode{fileread}
和
\indexcode{filewrite}）的函数。

前三个遵循现在熟悉的形式。
\lstinline{filealloc}
\lineref{kernel/file.c:/^filealloc/}
扫描文件表以查找未引用的文件
（\lstinline{f->ref}
\lstinline{==}
\lstinline{0}）
并返回新引用；
\indexcode{filedup}
\lineref{kernel/file.c:/^filedup/}
递增引用计数；
\indexcode{fileclose}
\lineref{kernel/file.c:/^fileclose/}
递减它。
当文件的引用计数达到零时，
\lstinline{fileclose}
根据类型释放底层管道或 inode。

函数
\indexcode{filestat}、
\indexcode{fileread}、
和
\indexcode{filewrite}
实现
\indexcode{stat}、
\indexcode{read}、
和
\indexcode{write}
文件操作。
\lstinline{filestat}
\lineref{kernel/file.c:/^filestat/}
仅允许在 inode 上并调用
\indexcode{stati}。
\lstinline{fileread}
和
\lstinline{filewrite}
检查打开模式是否允许操作，然后
将调用传递给管道或 inode 实现。
如果文件表示 inode，
\lstinline{fileread}
和
\lstinline{filewrite}
使用 I/O 偏移量作为操作的偏移量，然后推进它
\linerefs{kernel/file.c:/readi/,/./}
\linerefs{kernel/file.c:/writei/,/./}。
管道没有偏移量的概念。
回想一下 inode 函数要求调用者处理锁定
\linerefs{kernel/file.c:/stati/-1,/iunlock/}
\linerefs{kernel/file.c:/readi/-1,/iunlock/}
\linerefs{kernel/file.c:/writei\(f/-1,/iunlock/}。
inode 锁定有一个方便的副作用，即读写偏移量原子地更新，以便多个同时写入同一文件不能覆盖彼此的数据，尽管它们的写入可能最终交错。
%%
%%
%%
\section{代码：系统调用}

有了低层提供的函数，大多数系统调用的实现是微不足道的
（参见
\fileref{kernel/sysfile.c}）。
有一些调用值得仔细研究。

函数
\indexcode{sys_link}
和
\indexcode{sys_unlink}
编辑目录，创建或删除对 inode 的引用。
它们是使用事务的强大功能的另一个好例子。
\lstinline{sys_link}
\lineref{kernel/sysfile.c:/^sys_link/}
首先获取其参数，两个字符串
\lstinline{old}
和
\lstinline{new}
\lineref{kernel/sysfile.c:/argstr.*old.*new/}。
假设
\lstinline{old}
存在且不是目录
\linerefs{kernel/sysfile.c:/namei.old/,/^..}/}，
\lstinline{sys_link}
递增其
\lstinline{ip->nlink}
计数。
然后
\lstinline{sys_link}
调用
\indexcode{nameiparent}
以查找
\lstinline{new}
的父目录和最终路径元素
\lineref{kernel/sysfile.c:/nameiparent.new/}
并创建一个新的目录条目指向
\lstinline{old}
的 inode
\lineref{kernel/sysfile.c:/\|\| dirlink/}。
新的父目录必须存在且与现有 inode 在同一设备上：
inode 号只在单个磁盘上具有唯一含义。
如果发生这样的错误，
\indexcode{sys_link}
必须返回并递减
\lstinline{ip->nlink}。

事务简化了实现，因为它需要更新多个磁盘块，但我们不必担心执行它们的顺序。它们要么全部成功，要么都不成功。
例如，没有事务，在创建链接之前更新
\lstinline{ip->nlink}
将使文件系统暂时处于不安全状态，中间的崩溃可能导致严重破坏。
有了事务，我们不必担心这个。

\lstinline{sys_link}
为现有 inode 创建新名称。
函数
\indexcode{create}
\lineref{kernel/sysfile.c:/^create/}
为新 inode 创建新名称。
它是三个文件创建系统调用的概括：
带
\indexcode{O_CREATE}
标志的
\indexcode{open}
创建新普通文件，
\indexcode{mkdir}
创建新目录，
\indexcode{mkdev}
创建设备文件。
像
\indexcode{sys_link}
一样，
\indexcode{create}
首先调用
\indexcode{nameiparent}
以获取父目录的 inode。
然后它调用
\indexcode{dirlookup}
检查名称是否已存在
\lineref{kernel/sysfile.c:/dirlookup.*\[^=\]=.0/}。
如果名称确实存在，
\lstinline{create}
的行为取决于它为哪个系统调用使用：
\lstinline{open}
与
\indexcode{mkdir}
和
\indexcode{mkdev}
具有不同的语义。
如果
\lstinline{create}
代表
\lstinline{open}
使用
（\lstinline{type}
\lstinline{==}
\indexcode{T_FILE}）
且存在的名称本身是普通文件，
则
\lstinline{open}
将其视为成功，
因此
\lstinline{create}
也视为成功
\lineref{kernel/sysfile.c:/^......return.ip/}。
否则，它是错误
\linerefs{kernel/sysfile.c:/^......return.ip/+1,/return.0/}。
如果名称尚不存在，
\lstinline{create}
现在使用
\indexcode{ialloc}
分配新 inode
\lineref{kernel/sysfile.c:/ialloc/}。
如果新 inode 是目录，
\lstinline{create}
用
\indexcode{.}
和
\indexcode{..}
条目初始化它。
最后，既然数据已正确初始化，
\indexcode{create}
可以将其链接到父目录
\lineref{kernel/sysfile.c:/if.dirlink/}。
\lstinline{create}，
像
\indexcode{sys_link}
一样，同时持有两个 inode 锁：
\lstinline{ip}
和
\lstinline{dp}。
没有死锁的可能性，因为 inode
\lstinline{ip}
是新分配的：系统中的其他进程不会持有
\lstinline{ip}
的锁然后试图锁定
\lstinline{dp}。

使用
\lstinline{create}，
很容易实现
\indexcode{sys_open}、
\indexcode{sys_mkdir}、
和
\indexcode{sys_mknod}。
\lstinline{sys_open}
\lineref{kernel/sysfile.c:/^sys_open/}
是最复杂的，因为创建新文件只是它能做的一小部分。
如果
\indexcode{open}
被传递
\indexcode{O_CREATE}
标志，它调用
\lstinline{create}
\lineref{kernel/sysfile.c:/create.*T_FILE/}。
否则，它调用
\indexcode{namei}
\lineref{kernel/sysfile.c:/if..ip.=.namei.path/}。
\lstinline{create}
返回锁定的 inode，但
\lstinline{namei}
不返回，因此
\indexcode{sys_open}
必须自己锁定 inode。
这提供了一个方便的地方来检查目录是否只打开用于读取而不是写入。
假设 inode 以某种方式获得，
\lstinline{sys_open}
分配文件和文件描述符
\lineref{kernel/sysfile.c:/filealloc.*fdalloc/}
然后填充文件
\linerefs{kernel/sysfile.c:/type.=.FD_INODE/,/writable/}。
请注意，没有其他进程可以访问部分初始化的文件，因为它只在当前进程的表中。

第~\ref{CH:SLEEP} 章在我们拥有文件系统之前检查了管道的实现。
函数
\indexcode{sys_pipe}
通过提供创建管道对的方式将该实现连接到文件系统。
它的参数是指向两个整数空间的指针，
它将在其中记录两个新的文件描述符。
然后它分配管道并安装文件描述符。
%%
%%  -------------------------------------------
%%
\section{现实世界}

真实操作系统中的缓冲区缓存比 xv6 的复杂得多，但它服务于相同的两个目的：缓存和同步对磁盘的访问。
Xv6 的缓冲区缓存，像 V6 的一样，使用简单的最近最少使用（LRU）逐出策略；
可以实现许多更复杂的策略，每种策略对某些工作负载好，对其他工作负载不太好。
更高效的 LRU 缓存将消除链表，
而是使用哈希表进行查找，使用堆进行 LRU 逐出。
现代缓冲区缓存通常与虚拟内存系统集成以支持内存映射文件。

Xv6 的日志系统效率低下。
提交不能与文件系统系统调用并发发生。
系统记录整个块，即使块中只有几个字节被更改。它执行同步日志写入，一次一个块，每个都可能需要整个磁盘旋转时间。真正的日志系统解决了所有这些问题。

日志不是提供崩溃恢复的唯一方法。早期文件系统在重启期间使用清除器（例如，Unix
\indexcode{fsck}
程序）来检查每个文件和目录以及块和 inode 空闲列表，寻找并解决不一致性。清除大型文件系统可能需要数小时，而且有些情况下无法以使原始系统调用原子化的方式解决不一致性。从日志中恢复要快得多，并使系统调用在面对崩溃时是原子的。

Xv6 使用与早期 Unix 相同的 inode 和目录的基本磁盘布局；
这种方案多年来一直非常持久。
BSD 的 UFS/FFS 和 Linux 的 ext2/ext3 使用本质上相同的数据结构。
文件系统布局中效率最低的部分是目录，它需要在每次查找时对所有磁盘块进行线性扫描。
当目录只有几个磁盘块时这是合理的，但对于包含许多文件的目录来说很昂贵。
Microsoft Windows 的 NTFS、macOS 的 HFS 和 Solaris 的 ZFS，仅举几例，将目录实现为磁盘上块的平衡树。
这很复杂，但保证了对数时间的目录查找。

Xv6 对磁盘故障很天真：如果磁盘操作失败，xv6 发生 panic。
这是否合理取决于硬件：
如果操作系统位于使用冗余来屏蔽磁盘故障的特殊硬件之上，也许操作系统很少看到故障，因此 panic 是可以的。
另一方面，使用普通磁盘的操作系统应该预期故障并更优雅地处理它们，
以便一个文件中一个块的丢失不会影响文件系统其余部分的使用。

Xv6 要求文件系统适合一个磁盘设备且大小不变。
随着大型数据库和多媒体文件推动存储需求越来越高，操作系统正在开发消除``每个文件系统一个磁盘''瓶颈的方法。
基本方法是将许多磁盘组合成一个逻辑磁盘。硬件解决方案如 RAID 仍然是最受欢迎的，但当前趋势是尽可能在软件中实现这种逻辑。
这些软件实现通常允许丰富的功能，如通过动态添加或删除磁盘来增长或缩小逻辑设备。
当然，可以动态增长或缩小的存储层需要可以执行相同操作的文件系统：xv6 使用的 inode 块的固定大小数组在此类环境中效果不佳。
将磁盘管理与文件系统分离可能是最简洁的设计，但两者之间的复杂接口导致一些系统，如 Sun 的 ZFS，将它们组合在一起。

Xv6 的文件系统缺乏现代文件系统的许多其他功能；例如，它缺乏对快照和增量备份的支持。

现代 Unix 系统允许使用与磁盘存储相同的系统调用来访问许多类型的资源：
命名管道、网络连接、远程访问的网络文件系统，以及监视和控制接口，如
\lstinline{/proc}。
与 xv6 的
\lstinline{if}
语句在
\indexcode{fileread}
和
\indexcode{filewrite}
中不同，这些系统通常给每个打开的文件一个函数指针表，
每个操作一个，
并调用函数指针来调用该 inode 的调用实现。
网络文件系统和用户级文件系统提供将这些调用转换为网络 RPC 的函数，并在返回之前等待响应。
%%
%%  -------------------------------------------
%%
\section{练习}

\begin{enumerate}

\item 为什么
\lstinline{balloc}
会发生 panic？
xv6 能恢复吗？

\item 为什么
\lstinline{ialloc}
会发生 panic？
xv6 能恢复吗？

\item 为什么
\lstinline{filealloc}
在文件用完时不发生 panic？
为什么这更常见因此值得处理？

\item 假设与
\lstinline{ip}
对应的文件在
\lstinline{sys_link}
对
\lstinline{iunlock(ip)}
和
\lstinline{dirlink}
的调用之间被另一个进程取消链接。
链接会正确创建吗？
为什么或为什么不？

\item
\lstinline{create}
进行四个函数调用（一个到
\lstinline{ialloc}
和三个到
\lstinline{dirlink}）
它要求成功。
如果任何一个失败，
\lstinline{create}
调用
\lstinline{panic}。
为什么这是可接受的？
为什么这四个调用都不能失败？

\item
\lstinline{sys_chdir}
在
\lstinline{iput(cp->cwd)}
之前调用
\lstinline{iunlock(ip)}，
后者可能试图锁定
\lstinline{cp->cwd}，
但将
\lstinline{iunlock(ip)}
推迟到
\lstinline{iput}
之后不会导致死锁。
为什么不呢？

\item 实现
\lstinline{lseek}
系统调用。支持
\lstinline{lseek}
还需要修改
\lstinline{filewrite}
以用零填充文件中的空洞，如果
\lstinline{lseek}
设置
\lstinline{off}
超过
\lstinline{f->ip->size.}

\item 将
\lstinline{O_TRUNC}
和
\lstinline{O_APPEND}
添加到
\lstinline{open}，
以便
\lstinline{>}
和
\lstinline{>>}
操作符在 shell 中工作。

\item 修改文件系统以支持符号链接。

\item 修改文件系统以支持命名管道。

\item 修改文件和 VM 系统以支持内存映射文件。

\end{enumerate}

\chapter{并发重访}
\label{CH:LOCK2}

在内核设计中，同时实现良好的并行性能、确保并发正确性并保持代码可理解性，是一个巨大的挑战。直接使用锁是实现正确性的最佳途径，但并非总是可行。本章重点介绍 xv6 被迫以复杂方式使用锁的例子，以及 xv6 使用类锁技术但不用锁的例子。

\section{锁模式}

缓存项通常对加锁是一个挑战。例如，文件系统的块缓存 \lineref{kernel/bio.c:/^struct/} 存储最多 {\tt NBUF} 个磁盘块的副本。给定磁盘块在缓存中最多只有一个副本是至关重要的；否则，不同的进程可能会对应该是同一块的不同副本做出冲突的更改。每个缓存块存储在 {\tt struct buf} \fileref{kernel/buf.h} 中。{\tt struct buf} 有一个锁字段，有助于确保一次只有一个进程使用给定的磁盘块。然而，这个锁是不够的：如果块根本不在缓存中，两个进程想要同时使用它怎么办？没有 {\tt struct buf}（因为块尚未缓存），因此没有什么可以锁定的。Xv6 通过将额外的锁（{\tt bcache.lock}）与缓存块的身份集相关联来处理这种情况。需要检查块是否已缓存的代码（例如，{\tt bget} \lineref{kernel/bio.c:/^bget/}）或更改缓存块集的代码必须持有 {\tt bcache.lock}；在该代码找到块和它需要的 {\tt struct buf} 之后，它可以释放 {\tt bcache.lock} 并只锁定特定的块。这是一种常见模式：一个锁用于项集，加上每个项一个锁。

通常，获取锁的同一个函数会释放它。但更精确的看法是，锁在必须表现为原子的序列开始时获取，在该序列结束时释放。如果序列在不同的函数、不同的线程或不同的 CPU 上开始和结束，那么锁的获取和释放也必须这样做。锁的作用是强制其他使用者等待，而不是将某段数据固定到特定的执行者。一个例子是 {\tt yield} 中的 {\tt acquire}\lineref{kernel/proc.c:/^yield/}，它是在调度器线程中释放的，而不是在获取它的进程中。另一个例子是 {\tt ilock} 中的 {\tt acquiresleep}\lineref{kernel/fs.c:/^ilock/}；这段代码在读取磁盘时经常睡眠；它可能在不同的 CPU 上唤醒，这意味着锁可能在不同的 CPU 上获取和释放。

释放受嵌入在对象中的锁保护的对象是一项微妙的业务，因为拥有锁不足以保证释放是正确的。问题情况发生在其他线程在 {\tt acquire} 中等待使用对象时；释放对象隐式释放嵌入的锁，这将导致等待线程故障。一种解决方案是跟踪对对象的引用数量，以便仅在最后一个引用消失时才释放它。参见 {\tt pipeclose} \lineref{kernel/pipe.c:/^pipeclose/} 作为示例；{\tt pi->readopen} 和 {\tt pi->writeopen} 跟踪管道是否有文件描述符引用它。

通常，你会在对相关项目集进行读写序列周围看到锁；锁确保其他线程只看到已完成的更新序列（只要它们也加锁）。如果更新是对单个共享变量的简单写入呢？例如，\texttt{setkilled} 和 \texttt{killed} \lineref{kernel/proc.c:/^setkilled/} 在它们对 \lstinline{p->killed} 的简单使用周围加锁。如果没有锁，一个线程可能在另一个线程读取它的同时写入 \lstinline{p->killed}。这是一个 \indextextx{竞态（race condition）}，C 语言规范说竞态产生 \indextext{未定义行为（undefined behavior）}，这意味着程序可能崩溃或产生不正确的结果。锁防止竞态并避免未定义行为。

竞态可能破坏程序的一个原因是，如果没有锁或等效构造，编译器可能生成以与原始 C 代码截然不同的方式读写内存的机器代码。例如，调用 \texttt{killed} 的线程的机器代码可能将 \lstinline{p->killed} 复制到寄存器，并且只读取那个缓存值；这意味着该线程可能永远看不到对 \lstinline{p->killed} 的任何写入。锁防止这种缓存。

% sleep locks.
% hand-over-hand locking in namei.
%   namei's use of refcount to prevent changing underfoot,
%     and lock when actually using it.
% example of where deadlock is avoided? namei?
% spawn a thread to evade a lock order problem?

\section{类似锁的模式}

在许多地方，xv6 以类似锁的方式使用引用计数或标志来指示对象已分配且不应被释放或重新使用。进程的 {\tt p->state} 以这种方式工作，{\tt file}、{\tt inode} 和 {\tt buf} 结构中的引用计数也是如此。虽然在每种情况下锁都保护标志或引用计数，但正是后者防止对象被过早释放。

文件系统使用 {\tt struct inode} 引用计数作为一种可以被多个进程持有的共享锁，以避免如果代码使用普通锁会发生的死锁。例如，{\tt namex} \lineref{kernel/fs.c:/^namex/} 中的循环依次锁定路径名命名的每个目录。然而，{\tt namex} 必须在循环结束时释放每个锁，因为如果它持有多个锁，如果路径名包含点（例如，{\tt a/./b}），它可能会与自己死锁。它也可能与涉及目录和 {\tt ..} 的并发查找死锁。如 Chapter~\ref{CH:FS} 所解释的，解决方案是让循环将其引用计数递增的目录 inode 带到下一次迭代，但不锁定。

某些数据项在不同时间受不同机制保护，有时可能由 xv6 代码的结构隐式保护免受并发访问，而不是由显式锁保护。例如，当物理页面空闲时，它受 \texttt{kmem.lock} \lineref{kernel/kalloc.c:/^. kmem;/} 保护。如果页面然后被分配为管道 \lineref{kernel/pipe.c:/^pipealloc/}，它受不同的锁保护（嵌入的 \lstinline{pi->lock}）。如果页面被重新分配给新进程的用户内存，它根本不受锁保护。相反，分配器不会将该页面给予任何其他进程（直到它被释放）这一事实保护它免受并发访问。
新进程内存的所有权很复杂：首先在 {\tt fork} 中父进程分配和操作它，然后子进程使用它，（子进程退出后）父进程再次拥有内存并将其传递给 {\tt kfree}。这里有两个教训：数据对象在其生命周期的不同时间点可能以不同方式免受并发保护，保护可能采取隐式结构的形式而不是显式锁。

最后一个类似锁的例子是在调用 {\tt mycpu()} \lineref{kernel/proc.c:/^myproc/} 周围需要禁用中断。禁用中断使调用代码相对于可能强制上下文切换的定时器中断是原子的，从而将进程移动到不同的 CPU。

\section{完全没有锁}

有几个地方 xv6 完全没有锁就共享可变数据。一个是在自旋锁的实现中，尽管可以将 RISC-V 原子指令视为依赖于硬件中实现的锁。另一个是 {\tt main.c} 中的 {\tt started} 变量 \lineref{kernel/main.c:/^volatile/}，用于防止其他 CPU 运行，直到 CPU 零完成初始化 xv6；{\tt volatile} 确保编译器实际生成加载和存储指令。

Xv6 包含一个 CPU 或线程写入一些数据，另一个 CPU 或线程读取数据，但没有特定锁专门保护该数据的情况。例如，在 {\tt fork} 中，父进程写入子进程的用户内存页，子进程（不同的线程，可能在不同的 CPU 上）读取这些页面；没有锁显式保护这些页面。这严格来说不是锁问题，因为子进程在父进程完成写入后才开始执行。这是一个潜在的内存排序问题（参见 Chapter~\ref{CH:LOCK}），因为如果没有内存屏障，没有理由期望一个 CPU 看到另一个 CPU 的写入。然而，由于父进程释放锁，子进程在启动时获取锁，{\tt acquire} 和 {\tt release} 中的内存屏障确保子进程的 CPU 看到父进程的写入。

\section{并行性}

锁主要用于在正确性的利益下抑制并行性。因为性能也很重要，内核设计者通常必须考虑如何以既能实现正确性又能允许并行性的方式使用锁。虽然 xv6 不是为高性能系统设计的，但考虑哪些 xv6 操作可以并行执行，哪些可能在锁上冲突仍然值得。

Xv6 中的管道是适度良好并行性的一个例子。每个管道有自己的锁，因此不同进程可以在不同 CPU 上并行读写不同管道。对于给定管道，然而，写入者和读取者必须等待对方释放锁；它们不能同时读/写同一管道。从空管道读取（或写入满管道）必须阻塞的情况也是如此，但这不是由于锁定方案。

上下文切换是一个更复杂的例子。两个内核线程，每个在自己的 CPU 上执行，可以同时调用 {\tt yield}、{\tt sched} 和 {\tt swtch}，调用将并行执行。线程各自持有一个锁，但它们是不同的锁，因此它们不必等待对方。然而，一旦进入 {\tt scheduler}，两个 CPU 可能在搜索 {\tt RUNNABLE} 进程表时冲突在锁上。也就是说，xv6 可能在上下文切换期间从多个 CPU 获得性能提升，但也许不如它可能获得的那么多。

另一个例子是来自不同 CPU 上不同进程的 {\tt fork} 的并发调用。调用可能不得不等待彼此的 {\tt pid\_lock} 和 {\tt kmem\_lock}，以及在进程表中搜索 {\tt UNUSED} 进程所需的每个进程锁。另一方面，两个 fork 进程可以完全并行地复制用户内存页和格式化页表页。

上述每个例子中的锁定方案在某些情况下牺牲了并行性能。在每种情况下，可以使用更复杂的设计获得更多并行性。是否值得取决于细节：相关操作被调用的频率，代码在持有争用锁时花费多长时间，多少 CPU 可能同时运行冲突操作，代码的其他部分是否是更具限制性的瓶颈。很难猜测给定的锁定方案是否可能导致性能问题，或者新设计是否显著更好，因此通常需要在实际工作负载上进行测量。

\section{练习}

\begin{enumerate}

\item 修改 xv6 的管道实现，允许在不同 CPU 上同时读取和写入同一管道。

\item 修改 xv6 的 \texttt{scheduler()} 以减少当不同 CPU 同时查找可运行进程时的锁争用。

\item 消除 xv6 的 \texttt{fork()} 中的一些串行化。

\end{enumerate}

\chapter{总结}
\label{CH:SUM}

本文通过逐行研究一个操作系统 xv6，介绍了操作系统的主要思想。一些代码行体现了主要思想的本质（例如，上下文切换、用户/内核边界、锁等），每一行都很重要；其他代码行提供了如何实现特定操作系统思想的示例，可以很容易地以不同方式完成（例如，更好的调度算法、更好的磁盘数据结构来表示文件、更好的日志以允许并发事务等）。所有这些思想都在一个特定的、非常成功的系统调用接口（Unix 接口）的背景下进行了说明，但这些思想也适用于其他操作系统的设计。



{
% The following prevents latex from splitting a bibliography entry with a page
% break
\interlinepenalty=10000
% Since we're using natbib in numbers mode, we don't need plainnat,
% which exists to feed authors and years back in to natbib.  As a
% result, it complains about entries without years, which we don't
% care about.
%\bibliographystyle{plainnat}
\bibliographystyle{plain}
\bibliography{book}
}

\printindex

\end{document}
